% =====================================================================
%  Chapitre 4 : Portes quantiques et systèmes multi-qubits
%  Fichier importé par main.tex avec % =====================================================================
%  Chapitre 4 : Portes quantiques et systèmes multi-qubits
%  Fichier importé par main.tex avec % =====================================================================
%  Chapitre 4 : Portes quantiques et systèmes multi-qubits
%  Fichier importé par main.tex avec % =====================================================================
%  Chapitre 4 : Portes quantiques et systèmes multi-qubits
%  Fichier importé par main.tex avec \include{chapitre4}.
%  Il ne contient ni préambule ni \begin{document} : les environnements
%  (definition, resultat, remarque, aretenir, savoirfaire, exercice, correction),
%  les macros (\ii, \ee, \dd, \pd, \pdd, \Hop, \ket, \bra, \braket), les symboles
%  de circuits (ctl, cible, croix, mesure, porte) et les couleurs (insarouge,
%  insagris) sont définis dans main.tex. Un chapitre = une \section.
% =====================================================================

% Macros de Dirac (déjà définies dans main.tex ; ce bloc évite l'erreur
% « Undefined control sequence \ket » si main.tex est une ancienne version).
\providecommand{\ket}[1]{|#1\rangle}
\providecommand{\bra}[1]{\langle #1|}
\providecommand{\braket}[2]{\langle #1|#2\rangle}

% =====================================================================
\section{Portes quantiques et systèmes multi-qubits}\label{chap:portes}
% =====================================================================

Un calcul quantique est une suite de transformations unitaires appliquées à des qubits. Ce
chapitre décrit les portes à un qubit (Pauli, Hadamard, phase) et leur effet sur la sphère de
Bloch, la lecture d'un circuit, puis les portes qui agissent sur plusieurs qubits : CNOT, SWAP et
portes contrôlées.

% ---------------------------------------------------------------------
\subsection{Portes quantiques}\label{sec:portes}
% ---------------------------------------------------------------------

Un calcul quantique transforme des états. Les opérations sur un seul qubit sont appelées
\emph{portes} ; les sections suivantes passent à plusieurs qubits.

\begin{resultat}{Postulat 4 : évolution d'un système isolé}
L'évolution d'un système isolé est décrite par un opérateur unitaire : $\ket{\Psi}\to U\ket{\Psi}$
avec $U^\dagger U=\mathrm{Id}$.
\end{resultat}

C'est la forme de l'équation de Schrödinger du chapitre~\ref{chap:histoire}. Pour un hamiltonien
$\Hop$ indépendant du temps, $\ii\hbar\,\partial_t\ket{\Psi}=\Hop\ket{\Psi}$ a pour solution
$\ket{\Psi(t)}=U(t)\ket{\Psi(0)}$ avec $U(t)=\ee^{-\ii\Hop t/\hbar}$. Comme $\Hop$ est hermitien,
$U(t)^\dagger=\ee^{+\ii\Hop t/\hbar}$ et $U^\dagger U=\mathrm{Id}$ : l'évolution est unitaire. Une
porte est une telle évolution, obtenue en contrôlant $\Hop$ pendant une durée donnée.

\begin{definition}{Porte quantique}
Une porte à un qubit est un opérateur $U$ (matrice $2\times2$) \emph{unitaire} :
$U^\dagger U=\mathrm{Id}$. Elle transforme $\ket{\Psi}\to U\ket{\Psi}$, conserve la norme
($\bra{\Psi}U^\dagger U\ket{\Psi}=\braket{\Psi}{\Psi}=1$) et est réversible
($U^{-1}=U^\dagger$). Sur la sphère de Bloch (section~\ref{sec:bloch}), c'est une
\emph{rotation}.
\end{definition}

% ---------------------------------------------------------------------
\subsection{Portes de Pauli}\label{sec:pauli}
% ---------------------------------------------------------------------

Les matrices de Pauli de la section~\ref{sec:op-dirac} sont trois portes.

\begin{definition}{Portes $X$, $Y$, $Z$}
\begin{center}
\renewcommand{\arraystretch}{1}
\begin{tabular}{@{}lccl@{}}
\textbf{Porte} & \textbf{Matrice} & \textbf{Ket-bra} & \textbf{Nom} \\[0.8ex]
$X=\sigma_x$ & $\begin{pmatrix}0 & 1\\ 1 & 0\end{pmatrix}$
  & $\ket{0}\bra{1}+\ket{1}\bra{0}$ & bit flip \\[2.2ex]
$Y=\sigma_y$ & $\begin{pmatrix}0 & -\ii\\ \ii & 0\end{pmatrix}$
  & $-\ii\ket{0}\bra{1}+\ii\ket{1}\bra{0}$ &  \\[2.2ex]
$Z=\sigma_z$ & $\begin{pmatrix}1 & 0\\ 0 & -1\end{pmatrix}$
  & $\ket{0}\bra{0}-\ket{1}\bra{1}$ & phase flip \\
\end{tabular}
\end{center}
\end{definition}

\paragraph{Porte $X$ (inverseur).} Avec $\braket{0}{0}=\braket{1}{1}=1$ et
$\braket{0}{1}=\braket{1}{0}=0$ :
\begin{align*}
  X\ket{0}&=\ket{0}\braket{1}{0}+\ket{1}\braket{0}{0}=0\cdot\ket{0}+1\cdot\ket{1}=\ket{1},\\
  X\ket{1}&=\ket{0}\braket{1}{1}+\ket{1}\braket{0}{1}=1\cdot\ket{0}+0\cdot\ket{1}=\ket{0}.
\end{align*}
$X$ échange $\ket{0}$ et $\ket{1}$ : c'est le NOT quantique (\emph{bit flip}). Sur la
sphère de Bloch, c'est une rotation de $\pi$ autour de l'axe $x$ : le pôle nord devient
le pôle sud. Les états de l'axe $x$ ne bougent pas : $X\ket{+}=\ket{+}$ et
$X\ket{-}=-\ket{-}$.

\paragraph{Porte $Z$ (changement de phase).}
\begin{align*}
  Z\ket{0}&=\ket{0}\braket{0}{0}-\ket{1}\braket{1}{0}=\ket{0},\\
  Z\ket{1}&=\ket{0}\braket{0}{1}-\ket{1}\braket{1}{1}=-\ket{1},
\end{align*}
donc $Z$ change le signe de la composante sur $\ket{1}$ (\emph{phase flip}). Sur les
états de la base $X$, par linéarité :
\begin{align*}
  Z\ket{+}&=\tfrac{1}{\sqrt2}\big(Z\ket{0}+Z\ket{1}\big)=\tfrac{1}{\sqrt2}\big(\ket{0}-\ket{1}\big)=\ket{-},\\
  Z\ket{-}&=\tfrac{1}{\sqrt2}\big(Z\ket{0}-Z\ket{1}\big)=\tfrac{1}{\sqrt2}\big(\ket{0}+\ket{1}\big)=\ket{+}.
\end{align*}
Sous forme matricielle, $Z\ket{+}=\begin{pmatrix}1 & 0\\ 0 & -1\end{pmatrix}
\frac{1}{\sqrt2}\begin{pmatrix}1\\ 1\end{pmatrix}=\frac{1}{\sqrt2}\begin{pmatrix}1\\ -1\end{pmatrix}=\ket{-}$.
$Z$ est une rotation de $\pi$ autour de l'axe $z$ : elle laisse $\ket{0}$ et $\ket{1}$ en
place et échange $\ket{+}$ et $\ket{-}$.

\paragraph{Porte $Y$.} Avec $Y=-\ii\ket{0}\bra{1}+\ii\ket{1}\bra{0}$ :
\begin{align*}
  Y\ket{0}&=-\ii\ket{0}\braket{1}{0}+\ii\ket{1}\braket{0}{0}=\ii\ket{1},\\
  Y\ket{1}&=-\ii\ket{0}\braket{1}{1}+\ii\ket{1}\braket{0}{1}=-\ii\ket{0}.
\end{align*}
Sur les états de la base $X$ :
\begin{align*}
  Y\ket{+}&=\tfrac{1}{\sqrt2}\big(\ii\ket{1}-\ii\ket{0}\big)=-\ii\,\tfrac{1}{\sqrt2}\big(\ket{0}-\ket{1}\big)=-\ii\ket{-},\\
  Y\ket{-}&=\tfrac{1}{\sqrt2}\big(\ii\ket{1}+\ii\ket{0}\big)=\ii\,\tfrac{1}{\sqrt2}\big(\ket{0}+\ket{1}\big)=\ii\ket{+}.
\end{align*}
Un facteur global ($-\ii$ ou $\ii$) ne change pas l'état physique : $Y$ échange bien
$\ket{+}$ et $\ket{-}$. Sur les états de l'axe $y$ :
\begin{align*}
  Y\ket{{+}\ii}&=\tfrac{1}{\sqrt2}\big(Y\ket{0}+\ii Y\ket{1}\big)
   =\tfrac{1}{\sqrt2}\big(\ii\ket{1}+\ket{0}\big)=\ket{{+}\ii},\\
  Y\ket{{-}\ii}&=\tfrac{1}{\sqrt2}\big(Y\ket{0}-\ii Y\ket{1}\big)
   =\tfrac{1}{\sqrt2}\big(\ii\ket{1}-\ket{0}\big)=-\ket{{-}\ii}.
\end{align*}
$Y$ est une rotation de $\pi$ autour de l'axe $y$.

\begin{remarque}[Chaque Pauli fixe son axe]
$Z\ket{0}=\ket{0}$, $Z\ket{1}=-\ket{1}$ ; $X\ket{\pm}=\pm\ket{\pm}$ ; $Y\ket{{\pm}\ii}=\pm\ket{{\pm}\ii}$.
Les bases $Z$, $X$, $Y$ de la section~\ref{sec:bloch} sont les bases propres de
$Z$, $X$, $Y$. Une porte de Pauli laisse donc son propre axe invariant et retourne les
deux autres ; c'est pourquoi les exemples utilisent les états des deux autres axes.
\end{remarque}

% ---------------------------------------------------------------------
\subsection{Porte de Hadamard}\label{sec:hadamard}
% ---------------------------------------------------------------------

\begin{definition}{Porte de Hadamard}
\[
  H=\frac{1}{\sqrt2}\begin{pmatrix}1 & 1\\ 1 & -1\end{pmatrix}
  =\ket{+}\bra{0}+\ket{-}\bra{1}.
\]
\end{definition}

\begin{align*}
  H\ket{0}&=\frac{1}{\sqrt2}\begin{pmatrix}1 & 1\\ 1 & -1\end{pmatrix}\begin{pmatrix}1\\ 0\end{pmatrix}
    =\frac{1}{\sqrt2}\begin{pmatrix}1\\ 1\end{pmatrix}=\ket{+},\\[0.6ex]
  H\ket{1}&=\frac{1}{\sqrt2}\begin{pmatrix}1 & 1\\ 1 & -1\end{pmatrix}\begin{pmatrix}0\\ 1\end{pmatrix}
    =\frac{1}{\sqrt2}\begin{pmatrix}1\\ -1\end{pmatrix}=\ket{-},\\[0.6ex]
  H\ket{+}&=\frac{1}{2}\begin{pmatrix}1 & 1\\ 1 & -1\end{pmatrix}\begin{pmatrix}1\\ 1\end{pmatrix}
    =\frac{1}{2}\begin{pmatrix}2\\ 0\end{pmatrix}=\ket{0},\qquad
  H\ket{-}=\frac{1}{2}\begin{pmatrix}0\\ 2\end{pmatrix}=\ket{1}.
\end{align*}
Comme $H\ket{+}=\ket{0}$ et $H\ket{-}=\ket{1}$, on a $H^2=\mathrm{Id}$ ($H$ est son propre
inverse). Sur la sphère de Bloch, $H$ est une rotation de $\pi$ autour de l'axe
$(x+z)/\sqrt2$ : elle échange l'axe $x$ et l'axe $z$.

\begin{remarque}[À quoi sert $H$ ?]
\begin{itemize}[nosep]
  \item Elle \emph{crée une superposition} à partir d'un état de base :
        $H\ket{0}=\ket{+}$ donne $P(0)=P(1)=\frac12$. C'est la première étape de
        la plupart des algorithmes quantiques.
  \item Elle \emph{change de base} : elle envoie la base $Z$ sur la base $X$ et
        inversement. Mesurer dans la base $X$ revient à appliquer $H$ puis à mesurer
        dans la base $Z$ (figure~\ref{fig:mesurex}).
  \item Elle est réversible ($H^2=\mathrm{Id}$).
\end{itemize}
\end{remarque}

% ---------------------------------------------------------------------
\subsection{Portes de phase}\label{sec:phase}
% ---------------------------------------------------------------------

\begin{definition}{Portes de phase $P_\theta$, $S$ et $T$}
\[
  P_\theta=\begin{pmatrix}1 & 0\\ 0 & \ee^{\ii\theta}\end{pmatrix}
  =\ket{0}\bra{0}+\ee^{\ii\theta}\ket{1}\bra{1},
  \qquad
  S=P_{\pi/2}=\begin{pmatrix}1 & 0\\ 0 & \ii\end{pmatrix},
  \qquad
  T=P_{\pi/4}=\begin{pmatrix}1 & 0\\ 0 & \dfrac{1+\ii}{\sqrt2}\end{pmatrix}.
\]
\end{definition}

Ici $\ee^{\ii\pi/2}=\ii$ et $\ee^{\ii\pi/4}=\cos\frac\pi4+\ii\sin\frac\pi4=\frac{1+\ii}{\sqrt2}$.
$P_\theta$ laisse $\ket{0}$ inchangé et multiplie $\ket{1}$ par $\ee^{\ii\theta}$ :
\[
  P_\theta\big(\alpha\ket{0}+\beta\ket{1}\big)=\alpha\ket{0}+\ee^{\ii\theta}\beta\ket{1}.
\]
Dans la paramétrisation de Bloch, $\varphi\to\varphi+\theta$ : $P_\theta$ est une rotation
d'angle $\theta$ autour de l'axe $z$. On retrouve $P_\pi=Z$, et $S^2=Z$, $T^2=S$.

\paragraph{Exemples.}
\begin{align*}
  S\ket{+}&=\tfrac{1}{\sqrt2}\big(\ket{0}+\ii\ket{1}\big)=\ket{{+}\ii},\\
  S\ket{{+}\ii}&=\tfrac{1}{\sqrt2}\big(\ket{0}+\ii\cdot\ii\ket{1}\big)=\tfrac{1}{\sqrt2}\big(\ket{0}-\ket{1}\big)=\ket{-},\\
  T\ket{+}&=\tfrac{1}{\sqrt2}\begin{pmatrix}1 & 0\\ 0 & \frac{1+\ii}{\sqrt2}\end{pmatrix}\begin{pmatrix}1\\ 1\end{pmatrix}
    =\tfrac{1}{\sqrt2}\begin{pmatrix}1\\ \frac{1+\ii}{\sqrt2}\end{pmatrix}
    =\tfrac{1}{\sqrt2}\big(\ket{0}+\ee^{\ii\pi/4}\ket{1}\big)\quad\big(\text{angle polaire }\tfrac\pi2,\ \text{azimut }\tfrac\pi4\big).
\end{align*}
$S$ tourne donc $\ket{+}$ d'un quart de tour autour de $z$ (de l'axe $x$ à l'axe $y$),
puis d'un second quart de tour jusqu'à $\ket{-}$ ; $T$ fait un huitième de tour.

\begin{remarque}[Une porte de phase est une évolution libre]
Pour $\Hop=\frac{\hbar\omega}{2}Z$, $U(t)=\ee^{-\ii\Hop t/\hbar}=\mathrm{diag}\big(\ee^{-\ii\omega t/2},\ee^{\ii\omega t/2}\big)
=\ee^{-\ii\omega t/2}\,P_{\omega t}$. Au facteur global près, laisser évoluer le qubit pendant la
durée $t$ applique la porte de phase $P_{\omega t}$ : une rotation d'angle $\omega t$ autour de
l'axe $z$.
\end{remarque}

% ---------------------------------------------------------------------
\subsection{Portes et sphère de Bloch}\label{sec:portes-bloch}
% ---------------------------------------------------------------------

Toute porte à un qubit est une rotation de la sphère de Bloch. La figure~\ref{fig:portes} montre
les axes des portes usuelles et le tableau qui suit les résume.

\begin{figure}[!ht]
\centering
\begin{tikzpicture}[x={(-1.63cm,-0.56cm)}, y={(1.63cm,-0.56cm)}, z={(0cm,2.16cm)}, >=Stealth, scale=0.9]
  % sphère et grands cercles
  \shade[ball color=insagris!15, opacity=0.55] (0,0,0) circle[radius=2.3cm];
  \begin{scope}[insagris!75, thin]
    \draw[domain=-45:135, samples=60, variable=\t] plot ({cos(\t)},{sin(\t)},0);
    \draw[domain=135:315, samples=60, variable=\t, dashed] plot ({cos(\t)},{sin(\t)},0);
    \draw[domain=-62.7:117.3, samples=60, variable=\t] plot ({cos(\t)},0,{sin(\t)});
    \draw[domain=117.3:297.3, samples=60, variable=\t, dashed] plot ({cos(\t)},0,{sin(\t)});
    \draw[domain=-62.7:117.3, samples=60, variable=\t] plot (0,{cos(\t)},{sin(\t)});
    \draw[domain=117.3:297.3, samples=60, variable=\t, dashed] plot (0,{cos(\t)},{sin(\t)});
  \end{scope}
  % axes x, y, z (portes X, Y, Z)
  \draw[insagris, dashed] (0,0,0) -- (-1,0,0);
  \draw[insagris, dashed] (0,0,0) -- (0,-1,0);
  \draw[insagris, dashed] (0,0,0) -- (0,0,-1);
  \draw[insagris] (-1,0,0) -- (-1.3,0,0);
  \draw[insagris] (0,-1,0) -- (0,-1.3,0);
  \draw[insagris] (0,0,-1) -- (0,0,-1.3);
  \draw[insagris, -{Stealth[length=2.2mm]}] (0,0,0) -- (1.3,0,0);
  \draw[insagris, -{Stealth[length=2.2mm]}] (0,0,0) -- (0,1.3,0);
  \draw[insagris, -{Stealth[length=2.2mm]}] (0,0,0) -- (0,0,1.3);
  \node[anchor=east]  at (1.33,0,0) {$x$};
  \node[anchor=west]  at (0,1.33,0) {$y$};
  \node[anchor=south] at (0,0,1.33) {$z$};
  % axe de la porte H : (x+z)/sqrt(2)
  \draw[insarouge, dashed] (0,0,0) -- (-0.707,0,-0.707);
  \draw[insarouge, very thick, -{Stealth[length=2.6mm]}] (0,0,0) -- (0.98,0,0.98);
  \node[insarouge, anchor=south east] at (0.98,0,0.98) {axe de $H$};
  % rotation autour de z (portes de phase)
  \draw[insarouge, thick, -{Stealth[length=2.2mm]}, domain=15:105, samples=30, variable=\t]
    plot ({1.15*cos(\t)},{1.15*sin(\t)},0);
  \node[insarouge, anchor=north west, font=\small] at ({1.15*cos(60)},{1.15*sin(60)},0) {$P_\theta$};
\end{tikzpicture}
\caption{Portes et rotations sur la sphère de Bloch. Les portes $X$, $Y$, $Z$ sont des
rotations de $\pi$ autour des axes $x$, $y$, $z$ ; $H$ est une rotation de $\pi$ autour de
l'axe $(x+z)/\sqrt2$ (rouge) ; $P_\theta$ (donc $Z$, $S$, $T$) tourne d'un angle $\theta$
autour de $z$.}
\label{fig:portes}
\end{figure}

\begin{center}
\renewcommand{\arraystretch}{1}
\begin{tabular}{@{}lccc@{}}
\toprule
\textbf{Porte} & \textbf{Matrice} & \textbf{Effet sur la base} & \textbf{Rotation de Bloch} \\
\midrule
$X$ & $\begin{pmatrix}0 & 1\\ 1 & 0\end{pmatrix}$ & $\ket{0}\leftrightarrow\ket{1}$ & $\pi$ autour de $x$ \\[2.2ex]
$Y$ & $\begin{pmatrix}0 & -\ii\\ \ii & 0\end{pmatrix}$ & $\ket{0}\to\ii\ket{1},\ \ket{1}\to-\ii\ket{0}$ & $\pi$ autour de $y$ \\[2.2ex]
$Z$ & $\begin{pmatrix}1 & 0\\ 0 & -1\end{pmatrix}$ & $\ket{1}\to-\ket{1},\ \ket{\pm}\leftrightarrow\ket{\mp}$ & $\pi$ autour de $z$ \\[2.2ex]
$H$ & $\dfrac{1}{\sqrt2}\begin{pmatrix}1 & 1\\ 1 & -1\end{pmatrix}$ & $\ket{0}\leftrightarrow\ket{+},\ \ket{1}\leftrightarrow\ket{-}$ & $\pi$ autour de $\frac{x+z}{\sqrt2}$ \\[2.2ex]
$S$ & $\begin{pmatrix}1 & 0\\ 0 & \ii\end{pmatrix}$ & $\ket{1}\to\ii\ket{1}$ & $\frac\pi2$ autour de $z$ \\[2.2ex]
$T$ & $\begin{pmatrix}1 & 0\\ 0 & \frac{1+\ii}{\sqrt2}\end{pmatrix}$ & $\ket{1}\to\ee^{\ii\pi/4}\ket{1}$ & $\frac\pi4$ autour de $z$ \\
\bottomrule
\end{tabular}
\end{center}

\begin{figure}[!ht]
\centering
\begin{tikzpicture}[x={(-1.63cm,-0.56cm)}, y={(1.63cm,-0.56cm)}, z={(0cm,2.16cm)}, >=Stealth, scale=0.9]
  % sphère et grands cercles
  \shade[ball color=insagris!15, opacity=0.55] (0,0,0) circle[radius=2.3cm];
  \begin{scope}[insagris!75, thin]
    \draw[domain=-45:135, samples=60, variable=\t] plot ({cos(\t)},{sin(\t)},0);
    \draw[domain=135:315, samples=60, variable=\t, dashed] plot ({cos(\t)},{sin(\t)},0);
    \draw[domain=-62.7:117.3, samples=60, variable=\t] plot ({cos(\t)},0,{sin(\t)});
    \draw[domain=117.3:297.3, samples=60, variable=\t, dashed] plot ({cos(\t)},0,{sin(\t)});
    \draw[domain=-62.7:117.3, samples=60, variable=\t] plot (0,{cos(\t)},{sin(\t)});
    \draw[domain=117.3:297.3, samples=60, variable=\t, dashed] plot (0,{cos(\t)},{sin(\t)});
  \end{scope}
  % axes x, y, z (portes X, Y, Z)
  \draw[insagris, dashed] (0,0,0) -- (-1,0,0);
  \draw[insagris, dashed] (0,0,0) -- (0,-1,0);
  \draw[insagris, dashed] (0,0,0) -- (0,0,-1);
  \draw[insagris] (-1,0,0) -- (-1.3,0,0);
  \draw[insagris] (0,-1,0) -- (0,-1.3,0);
  \draw[insagris] (0,0,-1) -- (0,0,-1.3);
  \draw[insagris, -{Stealth[length=2.2mm]}] (0,0,0) -- (1.3,0,0);
  \draw[insagris, -{Stealth[length=2.2mm]}] (0,0,0) -- (0,1.3,0);
  \draw[insagris, -{Stealth[length=2.2mm]}] (0,0,0) -- (0,0,1.3);
  % trajet : |0> --H--> |+> --S--> |+i>
  \draw[insarouge, very thick, -{Stealth[length=2.6mm]}, domain=3:87, samples=30, variable=\t]
    plot ({sin(\t)},0,{cos(\t)});
  \draw[insarouge, very thick, -{Stealth[length=2.6mm]}, domain=3:87, samples=30, variable=\t]
    plot ({cos(\t)},{sin(\t)},0);
  \foreach \ax/\ay/\az in {0/0/1, 1/0/0, 0/1/0} \fill[insarouge] (\ax,\ay,\az) circle[radius=2pt];
  \node[insarouge, anchor=south] at (0,0,1.33) {$\ket{0}$};
  \node[insarouge, anchor=east] at (1.33,0,0) {$\ket{+}$};
  \node[insarouge, anchor=west] at (0,1.33,0) {$\ket{{+}\ii}$};
  \node[insarouge, font=\small, anchor=south west] at ({sin(45)},0,{cos(45)}) {$H$};
  \node[insarouge, font=\small, anchor=north west] at ({cos(45)},{sin(45)},0) {$S$};
\end{tikzpicture}
\caption{Trajet d'un état sous le circuit $H$ puis $S$ (flèches schématiques). $H$ envoie
$\ket{0}$ sur $\ket{+}$ ; $S$, rotation d'un quart de tour autour de $z$, envoie $\ket{+}$ sur
$\ket{{+}\ii}$. Les trois états sont alignés sur les axes $z$, $x$ et $y$.}
\label{fig:trajet}
\end{figure}

% ---------------------------------------------------------------------
\subsection{Circuits quantiques}\label{sec:circuits}
% ---------------------------------------------------------------------

Un \emph{circuit} représente un calcul : chaque qubit est un fil horizontal, le temps va de gauche
à droite, et chaque porte est un symbole posé sur les fils qu'elle concerne. La
figure~\ref{fig:symboles} donne les symboles utilisés dans ce document.

\begin{figure}[!ht]
\centering
\begin{adjustbox}{max width=\linewidth}
\begin{tikzpicture}[>=Stealth, font=\small,
    nom/.style={font=\footnotesize, align=center, anchor=north}]
  % (a) porte à un qubit
  \draw (0,0) -- (1.8,0);
  \node[porte] at (0.9,0) {$U$};
  \node[nom] at (0.9,-1.15) {porte $U$\\ à un qubit};
  % (b) CNOT
  \begin{scope}[xshift=3.2cm]
    \draw (0,0.6) -- (1.8,0.6);
    \draw (0,-0.6) -- (1.8,-0.6);
    \draw[insarouge, thick] (0.9,0.6) -- (0.9,-0.6);
    \pic at (0.9,0.6) {ctl};
    \pic at (0.9,-0.6) {cible};
    \node[nom] at (0.9,-1.15) {CNOT\\ (contrôle $\bullet$, cible $\oplus$)};
  \end{scope}
  % (c) SWAP
  \begin{scope}[xshift=6.4cm]
    \draw (0,0.6) -- (1.8,0.6);
    \draw (0,-0.6) -- (1.8,-0.6);
    \draw[insarouge, thick] (0.9,0.6) -- (0.9,-0.6);
    \pic at (0.9,0.6) {croix};
    \pic at (0.9,-0.6) {croix};
    \node[nom] at (0.9,-1.15) {SWAP\\ (échange)};
  \end{scope}
  % (d) contrôlé-U
  \begin{scope}[xshift=9.6cm]
    \draw (0,0.6) -- (1.8,0.6);
    \draw (0,-0.6) -- (1.8,-0.6);
    \draw[insarouge, thick] (0.9,-0.6) -- (0.9,0.3);
    \pic at (0.9,-0.6) {ctl};
    \node[porte] at (0.9,0.6) {$U$};
    \node[nom] at (0.9,-1.15) {contrôlé-$U$\\ (contrôle en bas)};
  \end{scope}
  % (e) mesure
  \begin{scope}[xshift=12.8cm]
    \draw (0,0) -- (0.6,0);
    \pic at (0.9,0) {mesure};
    \draw (1.2,0.03) -- (2.0,0.03);
    \draw (1.2,-0.03) -- (2.0,-0.03);
    \node[nom] at (1.0,-1.15) {mesure\\ (fil double : bit)};
  \end{scope}
\end{tikzpicture}
\end{adjustbox}
\caption{Symboles des circuits quantiques. Le point noir $\bullet$ marque un qubit de contrôle,
le symbole $\oplus$ une cible qui subit $X$, les deux croix un échange. Après une mesure, le fil
double transporte un bit classique.}
\label{fig:symboles}
\end{figure}

\begin{remarque}[Conventions de lecture]
\begin{itemize}[nosep]
  \item Le temps va de gauche à droite. Deux portes successives $U_1$ puis $U_2$ se multiplient
        dans l'ordre inverse : la matrice du circuit est $U_2U_1$, car $U_1$ agit d'abord sur le ket.
  \item Avec plusieurs fils, le fil du \emph{haut} porte le \emph{dernier} chiffre du ket et le
        fil du \emph{bas} le premier : si $A$ est le premier qubit et $B$ le second,
        $\ket{a,b}=\ket{a}_A\otimes\ket{b}_B$ et $B$ est dessiné au-dessus de $A$.
  \item Deux portes placées dans la même colonne, sur deux fils, agissent en parallèle : c'est un
        produit tensoriel $U_A\otimes U_B$ (section~\ref{sec:op2}).
\end{itemize}
\end{remarque}

\paragraph{Exemple.} Le circuit de la figure~\ref{fig:circuit-hs} applique $H$ puis $S$ à
$\ket{0}$. Sa matrice est $SH$ :
\[
  SH=\begin{pmatrix}1&0\\0&\ii\end{pmatrix}\frac{1}{\sqrt2}\begin{pmatrix}1&1\\1&-1\end{pmatrix}
  =\frac{1}{\sqrt2}\begin{pmatrix}1&1\\ \ii&-\ii\end{pmatrix},
  \qquad
  SH\ket{0}=\frac{1}{\sqrt2}\begin{pmatrix}1\\ \ii\end{pmatrix}=\ket{{+}\ii}.
\]
Le circuit prépare donc l'état $\ket{{+}\ii}$ de la base $Y$ à partir de $\ket{0}$
(figure~\ref{fig:trajet}). L'ordre compte : $HS\ket{0}=H\ket{0}=\ket{+}$, car $S\ket{0}=\ket{0}$.

\begin{figure}[!ht]
\centering
\begin{tikzpicture}[>=Stealth, font=\small]
  \draw (0,0) -- (6.6,0);
  \node[anchor=east] at (0,0) {$\ket{0}$};
  \node[porte] at (2.0,0) {$H$};
  \node[porte] at (4.3,0) {$S$};
  \node[anchor=west] at (6.6,0) {$\ket{{+}\ii}$};
  \foreach \x in {0.7, 3.15, 5.5} \draw[dashed, insagris!70] (\x,0.6) -- (\x,-0.95);
  \node[font=\footnotesize, anchor=north] at (0.7,-0.95) {$\ket{0}$};
  \node[font=\footnotesize, anchor=north] at (3.15,-0.95) {$\ket{+}$};
  \node[font=\footnotesize, anchor=north] at (5.5,-0.95) {$\ket{{+}\ii}$};
\end{tikzpicture}
\caption{Circuit $H$ puis $S$ appliqué à $\ket{0}$. Les états intermédiaires sont indiqués sous
les traits pointillés ; la matrice du circuit est $SH$.}
\label{fig:circuit-hs}
\end{figure}

\paragraph{Mesure.} Le symbole de mesure (compteur) donne le résultat $0$ ou $1$ dans la base $Z$.
Mesurer dans la base $X$ revient à appliquer $H$ puis à mesurer dans la base $Z$
(figure~\ref{fig:mesurex}) : comme $H\ket{\pm}=\ket{0}$ ou $\ket{1}$, le résultat $0$ correspond à
$\ket{+}$ et le résultat $1$ à $\ket{-}$, avec $P(\pm)=|\braket{\pm}{\Psi}|^2$.

\begin{figure}[!ht]
\centering
\begin{tikzpicture}[>=Stealth, font=\small]
  % (a) base Z
  \draw (0,0) -- (1.2,0);
  \node[anchor=east] at (0,0) {$\ket{\Psi}$};
  \pic at (1.5,0) {mesure};
  \draw (1.8,0.03) -- (2.6,0.03);
  \draw (1.8,-0.03) -- (2.6,-0.03);
  \node[anchor=west] at (2.7,0) {$0$ ou $1$};
  \node[font=\footnotesize] at (1.5,-0.95) {(a) mesure en base $Z$};
  % (b) base X
  \begin{scope}[xshift=7.5cm]
    \draw (0,0) -- (2.6,0);
    \node[anchor=east] at (0,0) {$\ket{\Psi}$};
    \node[porte] at (1.0,0) {$H$};
    \pic at (2.3,0) {mesure};
    \draw (2.6,0.03) -- (3.4,0.03);
    \draw (2.6,-0.03) -- (3.4,-0.03);
    \node[anchor=west] at (3.5,0) {$+$ ou $-$};
    \node[font=\footnotesize] at (1.9,-0.95) {(b) mesure en base $X$};
  \end{scope}
\end{tikzpicture}
\caption{(a) Mesure dans la base $Z$. (b) Mesure dans la base $X$ : porte de Hadamard, puis
mesure dans la base $Z$ ; les résultats $0$ et $1$ se lisent $+$ et $-$.}
\label{fig:mesurex}
\end{figure}

% ---------------------------------------------------------------------
\subsection{Portes locales sur deux qubits}\label{sec:op2}
% ---------------------------------------------------------------------

Une porte appliquée à un seul qubit d'un système de deux s'écrit avec un produit tensoriel
(section~\ref{sec:tenseur}) : $U\otimes I$ agit sur $A$ (premier facteur, fil du bas), $I\otimes U$
agit sur $B$, et $U_A\otimes U_B$ agit sur les deux en parallèle (figure~\ref{fig:local}).

\begin{figure}[!ht]
\centering
\begin{tikzpicture}[>=Stealth, font=\small]
  \foreach \xs/\haut/\bas/\tit in {0/{}/{$X$}/{(a) $X\otimes I$},
                                   5.4/{$X$}/{}/{(b) $I\otimes X$},
                                   10.8/{$H$}/{$H$}/{(c) $H\otimes H$}} {
    \begin{scope}[xshift=\xs cm]
      \draw (0,0.75) -- (2.6,0.75);
      \draw (0,0) -- (2.6,0);
      \node[anchor=east] at (0,0.75) {$B$};
      \node[anchor=east] at (0,0) {$A$};
      \ifx\haut\empty\else \node[porte] at (1.3,0.75) {\haut}; \fi
      \ifx\bas\empty\else \node[porte] at (1.3,0) {\bas}; \fi
      \node[font=\footnotesize] at (1.3,-0.85) {\tit};
    \end{scope}
  }
\end{tikzpicture}
\caption{Portes en parallèle sur deux qubits. Le qubit $A$, premier facteur du produit tensoriel,
est sur le fil du bas. (a) $X$ sur $A$ seulement. (b) $X$ sur $B$ seulement. (c) $H$ sur les deux.}
\label{fig:local}
\end{figure}

\subsubsection{Portes agissant sur un seul des deux qubits}

On note $I=\mathrm{Id}$ ($2\times2$). $X\otimes I$ agit sur $A$, $I\otimes X$ agit sur $B$ :
\begin{align*}
  I\otimes X&=\begin{pmatrix}1\cdot X & 0\cdot X\\ 0\cdot X & 1\cdot X\end{pmatrix}
    =\begin{pmatrix}0 & 1 & 0 & 0\\ 1 & 0 & 0 & 0\\ 0 & 0 & 0 & 1\\ 0 & 0 & 1 & 0\end{pmatrix},\\[1ex]
  X\otimes I&=\begin{pmatrix}0\cdot I & 1\cdot I\\ 1\cdot I & 0\cdot I\end{pmatrix}
    =\begin{pmatrix}0 & 0 & 1 & 0\\ 0 & 0 & 0 & 1\\ 1 & 0 & 0 & 0\\ 0 & 1 & 0 & 0\end{pmatrix}.
\end{align*}
Sur $\ket{00}=(1,0,0,0)^{\mathsf T}$ (première colonne de la matrice) :
$(I\otimes X)\ket{00}=(0,1,0,0)^{\mathsf T}=\ket{01}$ et
$(X\otimes I)\ket{00}=(0,0,1,0)^{\mathsf T}=\ket{10}$. Avec la règle :
$(I\otimes X)\ket{00}=I\ket{0}\otimes X\ket{0}=\ket{0}\otimes\ket{1}=\ket{01}$.

Pour la porte de Hadamard, avec $H=\frac{1}{\sqrt2}\begin{pmatrix}1 & 1\\ 1 & -1\end{pmatrix}$ :
\[
  H\otimes I=\frac{1}{\sqrt2}\begin{pmatrix}1\cdot I & 1\cdot I\\ 1\cdot I & -1\cdot I\end{pmatrix}
  =\frac{1}{\sqrt2}\begin{pmatrix}1 & 0 & 1 & 0\\ 0 & 1 & 0 & 1\\ 1 & 0 & -1 & 0\\ 0 & 1 & 0 & -1\end{pmatrix}.
\]
Ses colonnes donnent l'image des quatre états de base :
\begin{align*}
  (H\otimes I)\ket{00}&=\tfrac{1}{\sqrt2}(1,0,1,0)^{\mathsf T}=\ket{+}\otimes\ket{0},&
  (H\otimes I)\ket{01}&=\tfrac{1}{\sqrt2}(0,1,0,1)^{\mathsf T}=\ket{+}\otimes\ket{1},\\
  (H\otimes I)\ket{10}&=\tfrac{1}{\sqrt2}(1,0,-1,0)^{\mathsf T}=\ket{-}\otimes\ket{0},&
  (H\otimes I)\ket{11}&=\tfrac{1}{\sqrt2}(0,1,0,-1)^{\mathsf T}=\ket{-}\otimes\ket{1}.
\end{align*}

\subsubsection{Hadamard sur les deux qubits}

\[
  H\otimes H=\frac12\begin{pmatrix}1\cdot H' & 1\cdot H'\\ 1\cdot H' & -1\cdot H'\end{pmatrix},
  \quad H'=\begin{pmatrix}1 & 1\\ 1 & -1\end{pmatrix},
  \qquad
  H\otimes H=\frac12\begin{pmatrix}1 & 1 & 1 & 1\\ 1 & -1 & 1 & -1\\ 1 & 1 & -1 & -1\\ 1 & -1 & -1 & 1\end{pmatrix}.
\]
La première colonne donne $(H\otimes H)\ket{00}=\frac12(1,1,1,1)^{\mathsf T}=\ket{+}\otimes\ket{+}$
(cohérent avec $H\ket{0}\otimes H\ket{0}$) : les quatre résultats $00,01,10,11$ sont
équiprobables, chacun avec la probabilité $\frac14$. Sur $n$ qubits, $H^{\otimes n}\ket{0\ldots0}$
donne une superposition uniforme de $2^n$ états.

% ---------------------------------------------------------------------
\subsection{Porte CNOT}\label{sec:cnot}
% ---------------------------------------------------------------------

\begin{definition}{Porte CNOT (NOT contrôlé)}
La porte CNOT agit sur deux qubits : le qubit $B$ est le \emph{contrôle}, le qubit $A$ est la
\emph{cible}. La cible est inversée (porte $X$) si le contrôle vaut $1$ et reste inchangée
s'il vaut $0$. Avec $\ket{a,b}=\ket{a}_A\otimes\ket{b}_B$ :
\[
  \mathrm{CNOT}\,\ket{a,b}=\ket{a\oplus b,\,b},
  \qquad
  \mathrm{CNOT}=I\otimes\ket{0}\bra{0}+X\otimes\ket{1}\bra{1},
\]
où $\oplus$ est l'addition modulo $2$ ($0\oplus0=1\oplus1=0$ et $0\oplus1=1\oplus0=1$).
\end{definition}

Sur le circuit de la figure~\ref{fig:cnot}, le point noir posé sur le fil de $B$ (en haut) marque
le contrôle et le symbole $\oplus$ posé sur le fil de $A$ (en bas) marque la cible.

\begin{figure}[!ht]
\centering
\begin{tikzpicture}[>=Stealth, font=\small]
  \draw (0,1.2) -- (6,1.2);
  \draw (0,0) -- (6,0);
  \node[anchor=east] at (0,1.2) {$\ket{b}_B$};
  \node[anchor=east] at (0,0) {$\ket{a}_A$};
  \node[anchor=west] at (6,1.2) {$\ket{b}_B$};
  \node[anchor=west] at (6,0) {$\ket{a\oplus b}_A$};
  \draw[insarouge, thick] (3,1.2) -- (3,0);
  \pic at (3,1.2) {ctl};
  \pic at (3,0) {cible};
  \node[anchor=south, text=insagris] at (3,1.5) {qubit de contrôle};
  \node[anchor=north, text=insagris] at (3,-0.5) {qubit cible};
\end{tikzpicture}
\caption{Porte CNOT : le fil du haut ($B$) est le contrôle, le fil du bas ($A$) est la cible. Le
ket s'écrit $\ket{a,b}$ : le fil du haut donne le dernier chiffre.}
\label{fig:cnot}
\end{figure}

\subsubsection{Table de vérité}

\begin{center}
\renewcommand{\arraystretch}{1.3}
\begin{tabular}{@{}cccc@{}}
\toprule
Entrée $\ket{a,b}$ & Contrôle $b$ & Sortie $\ket{a\oplus b,b}$ & Effet sur la cible $A$ \\
\midrule
$\ket{00}$ & $0$ & $\ket{00}$ & aucun \\
$\ket{01}$ & $1$ & $\ket{11}$ & $0\to1$ \\
$\ket{10}$ & $0$ & $\ket{10}$ & aucun \\
$\ket{11}$ & $1$ & $\ket{01}$ & $1\to0$ \\
\bottomrule
\end{tabular}
\end{center}

\subsubsection{Matrice du CNOT}

La $j$-ième colonne d'une matrice $4\times4$ est l'image du $j$-ième vecteur de base
($\ket{00},\ket{01},\ket{10},\ket{11}$, dans cet ordre). La table donne :
\begin{align*}
  \mathrm{CNOT}\ket{00}&=\ket{00}=\begin{pmatrix}1\\0\\0\\0\end{pmatrix},&
  \mathrm{CNOT}\ket{01}&=\ket{11}=\begin{pmatrix}0\\0\\0\\1\end{pmatrix},\\[1.5ex]
  \mathrm{CNOT}\ket{10}&=\ket{10}=\begin{pmatrix}0\\0\\1\\0\end{pmatrix},&
  \mathrm{CNOT}\ket{11}&=\ket{01}=\begin{pmatrix}0\\1\\0\\0\end{pmatrix},
\end{align*}
d'où, en juxtaposant ces quatre colonnes,
\[
  \mathrm{CNOT}=\begin{pmatrix}
    1 & 0 & 0 & 0\\ 0 & 0 & 0 & 1\\ 0 & 0 & 1 & 0\\ 0 & 1 & 0 & 0
  \end{pmatrix}.
\]

On retrouve cette matrice avec la formule $I\otimes\ket{0}\bra{0}+X\otimes\ket{1}\bra{1}$ et
la règle $A\otimes B=(A_{ij}B)$ (section~\ref{sec:tenseur}), où
$\ket{0}\bra{0}=\begin{pmatrix}1&0\\0&0\end{pmatrix}$ et
$\ket{1}\bra{1}=\begin{pmatrix}0&0\\0&1\end{pmatrix}$ :
\[
  I\otimes\ket{0}\bra{0}=\begin{pmatrix}1&0&0&0\\0&0&0&0\\0&0&1&0\\0&0&0&0\end{pmatrix},
  \qquad
  X\otimes\ket{1}\bra{1}=\begin{pmatrix}0&0&0&0\\0&0&0&1\\0&0&0&0\\0&1&0&0\end{pmatrix},
\]
et la somme de ces deux matrices est la matrice du CNOT.

\subsubsection{Action sur un état quelconque}

Pour $\ket{\Psi}=c_{00}\ket{00}+c_{01}\ket{01}+c_{10}\ket{10}+c_{11}\ket{11}$ :
\[
  \mathrm{CNOT}\begin{pmatrix}c_{00}\\c_{01}\\c_{10}\\c_{11}\end{pmatrix}
  =\begin{pmatrix}
    1\cdot c_{00}+0\cdot c_{01}+0\cdot c_{10}+0\cdot c_{11}\\
    0\cdot c_{00}+0\cdot c_{01}+0\cdot c_{10}+1\cdot c_{11}\\
    0\cdot c_{00}+0\cdot c_{01}+1\cdot c_{10}+0\cdot c_{11}\\
    0\cdot c_{00}+1\cdot c_{01}+0\cdot c_{10}+0\cdot c_{11}
  \end{pmatrix}
  =\begin{pmatrix}c_{00}\\c_{11}\\c_{10}\\c_{01}\end{pmatrix}.
\]
Le CNOT échange les amplitudes de $\ket{01}$ et de $\ket{11}$, c'est-à-dire des deux états
dont le contrôle vaut $1$, et laisse les autres.

\paragraph{Exemple : contrôle en superposition.} Soit $\ket{0}_A\otimes\ket{+}_B
=\frac{1}{\sqrt2}\big(\ket{00}+\ket{01}\big)=\frac{1}{\sqrt2}(1,1,0,0)^{\mathsf T}$,
un état séparable. Alors
\[
  \mathrm{CNOT}\,\frac{1}{\sqrt2}\begin{pmatrix}1\\1\\0\\0\end{pmatrix}
  =\frac{1}{\sqrt2}\begin{pmatrix}1\\0\\0\\1\end{pmatrix}
  =\frac{1}{\sqrt2}\big(\ket{00}+\ket{11}\big).
\]
Le résultat est intriqué ($c_{00}c_{11}-c_{01}c_{10}=\frac12\neq0$ ; section~\ref{sec:separable}) :
un CNOT dont le contrôle est en superposition crée de l'intrication. C'est l'état de Bell
$\ket{\Phi^+}$ du chapitre~\ref{chap:bell}.

\paragraph{Réversibilité.} Deux CNOT successifs redonnent l'état initial :
$\ket{a,b}\to\ket{a\oplus b,b}\to\ket{a\oplus b\oplus b,b}=\ket{a,b}$, donc
$\mathrm{CNOT}^2=\mathrm{Id}$. La matrice est réelle et symétrique, donc
$\mathrm{CNOT}^\dagger\mathrm{CNOT}=\mathrm{CNOT}^2=\mathrm{Id}$ : elle est bien unitaire.

\begin{remarque}[CNOT et copie d'un bit]
Avec la cible dans l'état $\ket{0}$, $\mathrm{CNOT}\ket{0,b}=\ket{b,b}$ pour $b=0$ et $b=1$ : le CNOT
copie le bit classique $b$. En revanche $\mathrm{CNOT}\big(\ket{0}\otimes\ket{+}\big)$ est l'état
ci-dessus, et non $\ket{+}\otimes\ket{+}$ : un état quantique inconnu ne peut pas être copié
(théorème de non-clonage, section~\ref{sec:teleportation}).
\end{remarque}

\begin{remarque}[Contrôle en dernier ou en premier]
La matrice dépend de l'ordre d'écriture du ket. Notons $\mathrm{CNOT}_{21}$ le CNOT dont le
contrôle est le \emph{second} chiffre du ket : c'est la matrice ci-dessus, adoptée dans ce
document (fil du haut). Beaucoup d'ouvrages mettent le contrôle en \emph{premier} ; le CNOT
correspondant est
\[
  \mathrm{CNOT}_{12}=\ket{0}\bra{0}\otimes I+\ket{1}\bra{1}\otimes X
  =\begin{pmatrix}1&0&0&0\\0&1&0&0\\0&0&0&1\\0&0&1&0\end{pmatrix}.
\]
Le même circuit a donc deux matrices, selon que le contrôle est écrit en dernier ou en premier.
Les deux sont reliées par le SWAP (section~\ref{sec:swap}). La section~\ref{sec:controlees} sur les
portes contrôlées met le contrôle en premier, avec la matrice $\mathrm{diag}(I,U)$.
\end{remarque}

% ---------------------------------------------------------------------
\subsection{Porte SWAP}\label{sec:swap}
% ---------------------------------------------------------------------

\begin{definition}{Porte SWAP (échange)}
La porte SWAP échange les états des deux qubits :
\[
  \mathrm{SWAP}\big(\ket{\psi}\otimes\ket{\phi}\big)=\ket{\phi}\otimes\ket{\psi}.
\]
Sur les états de base, $\mathrm{SWAP}\ket{a,b}=\ket{b,a}$ ; on prolonge par linéarité.
\end{definition}

Son symbole est formé de deux croix reliées par un trait vertical
(figure~\ref{fig:swap}).

\begin{figure}[!ht]
\centering
\begin{tikzpicture}[>=Stealth, font=\small]
  \draw (0,1.2) -- (6,1.2);
  \draw (0,0) -- (6,0);
  \node[anchor=east] at (0,1.2) {$\ket{a}$};
  \node[anchor=east] at (0,0) {$\ket{b}$};
  \node[anchor=west] at (6,1.2) {$\ket{b}$};
  \node[anchor=west] at (6,0) {$\ket{a}$};
  \draw[insarouge, thick] (3,1.2) -- (3,0);
  \pic at (3,1.2) {croix};
  \pic at (3,0) {croix};
\end{tikzpicture}
\caption{Porte SWAP : les deux qubits échangent leurs états.}
\label{fig:swap}
\end{figure}

\subsubsection{Table et matrice}

\begin{center}
\renewcommand{\arraystretch}{1.3}
\begin{tabular}{@{}cc@{}}
\toprule
Entrée & Sortie \\
\midrule
$\ket{00}$ & $\ket{00}$ \\
$\ket{01}$ & $\ket{10}$ \\
$\ket{10}$ & $\ket{01}$ \\
$\ket{11}$ & $\ket{11}$ \\
\bottomrule
\end{tabular}
\end{center}

Les images des quatre vecteurs de base sont les colonnes de la matrice :
\[
  \mathrm{SWAP}\ket{00}=\begin{pmatrix}1\\0\\0\\0\end{pmatrix},\quad
  \mathrm{SWAP}\ket{01}=\begin{pmatrix}0\\0\\1\\0\end{pmatrix},\quad
  \mathrm{SWAP}\ket{10}=\begin{pmatrix}0\\1\\0\\0\end{pmatrix},\quad
  \mathrm{SWAP}\ket{11}=\begin{pmatrix}0\\0\\0\\1\end{pmatrix},
\]
\[
  \mathrm{SWAP}=\begin{pmatrix}
    1 & 0 & 0 & 0\\ 0 & 0 & 1 & 0\\ 0 & 1 & 0 & 0\\ 0 & 0 & 0 & 1
  \end{pmatrix}.
\]
Appliquée à $\ket{01}=(0,1,0,0)^{\mathsf T}$, cette matrice sélectionne sa deuxième colonne :
\[
  \begin{pmatrix}
    1 & 0 & 0 & 0\\ 0 & 0 & 1 & 0\\ 0 & 1 & 0 & 0\\ 0 & 0 & 0 & 1
  \end{pmatrix}
  \begin{pmatrix}0\\1\\0\\0\end{pmatrix}
  =\begin{pmatrix}0\\0\\1\\0\end{pmatrix}=\ket{10}.
\]
Sur un état quelconque, $\mathrm{SWAP}\,(c_{00},c_{01},c_{10},c_{11})^{\mathsf T}
=(c_{00},c_{10},c_{01},c_{11})^{\mathsf T}$ : les amplitudes de $\ket{01}$ et $\ket{10}$
sont échangées.

\subsubsection{Accord avec la définition}

Pour deux états $\ket{\psi}=(a_0,a_1)^{\mathsf T}$ et $\ket{\phi}=(b_0,b_1)^{\mathsf T}$ :
\[
  \mathrm{SWAP}\,(\ket{\psi}\otimes\ket{\phi})
  =\mathrm{SWAP}\begin{pmatrix}a_0b_0\\a_0b_1\\a_1b_0\\a_1b_1\end{pmatrix}
  =\begin{pmatrix}a_0b_0\\a_1b_0\\a_0b_1\\a_1b_1\end{pmatrix}
  =\begin{pmatrix}b_0\\b_1\end{pmatrix}\otimes\begin{pmatrix}a_0\\a_1\end{pmatrix}
  =\ket{\phi}\otimes\ket{\psi},
\]
ce qui est bien la définition. La porte s'écrit aussi avec des projecteurs :
\[
  \mathrm{SWAP}=\sum_{a,b\in\{0,1\}}\ket{a}\bra{b}\otimes\ket{b}\bra{a}.
\]
En effet, pour $\ket{c}\otimes\ket{d}$ on trouve
$\sum_{a,b}\ket{a}\braket{b}{c}\otimes\ket{b}\braket{a}{d}$, et seul le terme $b=c$, $a=d$
est non nul : le résultat est $\ket{d}\otimes\ket{c}$.

\paragraph{Exemples.}
\begin{itemize}[nosep]
  \item Deux états différents : $\mathrm{SWAP}\big(\ket{0}\otimes\ket{+}\big)
        =\mathrm{SWAP}\,\frac{1}{\sqrt2}(1,1,0,0)^{\mathsf T}=\frac{1}{\sqrt2}(1,0,1,0)^{\mathsf T}
        =\ket{+}\otimes\ket{0}$.
  \item L'état $\ket{\psi}=\frac{1}{\sqrt2}\ket{00}+\frac{1}{\sqrt2}\ket{11}$ ne change pas :
        $\mathrm{SWAP}\ket{\psi}=\frac{1}{\sqrt2}\ket{00}+\frac{1}{\sqrt2}\ket{11}=\ket{\psi}$,
        car $\ket{00}$ et $\ket{11}$ sont invariants. De même
        $\frac{1}{\sqrt2}(\ket{00}-\ket{11})$ et $\frac{1}{\sqrt2}(\ket{01}+\ket{10})$ sont invariants,
        alors que $\mathrm{SWAP}\,\frac{1}{\sqrt2}(\ket{01}-\ket{10})=\frac{1}{\sqrt2}(\ket{10}-\ket{01})$
        est l'opposé de l'état de départ. Ce sont les quatre états de Bell (chapitre~\ref{chap:bell}).
\end{itemize}

\paragraph{Propriétés.} Échanger deux fois redonne l'état de départ :
$\mathrm{SWAP}^2=\mathrm{Id}$ ; la matrice est réelle et symétrique, donc
$\mathrm{SWAP}^\dagger\mathrm{SWAP}=\mathrm{SWAP}^2=\mathrm{Id}$ : elle est unitaire.

\subsubsection{Lien avec le CNOT}

\begin{resultat}{SWAP échange les rôles de contrôle et de cible}
\[
  \mathrm{CNOT}_{21}=\mathrm{SWAP}\;\mathrm{CNOT}_{12}\;\mathrm{SWAP}.
\]
\end{resultat}

\emph{Sur les états de base.}
$\mathrm{SWAP}\,\mathrm{CNOT}_{12}\,\mathrm{SWAP}\ket{a,b}
=\mathrm{SWAP}\,\mathrm{CNOT}_{12}\ket{b,a}
=\mathrm{SWAP}\ket{b,\,b\oplus a}=\ket{b\oplus a,\,b}=\mathrm{CNOT}_{21}\ket{a,b}$.

\emph{Par les matrices.} Multiplier à gauche par $\mathrm{SWAP}$ échange les lignes $2$ et
$3$, multiplier à droite l'échange les colonnes $2$ et $3$ :
\[
  \mathrm{SWAP}\,\mathrm{CNOT}_{12}
  =\begin{pmatrix}1&0&0&0\\0&0&0&1\\0&1&0&0\\0&0&1&0\end{pmatrix},
  \qquad
  \mathrm{SWAP}\,\mathrm{CNOT}_{12}\,\mathrm{SWAP}
  =\begin{pmatrix}1&0&0&0\\0&0&0&1\\0&0&1&0\\0&1&0&0\end{pmatrix}=\mathrm{CNOT}_{21}.
\]

\begin{remarque}[SWAP avec trois CNOT]
$\mathrm{SWAP}=\mathrm{CNOT}_{12}\,\mathrm{CNOT}_{21}\,\mathrm{CNOT}_{12}$. En effet, sur
$\ket{a,b}$ (on applique les portes de droite à gauche) :
\[
  \ket{a,b}\ \xrightarrow{\ \mathrm{CNOT}_{12}\ }\ \ket{a,\,a\oplus b}
  \ \xrightarrow{\ \mathrm{CNOT}_{21}\ }\ \ket{a\oplus(a\oplus b),\,a\oplus b}=\ket{b,\,a\oplus b}
  \ \xrightarrow{\ \mathrm{CNOT}_{12}\ }\ \ket{b,\,b\oplus(a\oplus b)}=\ket{b,a}.
\]
Les deux membres coïncident sur les états de base, donc sur tous les états par linéarité.
\end{remarque}

% ---------------------------------------------------------------------
\subsection{Portes contrôlées : contrôlé-\texorpdfstring{$U$}{U}, Fredkin et Toffoli}\label{sec:controlees}
% ---------------------------------------------------------------------

Le CNOT est le cas le plus simple d'une idée générale : appliquer une opération à un ou
plusieurs qubits \emph{cibles} seulement si un ou plusieurs qubits de \emph{contrôle} valent $1$.
Dans cette section, le contrôle est le \emph{premier} chiffre
du ket (remarque de la section~\ref{sec:cnot}). Les qubits sont notés $x$, $y$, $z$ (le qubit $x$ est le contrôle) et
$\ket{x,y,z}=\ket{x}\otimes\ket{y}\otimes\ket{z}$ ; d'après la convention de la section~\ref{sec:circuits}, le
contrôle est dessiné en bas (figure~\ref{fig:controlees}). Pour éviter la confusion avec les
matrices de Pauli, on écrit ici $\sigma_x=X$ pour la porte NOT.

\begin{figure}[!ht]
\centering
\begin{adjustbox}{max width=\linewidth}
\begin{tikzpicture}[>=Stealth, font=\small]
  % (a) contrôlé-U : contrôle x en bas, cible y en haut
  \draw (0,1.2) -- (3.4,1.2);
  \draw (0,0) -- (3.4,0);
  \node[anchor=east] at (0,1.2) {$y$};
  \node[anchor=east] at (0,0) {$x$};
  \draw[insarouge, thick] (1.7,0) -- (1.7,0.85);
  \pic at (1.7,0) {ctl};
  \node[porte] at (1.7,1.2) {$U$};
  \node[anchor=north, font=\footnotesize] at (1.7,-0.5) {(a) contrôlé-$U$};
  % (b) contrôlé-SWAP (Fredkin) : x contrôle, y et z échangés
  \draw (5,2.4) -- (8.4,2.4);
  \draw (5,1.2) -- (8.4,1.2);
  \draw (5,0) -- (8.4,0);
  \node[anchor=east] at (5,2.4) {$z$};
  \node[anchor=east] at (5,1.2) {$y$};
  \node[anchor=east] at (5,0) {$x$};
  \draw[insarouge, thick] (6.7,0) -- (6.7,2.4);
  \pic at (6.7,0) {ctl};
  \pic at (6.7,1.2) {croix};
  \pic at (6.7,2.4) {croix};
  \node[anchor=north, font=\footnotesize] at (6.7,-0.5) {(b) Fredkin (C-SWAP)};
  % (c) Toffoli (CCNOT) : x et y contrôlent, z cible
  \draw (10,2.4) -- (13.4,2.4);
  \draw (10,1.2) -- (13.4,1.2);
  \draw (10,0) -- (13.4,0);
  \node[anchor=east] at (10,2.4) {$z$};
  \node[anchor=east] at (10,1.2) {$y$};
  \node[anchor=east] at (10,0) {$x$};
  \draw[insarouge, thick] (11.7,0) -- (11.7,2.4);
  \pic at (11.7,0) {ctl};
  \pic at (11.7,1.2) {ctl};
  \pic at (11.7,2.4) {cible};
  \node[anchor=north, font=\footnotesize] at (11.7,-0.5) {(c) Toffoli (CCNOT)};
\end{tikzpicture}
\end{adjustbox}
\caption{Portes contrôlées. Le contrôle $x$ est le premier chiffre du ket, donc le fil du bas.
(a) Le qubit $y$ subit $U$ si $x=1$ ; pour $U=X$ on retrouve le CNOT. (b) Si $x=1$, les qubits
$y$ et $z$ sont échangés. (c) Si $x=y=1$, le qubit $z$ est inversé.}
\label{fig:controlees}
\end{figure}

\subsubsection{Contrôlé-\texorpdfstring{$U$}{U}}

\begin{definition}{Porte contrôlée-$U$}
Soit $U$ une matrice unitaire $2\times2$. Sur le système $(x,y)$, de contrôle $x$ et de cible $y$,
\[
  \mathrm{C}U=\ket{0}\bra{0}_x\otimes I_y+\ket{1}\bra{1}_x\otimes U_y .
\]
Le premier terme agit quand le contrôle est $\ket{0}$ : la cible ne change pas ($I$). Le second
agit quand le contrôle est $\ket{1}$ : la cible subit $U$.
\end{definition}

Comme $\ket{0}\bra{0}\ket{x}=\delta_{x0}\ket{0}$ et $\ket{1}\bra{1}\ket{x}=\delta_{x1}\ket{1}$,
\[
  \mathrm{C}U\big(\ket{x}\otimes\ket{y}\big)=\ket{x}\otimes U^x\ket{y},\qquad U^0=I,\ U^1=U .
\]
Par la règle $A\otimes B=(A_{ij}B)$, $\ket{0}\bra{0}\otimes I=\begin{pmatrix}I&0\\0&0\end{pmatrix}$
et $\ket{1}\bra{1}\otimes U=\begin{pmatrix}0&0\\0&U\end{pmatrix}$ (blocs $2\times2$), donc
\[
  \mathrm{C}U=\begin{pmatrix}I&0\\0&U\end{pmatrix}
  =\begin{pmatrix}1&0&0&0\\0&1&0&0\\0&0&u_{00}&u_{01}\\0&0&u_{10}&u_{11}\end{pmatrix}
  \qquad\text{pour }U=\begin{pmatrix}u_{00}&u_{01}\\u_{10}&u_{11}\end{pmatrix}.
\]
Cette matrice est unitaire car $I^\dagger I=I$ et $U^\dagger U=I$.

\paragraph{Le CNOT est un cas particulier.} Pour $U=\sigma_x$,
$\mathrm{C}\sigma_x=\ket{0}\bra{0}_x\otimes I_y+\ket{1}\bra{1}_x\otimes\sigma_x$ est le CNOT de
contrôle $x$ (premier chiffre), c'est-à-dire $\mathrm{CNOT}_{12}$ de la section~\ref{sec:cnot}.
On calcule l'action sur les quatre états de base, avec $\braket{0}{0}=\braket{1}{1}=1$,
$\braket{0}{1}=\braket{1}{0}=0$, $\sigma_x\ket{0}=\ket{1}$ et $\sigma_x\ket{1}=\ket{0}$ :
\begin{align*}
  \mathrm{C}\sigma_x\ket{00}&=\big[\ket{0}\bra{0}_x\otimes I_y\big]\ket{0}_x\ket{0}_y
      +\big[\ket{1}\bra{1}_x\otimes\sigma_x\big]\ket{0}_x\ket{0}_y\\
    &=\ket{0}\braket{0}{0}\otimes I\ket{0}+\ket{1}\braket{1}{0}\otimes\sigma_x\ket{0}
     =\ket{0}\otimes\ket{0}+0=\ket{00},\\[1.2ex]
  \mathrm{C}\sigma_x\ket{01}&=\ket{0}\braket{0}{0}\otimes I\ket{1}+\ket{1}\braket{1}{0}\otimes\sigma_x\ket{1}
     =\ket{0}\otimes\ket{1}+0=\ket{01},\\[1.2ex]
  \mathrm{C}\sigma_x\ket{10}&=\ket{0}\braket{0}{1}\otimes I\ket{0}+\ket{1}\braket{1}{1}\otimes\sigma_x\ket{0}
     =0+\ket{1}\otimes\ket{1}=\ket{11},\\[1.2ex]
  \mathrm{C}\sigma_x\ket{11}&=\ket{0}\braket{0}{1}\otimes I\ket{1}+\ket{1}\braket{1}{1}\otimes\sigma_x\ket{1}
     =0+\ket{1}\otimes\ket{0}=\ket{10}.
\end{align*}
On retrouve la table du CNOT : la cible n'est inversée que si le contrôle vaut $1$.

\paragraph{Un autre exemple : $U=Z$.} Ici
\[
  \mathrm{C}Z=\begin{pmatrix}I&0\\0&Z\end{pmatrix}=\mathrm{diag}(1,1,1,-1) :
\]
la porte multiplie $\ket{11}$ par $-1$ et ne touche pas les autres états de base. Elle est
symétrique entre les deux qubits. Elle se déduit du CNOT par
$\mathrm{C}Z=(I\otimes H)\,\mathrm{C}X\,(I\otimes H)$, où $\mathrm{C}X$ est le contrôlé-$X$.
En effet, $I\otimes H=\mathrm{diag}(H,H)$ (deux blocs $2\times2$ égaux à $H$), et le produit de
matrices diagonales par blocs se fait bloc par bloc :
\[
  (I\otimes H)\begin{pmatrix}I&0\\0&X\end{pmatrix}(I\otimes H)
  =\begin{pmatrix}H\,I\,H&0\\0&H\,X\,H\end{pmatrix}
  =\begin{pmatrix}I&0\\0&HXH\end{pmatrix},
\]
car $H^2=I$, et
\[
  HXH=\frac12\begin{pmatrix}1&1\\1&-1\end{pmatrix}\begin{pmatrix}0&1\\1&0\end{pmatrix}\begin{pmatrix}1&1\\1&-1\end{pmatrix}
  =\frac12\begin{pmatrix}1&1\\1&-1\end{pmatrix}\begin{pmatrix}1&-1\\1&1\end{pmatrix}
  =\frac12\begin{pmatrix}2&0\\0&-2\end{pmatrix}=Z
\]
(section~\ref{sec:hadamard} : $H$ échange les axes $x$ et $z$).

\subsubsection{Contrôlé-SWAP : la porte de Fredkin}

\begin{definition}{Porte de Fredkin (C-SWAP)}
Sur trois qubits $(x,y,z)$ (espace de dimension $2^3=8$), avec $x$ pour contrôle,
\[
  \mathrm{C\text{-}SWAP}=\ket{0}\bra{0}_x\otimes I_y\otimes I_z+\ket{1}\bra{1}_x\otimes\mathrm{SWAP}_{yz}.
\]
Si $x=0$, rien ne change ; si $x=1$, les qubits $y$ et $z$ sont échangés.
\end{definition}

Calcul sur quatre états de base (avec $I\ket{ab}=\ket{ab}$ et $\mathrm{SWAP}\ket{ab}=\ket{ba}$) :
\begin{align*}
  \mathrm{C\text{-}SWAP}\ket{000}&=\ket{0}\braket{0}{0}\otimes I\ket{00}+\ket{1}\braket{1}{0}\otimes\mathrm{SWAP}\ket{00}
     =\ket{0}\otimes\ket{00}+0=\ket{000},\\[1.2ex]
  \mathrm{C\text{-}SWAP}\ket{101}&=\ket{0}\braket{0}{1}\otimes I\ket{01}+\ket{1}\braket{1}{1}\otimes\mathrm{SWAP}\ket{01}
     =0+\ket{1}\otimes\ket{10}=\ket{110},\\[1.2ex]
  \mathrm{C\text{-}SWAP}\ket{110}&=\ket{0}\braket{0}{1}\otimes I\ket{10}+\ket{1}\braket{1}{1}\otimes\mathrm{SWAP}\ket{10}
     =0+\ket{1}\otimes\ket{01}=\ket{101},\\[1.2ex]
  \mathrm{C\text{-}SWAP}\ket{010}&=\ket{0}\braket{0}{0}\otimes I\ket{10}+\ket{1}\braket{1}{0}\otimes\mathrm{SWAP}\ket{10}
     =\ket{0}\otimes\ket{10}+0=\ket{010}.
\end{align*}
Les huit états de base donnent :
\begin{center}
\renewcommand{\arraystretch}{1.3}
\begin{tabular}{@{}ccccccccc@{}}
\toprule
Entrée & $\ket{000}$ & $\ket{001}$ & $\ket{010}$ & $\ket{011}$ & $\ket{100}$ & $\ket{101}$ & $\ket{110}$ & $\ket{111}$ \\
\midrule
Sortie & $\ket{000}$ & $\ket{001}$ & $\ket{010}$ & $\ket{011}$ & $\ket{100}$ & $\ket{110}$ & $\ket{101}$ & $\ket{111}$ \\
\bottomrule
\end{tabular}
\end{center}
Seuls $\ket{101}$ et $\ket{110}$ changent (contrôle à $1$ et $y\neq z$) ; pour $\ket{100}$ et
$\ket{111}$ l'échange de deux qubits égaux ne change rien. Dans la base
$\ket{000},\ket{001},\dots,\ket{111}$ (ordre binaire), la matrice est
$\mathrm{diag}(I_4,\mathrm{SWAP})$ (deux blocs $4\times4$) :
\[
  \mathrm{C\text{-}SWAP}=\begingroup\setlength{\arraycolsep}{3pt}\begin{pmatrix}
    1&0&0&0&0&0&0&0\\
    0&1&0&0&0&0&0&0\\
    0&0&1&0&0&0&0&0\\
    0&0&0&1&0&0&0&0\\
    0&0&0&0&1&0&0&0\\
    0&0&0&0&0&0&1&0\\
    0&0&0&0&0&1&0&0\\
    0&0&0&0&0&0&0&1
  \end{pmatrix}\endgroup .
\]
C'est l'identité dont les lignes $6$ et $7$ sont échangées.

\paragraph{Contrôle en superposition.} Avec $\ket{+}_x\otimes\ket{0}_y\otimes\ket{1}_z
=\frac{1}{\sqrt2}\big(\ket{001}+\ket{101}\big)$ :
\[
  \mathrm{C\text{-}SWAP}\,\frac{1}{\sqrt2}\big(\ket{001}+\ket{101}\big)
  =\frac{1}{\sqrt2}\big(\ket{001}+\ket{110}\big).
\]
Le résultat n'est pas un produit : les coefficients de $\ket{x}\otimes\ket{yz}$ forment la
matrice (lignes $x=0,1$, colonnes $yz=00,01,10,11$)
\[
  \frac{1}{\sqrt2}\begin{pmatrix}0&1&0&0\\0&0&1&0\end{pmatrix},
\]
dont les deux lignes ne sont pas proportionnelles. Le qubit $x$ est intriqué avec la paire
$(y,z)$.

\subsubsection{Contrôlé-contrôlé-NOT : la porte de Toffoli}

\begin{definition}{Porte de Toffoli (CCNOT)}
Sur trois qubits $(x,y,z)$, avec $x$ et $y$ pour contrôles et $z$ pour cible,
\[
  \mathrm{CCNOT}=\ket{0}\bra{0}_x\otimes I_y\otimes I_z
   +\ket{1}\bra{1}_x\otimes\Big[\ket{0}\bra{0}_y\otimes I_z+\ket{1}\bra{1}_y\otimes\sigma_x\Big].
\]
Le crochet est un CNOT (contrôle $y$, cible $z$), appliqué seulement si $x=1$ : si $x=0$, rien ;
si $x=1$ et $y=0$, rien ; si $x=1$ et $y=1$, le qubit $z$ est inversé.
\end{definition}

Calcul de $\mathrm{CCNOT}\ket{110}$ :
\begin{align*}
  \mathrm{CCNOT}\ket{110}&=\ket{0}\braket{0}{1}\otimes I\ket{1}\otimes I\ket{0}
     +\ket{1}\braket{1}{1}\otimes\Big[\ket{0}\braket{0}{1}\otimes I\ket{0}
       +\ket{1}\braket{1}{1}\otimes\sigma_x\ket{0}\Big]\\
  &=0+\ket{1}\otimes\Big[0+\ket{1}\otimes\ket{1}\Big]=\ket{111}.
\end{align*}
Pour les huit états de base, $\mathrm{CCNOT}\ket{x,y,z}=\ket{x,\,y,\,z\oplus xy}$ :
\begin{center}
\renewcommand{\arraystretch}{1.3}
\begin{tabular}{@{}ccccccccc@{}}
\toprule
Entrée & $\ket{000}$ & $\ket{001}$ & $\ket{010}$ & $\ket{011}$ & $\ket{100}$ & $\ket{101}$ & $\ket{110}$ & $\ket{111}$ \\
\midrule
Sortie & $\ket{000}$ & $\ket{001}$ & $\ket{010}$ & $\ket{011}$ & $\ket{100}$ & $\ket{101}$ & $\ket{111}$ & $\ket{110}$ \\
\bottomrule
\end{tabular}
\end{center}
Seuls $\ket{110}$ et $\ket{111}$ sont échangés : la matrice est $\mathrm{diag}(I_6,X)$
(un bloc $6\times6$ égal à l'identité, puis $X$),
\[
  \mathrm{CCNOT}=\begingroup\setlength{\arraycolsep}{3pt}\begin{pmatrix}
    1&0&0&0&0&0&0&0\\
    0&1&0&0&0&0&0&0\\
    0&0&1&0&0&0&0&0\\
    0&0&0&1&0&0&0&0\\
    0&0&0&0&1&0&0&0\\
    0&0&0&0&0&1&0&0\\
    0&0&0&0&0&0&0&1\\
    0&0&0&0&0&0&1&0
  \end{pmatrix}\endgroup .
\]

\begin{remarque}[Portes réversibles]
Les portes de Fredkin et de Toffoli sont des permutations des états de base qui échangent
seulement deux d'entre eux : elles sont leurs propres inverses ($\mathrm{C\text{-}SWAP}^2=\mathrm{CCNOT}^2=\mathrm{Id}$)
et unitaires. Avec $z=0$, la Toffoli calcule le ET logique $x\wedge y$ sans rien effacer :
$\ket{x,y,0}\to\ket{x,y,xy}$ ; avec $z=1$, elle calcule le NON-ET. Elle sert à écrire tout calcul
classique de façon réversible.
\end{remarque}

\clearpage
% ---------------------------------------------------------------------
\subsection{Récapitulatif du chapitre}\label{sec:recap4}
% ---------------------------------------------------------------------

\begin{center}
\adjustbox{max width=\linewidth}{%
\begin{tikzpicture}[>=Stealth, font=\footnotesize,
    boite/.style={draw=insarouge, fill=insarouge!6, rounded corners=2pt, thick,
                  minimum width=4.4cm, minimum height=1.5cm, align=center, inner sep=3pt},
    lien/.style={->, very thick, insagris},
    etiq/.style={font=\scriptsize\itshape, text=insagris, inner sep=2pt}]
  \node[boite] (uni) at (0,0)
    {\textbf{Porte unitaire}\\ $U^\dagger U=\mathrm{Id}$, \ $\ket{\Psi}\to U\ket{\Psi}$\\ {\scriptsize\S\,\ref{sec:portes}}};
  \node[boite] (un) at (5.9,0)
    {\textbf{Portes à un qubit}\\ $X$, $Y$, $Z$, $H$, $S$, $T$\\ {\scriptsize\S\,\ref{sec:pauli}--\ref{sec:phase}}};
  \node[boite] (blo) at (11.8,0)
    {\textbf{Rotations de Bloch}\\ $X,Y,Z,H$ : angle $\pi$ ; $P_\theta$ autour de $z$\\
     {\scriptsize\S\,\ref{sec:portes-bloch}}};
  \node[boite] (cir) at (11.8,-2.5)
    {\textbf{Circuits}\\ fils, boîtes, $U_2U_1$\\ {\scriptsize\S\,\ref{sec:circuits}, \ref{sec:op2}}};
  \node[boite] (cn) at (5.9,-2.5)
    {\textbf{CNOT et SWAP}\\ $\ket{a,b}\to\ket{a\oplus b,b}$\\ {\scriptsize\S\,\ref{sec:cnot}, \ref{sec:swap}}};
  \node[boite] (cu) at (0,-2.5)
    {\textbf{Portes contrôlées}\\ $\mathrm{C}U=\mathrm{diag}(I,U)$, Fredkin, Toffoli\\
     {\scriptsize\S\,\ref{sec:controlees}}};
  \draw[lien] (uni) -- (un);
  \draw[lien] (un) -- (blo);
  \draw[lien] (blo) -- (cir);
  \draw[lien] (cir) -- (cn) node[etiq, midway, above=1pt] {2 qubits};
  \draw[lien] (cn) -- (cu) node[etiq, midway, above=1pt] {contrôle};
\end{tikzpicture}}
\end{center}

\begin{aretenir}
\begin{itemize}
  \item Une porte est un opérateur unitaire : elle conserve la norme, elle est réversible
        ($U^{-1}=U^\dagger$) et agit sur la sphère de Bloch comme une rotation. C'est l'évolution de
        Schrödinger, $U=\ee^{-\ii\Hop t/\hbar}$.
  \item $X$ échange $\ket{0}$ et $\ket{1}$, $Z$ change le signe de $\ket{1}$, $Y=\ii XZ$ fait les deux ;
        $H$ échange les bases $Z$ et $X$ et crée des superpositions ; $P_\theta$, $S=P_{\pi/2}$ et
        $T=P_{\pi/4}$ font tourner autour de $z$ ($T^2=S$, $S^2=Z$).
  \item Dans un circuit, le temps va de gauche à droite et la matrice est le produit dans l'ordre
        inverse, $U_2U_1$. Le fil du haut porte le dernier chiffre du ket. Des portes en parallèle
        forment un produit tensoriel $U_A\otimes U_B$.
  \item Le CNOT inverse la cible si le contrôle vaut $1$ : $\ket{a,b}\to\ket{a\oplus b,b}$. Le SWAP
        échange deux qubits ; $\mathrm{CNOT}_{21}=\mathrm{SWAP}\,\mathrm{CNOT}_{12}\,\mathrm{SWAP}$.
  \item Une porte contrôlée a pour matrice $\mathrm{diag}(I,U)$. Fredkin (contrôlé-SWAP) et
        Toffoli (CCNOT) sont des permutations des états de base : réversibles et universelles
        pour le calcul classique.
\end{itemize}
\end{aretenir}

\begin{center}
\renewcommand{\arraystretch}{1.45}
\begin{tabular}{@{}l l r@{}}
\toprule
\textbf{Notion} & \textbf{Formule} & \textbf{§} \\
\midrule
Porte unitaire & $U^\dagger U=\mathrm{Id}$, \quad $\bra{\Psi}U^\dagger U\ket{\Psi}=1$ & \ref{sec:portes} \\
Pauli & $X\ket{0}=\ket{1}$, \quad $Z\ket{1}=-\ket{1}$, \quad $X^2=Y^2=Z^2=\mathrm{Id}$ & \ref{sec:pauli} \\
Hadamard & $H\ket{0}=\ket{+}$, \quad $H\ket{1}=\ket{-}$, \quad $H^2=\mathrm{Id}$ & \ref{sec:hadamard} \\
Phase & $P_\theta=\mathrm{diag}(1,\ee^{\ii\theta})$, \quad $S=P_{\pi/2}$, \quad $T=P_{\pi/4}$ & \ref{sec:phase} \\
Circuit & $U_1$ puis $U_2$ : \quad $U_2U_1$ & \ref{sec:circuits} \\
Portes locales & $(U\otimes I)(\ket{u}\otimes\ket{v})=U\ket{u}\otimes\ket{v}$ & \ref{sec:op2} \\
CNOT & $\ket{a,b}\to\ket{a\oplus b,b}$, \quad $\mathrm{CNOT}=I\otimes\ket{0}\bra{0}+X\otimes\ket{1}\bra{1}$
  & \ref{sec:cnot} \\
SWAP & $\mathrm{SWAP}\ket{a,b}=\ket{b,a}$, \quad $\mathrm{SWAP}=\mathrm{CNOT}_{12}\mathrm{CNOT}_{21}\mathrm{CNOT}_{12}$
  & \ref{sec:swap} \\
Contrôlé-$U$ & $\mathrm{C}U=\ket{0}\bra{0}\otimes I+\ket{1}\bra{1}\otimes U=\mathrm{diag}(I,U)$ & \ref{sec:controlees} \\
Toffoli & $\mathrm{CCNOT}\ket{x,y,z}=\ket{x,y,z\oplus xy}$ & \ref{sec:controlees} \\
\bottomrule
\end{tabular}
\end{center}

\begin{savoirfaire}
\begin{itemize}
  \item Appliquer une porte à un état (par la matrice ou par les ket-bra) et identifier sa rotation
        de Bloch (exercices~\ref{ex:conj} et~\ref{ex:phases}).
  \item Lire un circuit, écrire sa matrice dans le bon ordre et suivre l'état à chaque étape
        (exercice~\ref{ex:kickback}).
  \item Écrire la matrice de $U\otimes I$, du CNOT, du SWAP ou d'une porte contrôlée, et calculer
        leur action sur les états de base.
  \item Reconnaître une porte réversible qui calcule une fonction booléenne
        (exercice~\ref{ex:toffoli}).
\end{itemize}
\end{savoirfaire}
.
%  Il ne contient ni préambule ni \begin{document} : les environnements
%  (definition, resultat, remarque, aretenir, savoirfaire, exercice, correction),
%  les macros (\ii, \ee, \dd, \pd, \pdd, \Hop, \ket, \bra, \braket), les symboles
%  de circuits (ctl, cible, croix, mesure, porte) et les couleurs (insarouge,
%  insagris) sont définis dans main.tex. Un chapitre = une \section.
% =====================================================================

% Macros de Dirac (déjà définies dans main.tex ; ce bloc évite l'erreur
% « Undefined control sequence \ket » si main.tex est une ancienne version).
\providecommand{\ket}[1]{|#1\rangle}
\providecommand{\bra}[1]{\langle #1|}
\providecommand{\braket}[2]{\langle #1|#2\rangle}

% =====================================================================
\section{Portes quantiques et systèmes multi-qubits}\label{chap:portes}
% =====================================================================

Un calcul quantique est une suite de transformations unitaires appliquées à des qubits. Ce
chapitre décrit les portes à un qubit (Pauli, Hadamard, phase) et leur effet sur la sphère de
Bloch, la lecture d'un circuit, puis les portes qui agissent sur plusieurs qubits : CNOT, SWAP et
portes contrôlées.

% ---------------------------------------------------------------------
\subsection{Portes quantiques}\label{sec:portes}
% ---------------------------------------------------------------------

Un calcul quantique transforme des états. Les opérations sur un seul qubit sont appelées
\emph{portes} ; les sections suivantes passent à plusieurs qubits.

\begin{resultat}{Postulat 4 : évolution d'un système isolé}
L'évolution d'un système isolé est décrite par un opérateur unitaire : $\ket{\Psi}\to U\ket{\Psi}$
avec $U^\dagger U=\mathrm{Id}$.
\end{resultat}

C'est la forme de l'équation de Schrödinger du chapitre~\ref{chap:histoire}. Pour un hamiltonien
$\Hop$ indépendant du temps, $\ii\hbar\,\partial_t\ket{\Psi}=\Hop\ket{\Psi}$ a pour solution
$\ket{\Psi(t)}=U(t)\ket{\Psi(0)}$ avec $U(t)=\ee^{-\ii\Hop t/\hbar}$. Comme $\Hop$ est hermitien,
$U(t)^\dagger=\ee^{+\ii\Hop t/\hbar}$ et $U^\dagger U=\mathrm{Id}$ : l'évolution est unitaire. Une
porte est une telle évolution, obtenue en contrôlant $\Hop$ pendant une durée donnée.

\begin{definition}{Porte quantique}
Une porte à un qubit est un opérateur $U$ (matrice $2\times2$) \emph{unitaire} :
$U^\dagger U=\mathrm{Id}$. Elle transforme $\ket{\Psi}\to U\ket{\Psi}$, conserve la norme
($\bra{\Psi}U^\dagger U\ket{\Psi}=\braket{\Psi}{\Psi}=1$) et est réversible
($U^{-1}=U^\dagger$). Sur la sphère de Bloch (section~\ref{sec:bloch}), c'est une
\emph{rotation}.
\end{definition}

% ---------------------------------------------------------------------
\subsection{Portes de Pauli}\label{sec:pauli}
% ---------------------------------------------------------------------

Les matrices de Pauli de la section~\ref{sec:op-dirac} sont trois portes.

\begin{definition}{Portes $X$, $Y$, $Z$}
\begin{center}
\renewcommand{\arraystretch}{1}
\begin{tabular}{@{}lccl@{}}
\textbf{Porte} & \textbf{Matrice} & \textbf{Ket-bra} & \textbf{Nom} \\[0.8ex]
$X=\sigma_x$ & $\begin{pmatrix}0 & 1\\ 1 & 0\end{pmatrix}$
  & $\ket{0}\bra{1}+\ket{1}\bra{0}$ & bit flip \\[2.2ex]
$Y=\sigma_y$ & $\begin{pmatrix}0 & -\ii\\ \ii & 0\end{pmatrix}$
  & $-\ii\ket{0}\bra{1}+\ii\ket{1}\bra{0}$ &  \\[2.2ex]
$Z=\sigma_z$ & $\begin{pmatrix}1 & 0\\ 0 & -1\end{pmatrix}$
  & $\ket{0}\bra{0}-\ket{1}\bra{1}$ & phase flip \\
\end{tabular}
\end{center}
\end{definition}

\paragraph{Porte $X$ (inverseur).} Avec $\braket{0}{0}=\braket{1}{1}=1$ et
$\braket{0}{1}=\braket{1}{0}=0$ :
\begin{align*}
  X\ket{0}&=\ket{0}\braket{1}{0}+\ket{1}\braket{0}{0}=0\cdot\ket{0}+1\cdot\ket{1}=\ket{1},\\
  X\ket{1}&=\ket{0}\braket{1}{1}+\ket{1}\braket{0}{1}=1\cdot\ket{0}+0\cdot\ket{1}=\ket{0}.
\end{align*}
$X$ échange $\ket{0}$ et $\ket{1}$ : c'est le NOT quantique (\emph{bit flip}). Sur la
sphère de Bloch, c'est une rotation de $\pi$ autour de l'axe $x$ : le pôle nord devient
le pôle sud. Les états de l'axe $x$ ne bougent pas : $X\ket{+}=\ket{+}$ et
$X\ket{-}=-\ket{-}$.

\paragraph{Porte $Z$ (changement de phase).}
\begin{align*}
  Z\ket{0}&=\ket{0}\braket{0}{0}-\ket{1}\braket{1}{0}=\ket{0},\\
  Z\ket{1}&=\ket{0}\braket{0}{1}-\ket{1}\braket{1}{1}=-\ket{1},
\end{align*}
donc $Z$ change le signe de la composante sur $\ket{1}$ (\emph{phase flip}). Sur les
états de la base $X$, par linéarité :
\begin{align*}
  Z\ket{+}&=\tfrac{1}{\sqrt2}\big(Z\ket{0}+Z\ket{1}\big)=\tfrac{1}{\sqrt2}\big(\ket{0}-\ket{1}\big)=\ket{-},\\
  Z\ket{-}&=\tfrac{1}{\sqrt2}\big(Z\ket{0}-Z\ket{1}\big)=\tfrac{1}{\sqrt2}\big(\ket{0}+\ket{1}\big)=\ket{+}.
\end{align*}
Sous forme matricielle, $Z\ket{+}=\begin{pmatrix}1 & 0\\ 0 & -1\end{pmatrix}
\frac{1}{\sqrt2}\begin{pmatrix}1\\ 1\end{pmatrix}=\frac{1}{\sqrt2}\begin{pmatrix}1\\ -1\end{pmatrix}=\ket{-}$.
$Z$ est une rotation de $\pi$ autour de l'axe $z$ : elle laisse $\ket{0}$ et $\ket{1}$ en
place et échange $\ket{+}$ et $\ket{-}$.

\paragraph{Porte $Y$.} Avec $Y=-\ii\ket{0}\bra{1}+\ii\ket{1}\bra{0}$ :
\begin{align*}
  Y\ket{0}&=-\ii\ket{0}\braket{1}{0}+\ii\ket{1}\braket{0}{0}=\ii\ket{1},\\
  Y\ket{1}&=-\ii\ket{0}\braket{1}{1}+\ii\ket{1}\braket{0}{1}=-\ii\ket{0}.
\end{align*}
Sur les états de la base $X$ :
\begin{align*}
  Y\ket{+}&=\tfrac{1}{\sqrt2}\big(\ii\ket{1}-\ii\ket{0}\big)=-\ii\,\tfrac{1}{\sqrt2}\big(\ket{0}-\ket{1}\big)=-\ii\ket{-},\\
  Y\ket{-}&=\tfrac{1}{\sqrt2}\big(\ii\ket{1}+\ii\ket{0}\big)=\ii\,\tfrac{1}{\sqrt2}\big(\ket{0}+\ket{1}\big)=\ii\ket{+}.
\end{align*}
Un facteur global ($-\ii$ ou $\ii$) ne change pas l'état physique : $Y$ échange bien
$\ket{+}$ et $\ket{-}$. Sur les états de l'axe $y$ :
\begin{align*}
  Y\ket{{+}\ii}&=\tfrac{1}{\sqrt2}\big(Y\ket{0}+\ii Y\ket{1}\big)
   =\tfrac{1}{\sqrt2}\big(\ii\ket{1}+\ket{0}\big)=\ket{{+}\ii},\\
  Y\ket{{-}\ii}&=\tfrac{1}{\sqrt2}\big(Y\ket{0}-\ii Y\ket{1}\big)
   =\tfrac{1}{\sqrt2}\big(\ii\ket{1}-\ket{0}\big)=-\ket{{-}\ii}.
\end{align*}
$Y$ est une rotation de $\pi$ autour de l'axe $y$.

\begin{remarque}[Chaque Pauli fixe son axe]
$Z\ket{0}=\ket{0}$, $Z\ket{1}=-\ket{1}$ ; $X\ket{\pm}=\pm\ket{\pm}$ ; $Y\ket{{\pm}\ii}=\pm\ket{{\pm}\ii}$.
Les bases $Z$, $X$, $Y$ de la section~\ref{sec:bloch} sont les bases propres de
$Z$, $X$, $Y$. Une porte de Pauli laisse donc son propre axe invariant et retourne les
deux autres ; c'est pourquoi les exemples utilisent les états des deux autres axes.
\end{remarque}

% ---------------------------------------------------------------------
\subsection{Porte de Hadamard}\label{sec:hadamard}
% ---------------------------------------------------------------------

\begin{definition}{Porte de Hadamard}
\[
  H=\frac{1}{\sqrt2}\begin{pmatrix}1 & 1\\ 1 & -1\end{pmatrix}
  =\ket{+}\bra{0}+\ket{-}\bra{1}.
\]
\end{definition}

\begin{align*}
  H\ket{0}&=\frac{1}{\sqrt2}\begin{pmatrix}1 & 1\\ 1 & -1\end{pmatrix}\begin{pmatrix}1\\ 0\end{pmatrix}
    =\frac{1}{\sqrt2}\begin{pmatrix}1\\ 1\end{pmatrix}=\ket{+},\\[0.6ex]
  H\ket{1}&=\frac{1}{\sqrt2}\begin{pmatrix}1 & 1\\ 1 & -1\end{pmatrix}\begin{pmatrix}0\\ 1\end{pmatrix}
    =\frac{1}{\sqrt2}\begin{pmatrix}1\\ -1\end{pmatrix}=\ket{-},\\[0.6ex]
  H\ket{+}&=\frac{1}{2}\begin{pmatrix}1 & 1\\ 1 & -1\end{pmatrix}\begin{pmatrix}1\\ 1\end{pmatrix}
    =\frac{1}{2}\begin{pmatrix}2\\ 0\end{pmatrix}=\ket{0},\qquad
  H\ket{-}=\frac{1}{2}\begin{pmatrix}0\\ 2\end{pmatrix}=\ket{1}.
\end{align*}
Comme $H\ket{+}=\ket{0}$ et $H\ket{-}=\ket{1}$, on a $H^2=\mathrm{Id}$ ($H$ est son propre
inverse). Sur la sphère de Bloch, $H$ est une rotation de $\pi$ autour de l'axe
$(x+z)/\sqrt2$ : elle échange l'axe $x$ et l'axe $z$.

\begin{remarque}[À quoi sert $H$ ?]
\begin{itemize}[nosep]
  \item Elle \emph{crée une superposition} à partir d'un état de base :
        $H\ket{0}=\ket{+}$ donne $P(0)=P(1)=\frac12$. C'est la première étape de
        la plupart des algorithmes quantiques.
  \item Elle \emph{change de base} : elle envoie la base $Z$ sur la base $X$ et
        inversement. Mesurer dans la base $X$ revient à appliquer $H$ puis à mesurer
        dans la base $Z$ (figure~\ref{fig:mesurex}).
  \item Elle est réversible ($H^2=\mathrm{Id}$).
\end{itemize}
\end{remarque}

% ---------------------------------------------------------------------
\subsection{Portes de phase}\label{sec:phase}
% ---------------------------------------------------------------------

\begin{definition}{Portes de phase $P_\theta$, $S$ et $T$}
\[
  P_\theta=\begin{pmatrix}1 & 0\\ 0 & \ee^{\ii\theta}\end{pmatrix}
  =\ket{0}\bra{0}+\ee^{\ii\theta}\ket{1}\bra{1},
  \qquad
  S=P_{\pi/2}=\begin{pmatrix}1 & 0\\ 0 & \ii\end{pmatrix},
  \qquad
  T=P_{\pi/4}=\begin{pmatrix}1 & 0\\ 0 & \dfrac{1+\ii}{\sqrt2}\end{pmatrix}.
\]
\end{definition}

Ici $\ee^{\ii\pi/2}=\ii$ et $\ee^{\ii\pi/4}=\cos\frac\pi4+\ii\sin\frac\pi4=\frac{1+\ii}{\sqrt2}$.
$P_\theta$ laisse $\ket{0}$ inchangé et multiplie $\ket{1}$ par $\ee^{\ii\theta}$ :
\[
  P_\theta\big(\alpha\ket{0}+\beta\ket{1}\big)=\alpha\ket{0}+\ee^{\ii\theta}\beta\ket{1}.
\]
Dans la paramétrisation de Bloch, $\varphi\to\varphi+\theta$ : $P_\theta$ est une rotation
d'angle $\theta$ autour de l'axe $z$. On retrouve $P_\pi=Z$, et $S^2=Z$, $T^2=S$.

\paragraph{Exemples.}
\begin{align*}
  S\ket{+}&=\tfrac{1}{\sqrt2}\big(\ket{0}+\ii\ket{1}\big)=\ket{{+}\ii},\\
  S\ket{{+}\ii}&=\tfrac{1}{\sqrt2}\big(\ket{0}+\ii\cdot\ii\ket{1}\big)=\tfrac{1}{\sqrt2}\big(\ket{0}-\ket{1}\big)=\ket{-},\\
  T\ket{+}&=\tfrac{1}{\sqrt2}\begin{pmatrix}1 & 0\\ 0 & \frac{1+\ii}{\sqrt2}\end{pmatrix}\begin{pmatrix}1\\ 1\end{pmatrix}
    =\tfrac{1}{\sqrt2}\begin{pmatrix}1\\ \frac{1+\ii}{\sqrt2}\end{pmatrix}
    =\tfrac{1}{\sqrt2}\big(\ket{0}+\ee^{\ii\pi/4}\ket{1}\big)\quad\big(\text{angle polaire }\tfrac\pi2,\ \text{azimut }\tfrac\pi4\big).
\end{align*}
$S$ tourne donc $\ket{+}$ d'un quart de tour autour de $z$ (de l'axe $x$ à l'axe $y$),
puis d'un second quart de tour jusqu'à $\ket{-}$ ; $T$ fait un huitième de tour.

\begin{remarque}[Une porte de phase est une évolution libre]
Pour $\Hop=\frac{\hbar\omega}{2}Z$, $U(t)=\ee^{-\ii\Hop t/\hbar}=\mathrm{diag}\big(\ee^{-\ii\omega t/2},\ee^{\ii\omega t/2}\big)
=\ee^{-\ii\omega t/2}\,P_{\omega t}$. Au facteur global près, laisser évoluer le qubit pendant la
durée $t$ applique la porte de phase $P_{\omega t}$ : une rotation d'angle $\omega t$ autour de
l'axe $z$.
\end{remarque}

% ---------------------------------------------------------------------
\subsection{Portes et sphère de Bloch}\label{sec:portes-bloch}
% ---------------------------------------------------------------------

Toute porte à un qubit est une rotation de la sphère de Bloch. La figure~\ref{fig:portes} montre
les axes des portes usuelles et le tableau qui suit les résume.

\begin{figure}[!ht]
\centering
\begin{tikzpicture}[x={(-1.63cm,-0.56cm)}, y={(1.63cm,-0.56cm)}, z={(0cm,2.16cm)}, >=Stealth, scale=0.9]
  % sphère et grands cercles
  \shade[ball color=insagris!15, opacity=0.55] (0,0,0) circle[radius=2.3cm];
  \begin{scope}[insagris!75, thin]
    \draw[domain=-45:135, samples=60, variable=\t] plot ({cos(\t)},{sin(\t)},0);
    \draw[domain=135:315, samples=60, variable=\t, dashed] plot ({cos(\t)},{sin(\t)},0);
    \draw[domain=-62.7:117.3, samples=60, variable=\t] plot ({cos(\t)},0,{sin(\t)});
    \draw[domain=117.3:297.3, samples=60, variable=\t, dashed] plot ({cos(\t)},0,{sin(\t)});
    \draw[domain=-62.7:117.3, samples=60, variable=\t] plot (0,{cos(\t)},{sin(\t)});
    \draw[domain=117.3:297.3, samples=60, variable=\t, dashed] plot (0,{cos(\t)},{sin(\t)});
  \end{scope}
  % axes x, y, z (portes X, Y, Z)
  \draw[insagris, dashed] (0,0,0) -- (-1,0,0);
  \draw[insagris, dashed] (0,0,0) -- (0,-1,0);
  \draw[insagris, dashed] (0,0,0) -- (0,0,-1);
  \draw[insagris] (-1,0,0) -- (-1.3,0,0);
  \draw[insagris] (0,-1,0) -- (0,-1.3,0);
  \draw[insagris] (0,0,-1) -- (0,0,-1.3);
  \draw[insagris, -{Stealth[length=2.2mm]}] (0,0,0) -- (1.3,0,0);
  \draw[insagris, -{Stealth[length=2.2mm]}] (0,0,0) -- (0,1.3,0);
  \draw[insagris, -{Stealth[length=2.2mm]}] (0,0,0) -- (0,0,1.3);
  \node[anchor=east]  at (1.33,0,0) {$x$};
  \node[anchor=west]  at (0,1.33,0) {$y$};
  \node[anchor=south] at (0,0,1.33) {$z$};
  % axe de la porte H : (x+z)/sqrt(2)
  \draw[insarouge, dashed] (0,0,0) -- (-0.707,0,-0.707);
  \draw[insarouge, very thick, -{Stealth[length=2.6mm]}] (0,0,0) -- (0.98,0,0.98);
  \node[insarouge, anchor=south east] at (0.98,0,0.98) {axe de $H$};
  % rotation autour de z (portes de phase)
  \draw[insarouge, thick, -{Stealth[length=2.2mm]}, domain=15:105, samples=30, variable=\t]
    plot ({1.15*cos(\t)},{1.15*sin(\t)},0);
  \node[insarouge, anchor=north west, font=\small] at ({1.15*cos(60)},{1.15*sin(60)},0) {$P_\theta$};
\end{tikzpicture}
\caption{Portes et rotations sur la sphère de Bloch. Les portes $X$, $Y$, $Z$ sont des
rotations de $\pi$ autour des axes $x$, $y$, $z$ ; $H$ est une rotation de $\pi$ autour de
l'axe $(x+z)/\sqrt2$ (rouge) ; $P_\theta$ (donc $Z$, $S$, $T$) tourne d'un angle $\theta$
autour de $z$.}
\label{fig:portes}
\end{figure}

\begin{center}
\renewcommand{\arraystretch}{1}
\begin{tabular}{@{}lccc@{}}
\toprule
\textbf{Porte} & \textbf{Matrice} & \textbf{Effet sur la base} & \textbf{Rotation de Bloch} \\
\midrule
$X$ & $\begin{pmatrix}0 & 1\\ 1 & 0\end{pmatrix}$ & $\ket{0}\leftrightarrow\ket{1}$ & $\pi$ autour de $x$ \\[2.2ex]
$Y$ & $\begin{pmatrix}0 & -\ii\\ \ii & 0\end{pmatrix}$ & $\ket{0}\to\ii\ket{1},\ \ket{1}\to-\ii\ket{0}$ & $\pi$ autour de $y$ \\[2.2ex]
$Z$ & $\begin{pmatrix}1 & 0\\ 0 & -1\end{pmatrix}$ & $\ket{1}\to-\ket{1},\ \ket{\pm}\leftrightarrow\ket{\mp}$ & $\pi$ autour de $z$ \\[2.2ex]
$H$ & $\dfrac{1}{\sqrt2}\begin{pmatrix}1 & 1\\ 1 & -1\end{pmatrix}$ & $\ket{0}\leftrightarrow\ket{+},\ \ket{1}\leftrightarrow\ket{-}$ & $\pi$ autour de $\frac{x+z}{\sqrt2}$ \\[2.2ex]
$S$ & $\begin{pmatrix}1 & 0\\ 0 & \ii\end{pmatrix}$ & $\ket{1}\to\ii\ket{1}$ & $\frac\pi2$ autour de $z$ \\[2.2ex]
$T$ & $\begin{pmatrix}1 & 0\\ 0 & \frac{1+\ii}{\sqrt2}\end{pmatrix}$ & $\ket{1}\to\ee^{\ii\pi/4}\ket{1}$ & $\frac\pi4$ autour de $z$ \\
\bottomrule
\end{tabular}
\end{center}

\begin{figure}[!ht]
\centering
\begin{tikzpicture}[x={(-1.63cm,-0.56cm)}, y={(1.63cm,-0.56cm)}, z={(0cm,2.16cm)}, >=Stealth, scale=0.9]
  % sphère et grands cercles
  \shade[ball color=insagris!15, opacity=0.55] (0,0,0) circle[radius=2.3cm];
  \begin{scope}[insagris!75, thin]
    \draw[domain=-45:135, samples=60, variable=\t] plot ({cos(\t)},{sin(\t)},0);
    \draw[domain=135:315, samples=60, variable=\t, dashed] plot ({cos(\t)},{sin(\t)},0);
    \draw[domain=-62.7:117.3, samples=60, variable=\t] plot ({cos(\t)},0,{sin(\t)});
    \draw[domain=117.3:297.3, samples=60, variable=\t, dashed] plot ({cos(\t)},0,{sin(\t)});
    \draw[domain=-62.7:117.3, samples=60, variable=\t] plot (0,{cos(\t)},{sin(\t)});
    \draw[domain=117.3:297.3, samples=60, variable=\t, dashed] plot (0,{cos(\t)},{sin(\t)});
  \end{scope}
  % axes x, y, z (portes X, Y, Z)
  \draw[insagris, dashed] (0,0,0) -- (-1,0,0);
  \draw[insagris, dashed] (0,0,0) -- (0,-1,0);
  \draw[insagris, dashed] (0,0,0) -- (0,0,-1);
  \draw[insagris] (-1,0,0) -- (-1.3,0,0);
  \draw[insagris] (0,-1,0) -- (0,-1.3,0);
  \draw[insagris] (0,0,-1) -- (0,0,-1.3);
  \draw[insagris, -{Stealth[length=2.2mm]}] (0,0,0) -- (1.3,0,0);
  \draw[insagris, -{Stealth[length=2.2mm]}] (0,0,0) -- (0,1.3,0);
  \draw[insagris, -{Stealth[length=2.2mm]}] (0,0,0) -- (0,0,1.3);
  % trajet : |0> --H--> |+> --S--> |+i>
  \draw[insarouge, very thick, -{Stealth[length=2.6mm]}, domain=3:87, samples=30, variable=\t]
    plot ({sin(\t)},0,{cos(\t)});
  \draw[insarouge, very thick, -{Stealth[length=2.6mm]}, domain=3:87, samples=30, variable=\t]
    plot ({cos(\t)},{sin(\t)},0);
  \foreach \ax/\ay/\az in {0/0/1, 1/0/0, 0/1/0} \fill[insarouge] (\ax,\ay,\az) circle[radius=2pt];
  \node[insarouge, anchor=south] at (0,0,1.33) {$\ket{0}$};
  \node[insarouge, anchor=east] at (1.33,0,0) {$\ket{+}$};
  \node[insarouge, anchor=west] at (0,1.33,0) {$\ket{{+}\ii}$};
  \node[insarouge, font=\small, anchor=south west] at ({sin(45)},0,{cos(45)}) {$H$};
  \node[insarouge, font=\small, anchor=north west] at ({cos(45)},{sin(45)},0) {$S$};
\end{tikzpicture}
\caption{Trajet d'un état sous le circuit $H$ puis $S$ (flèches schématiques). $H$ envoie
$\ket{0}$ sur $\ket{+}$ ; $S$, rotation d'un quart de tour autour de $z$, envoie $\ket{+}$ sur
$\ket{{+}\ii}$. Les trois états sont alignés sur les axes $z$, $x$ et $y$.}
\label{fig:trajet}
\end{figure}

% ---------------------------------------------------------------------
\subsection{Circuits quantiques}\label{sec:circuits}
% ---------------------------------------------------------------------

Un \emph{circuit} représente un calcul : chaque qubit est un fil horizontal, le temps va de gauche
à droite, et chaque porte est un symbole posé sur les fils qu'elle concerne. La
figure~\ref{fig:symboles} donne les symboles utilisés dans ce document.

\begin{figure}[!ht]
\centering
\begin{adjustbox}{max width=\linewidth}
\begin{tikzpicture}[>=Stealth, font=\small,
    nom/.style={font=\footnotesize, align=center, anchor=north}]
  % (a) porte à un qubit
  \draw (0,0) -- (1.8,0);
  \node[porte] at (0.9,0) {$U$};
  \node[nom] at (0.9,-1.15) {porte $U$\\ à un qubit};
  % (b) CNOT
  \begin{scope}[xshift=3.2cm]
    \draw (0,0.6) -- (1.8,0.6);
    \draw (0,-0.6) -- (1.8,-0.6);
    \draw[insarouge, thick] (0.9,0.6) -- (0.9,-0.6);
    \pic at (0.9,0.6) {ctl};
    \pic at (0.9,-0.6) {cible};
    \node[nom] at (0.9,-1.15) {CNOT\\ (contrôle $\bullet$, cible $\oplus$)};
  \end{scope}
  % (c) SWAP
  \begin{scope}[xshift=6.4cm]
    \draw (0,0.6) -- (1.8,0.6);
    \draw (0,-0.6) -- (1.8,-0.6);
    \draw[insarouge, thick] (0.9,0.6) -- (0.9,-0.6);
    \pic at (0.9,0.6) {croix};
    \pic at (0.9,-0.6) {croix};
    \node[nom] at (0.9,-1.15) {SWAP\\ (échange)};
  \end{scope}
  % (d) contrôlé-U
  \begin{scope}[xshift=9.6cm]
    \draw (0,0.6) -- (1.8,0.6);
    \draw (0,-0.6) -- (1.8,-0.6);
    \draw[insarouge, thick] (0.9,-0.6) -- (0.9,0.3);
    \pic at (0.9,-0.6) {ctl};
    \node[porte] at (0.9,0.6) {$U$};
    \node[nom] at (0.9,-1.15) {contrôlé-$U$\\ (contrôle en bas)};
  \end{scope}
  % (e) mesure
  \begin{scope}[xshift=12.8cm]
    \draw (0,0) -- (0.6,0);
    \pic at (0.9,0) {mesure};
    \draw (1.2,0.03) -- (2.0,0.03);
    \draw (1.2,-0.03) -- (2.0,-0.03);
    \node[nom] at (1.0,-1.15) {mesure\\ (fil double : bit)};
  \end{scope}
\end{tikzpicture}
\end{adjustbox}
\caption{Symboles des circuits quantiques. Le point noir $\bullet$ marque un qubit de contrôle,
le symbole $\oplus$ une cible qui subit $X$, les deux croix un échange. Après une mesure, le fil
double transporte un bit classique.}
\label{fig:symboles}
\end{figure}

\begin{remarque}[Conventions de lecture]
\begin{itemize}[nosep]
  \item Le temps va de gauche à droite. Deux portes successives $U_1$ puis $U_2$ se multiplient
        dans l'ordre inverse : la matrice du circuit est $U_2U_1$, car $U_1$ agit d'abord sur le ket.
  \item Avec plusieurs fils, le fil du \emph{haut} porte le \emph{dernier} chiffre du ket et le
        fil du \emph{bas} le premier : si $A$ est le premier qubit et $B$ le second,
        $\ket{a,b}=\ket{a}_A\otimes\ket{b}_B$ et $B$ est dessiné au-dessus de $A$.
  \item Deux portes placées dans la même colonne, sur deux fils, agissent en parallèle : c'est un
        produit tensoriel $U_A\otimes U_B$ (section~\ref{sec:op2}).
\end{itemize}
\end{remarque}

\paragraph{Exemple.} Le circuit de la figure~\ref{fig:circuit-hs} applique $H$ puis $S$ à
$\ket{0}$. Sa matrice est $SH$ :
\[
  SH=\begin{pmatrix}1&0\\0&\ii\end{pmatrix}\frac{1}{\sqrt2}\begin{pmatrix}1&1\\1&-1\end{pmatrix}
  =\frac{1}{\sqrt2}\begin{pmatrix}1&1\\ \ii&-\ii\end{pmatrix},
  \qquad
  SH\ket{0}=\frac{1}{\sqrt2}\begin{pmatrix}1\\ \ii\end{pmatrix}=\ket{{+}\ii}.
\]
Le circuit prépare donc l'état $\ket{{+}\ii}$ de la base $Y$ à partir de $\ket{0}$
(figure~\ref{fig:trajet}). L'ordre compte : $HS\ket{0}=H\ket{0}=\ket{+}$, car $S\ket{0}=\ket{0}$.

\begin{figure}[!ht]
\centering
\begin{tikzpicture}[>=Stealth, font=\small]
  \draw (0,0) -- (6.6,0);
  \node[anchor=east] at (0,0) {$\ket{0}$};
  \node[porte] at (2.0,0) {$H$};
  \node[porte] at (4.3,0) {$S$};
  \node[anchor=west] at (6.6,0) {$\ket{{+}\ii}$};
  \foreach \x in {0.7, 3.15, 5.5} \draw[dashed, insagris!70] (\x,0.6) -- (\x,-0.95);
  \node[font=\footnotesize, anchor=north] at (0.7,-0.95) {$\ket{0}$};
  \node[font=\footnotesize, anchor=north] at (3.15,-0.95) {$\ket{+}$};
  \node[font=\footnotesize, anchor=north] at (5.5,-0.95) {$\ket{{+}\ii}$};
\end{tikzpicture}
\caption{Circuit $H$ puis $S$ appliqué à $\ket{0}$. Les états intermédiaires sont indiqués sous
les traits pointillés ; la matrice du circuit est $SH$.}
\label{fig:circuit-hs}
\end{figure}

\paragraph{Mesure.} Le symbole de mesure (compteur) donne le résultat $0$ ou $1$ dans la base $Z$.
Mesurer dans la base $X$ revient à appliquer $H$ puis à mesurer dans la base $Z$
(figure~\ref{fig:mesurex}) : comme $H\ket{\pm}=\ket{0}$ ou $\ket{1}$, le résultat $0$ correspond à
$\ket{+}$ et le résultat $1$ à $\ket{-}$, avec $P(\pm)=|\braket{\pm}{\Psi}|^2$.

\begin{figure}[!ht]
\centering
\begin{tikzpicture}[>=Stealth, font=\small]
  % (a) base Z
  \draw (0,0) -- (1.2,0);
  \node[anchor=east] at (0,0) {$\ket{\Psi}$};
  \pic at (1.5,0) {mesure};
  \draw (1.8,0.03) -- (2.6,0.03);
  \draw (1.8,-0.03) -- (2.6,-0.03);
  \node[anchor=west] at (2.7,0) {$0$ ou $1$};
  \node[font=\footnotesize] at (1.5,-0.95) {(a) mesure en base $Z$};
  % (b) base X
  \begin{scope}[xshift=7.5cm]
    \draw (0,0) -- (2.6,0);
    \node[anchor=east] at (0,0) {$\ket{\Psi}$};
    \node[porte] at (1.0,0) {$H$};
    \pic at (2.3,0) {mesure};
    \draw (2.6,0.03) -- (3.4,0.03);
    \draw (2.6,-0.03) -- (3.4,-0.03);
    \node[anchor=west] at (3.5,0) {$+$ ou $-$};
    \node[font=\footnotesize] at (1.9,-0.95) {(b) mesure en base $X$};
  \end{scope}
\end{tikzpicture}
\caption{(a) Mesure dans la base $Z$. (b) Mesure dans la base $X$ : porte de Hadamard, puis
mesure dans la base $Z$ ; les résultats $0$ et $1$ se lisent $+$ et $-$.}
\label{fig:mesurex}
\end{figure}

% ---------------------------------------------------------------------
\subsection{Portes locales sur deux qubits}\label{sec:op2}
% ---------------------------------------------------------------------

Une porte appliquée à un seul qubit d'un système de deux s'écrit avec un produit tensoriel
(section~\ref{sec:tenseur}) : $U\otimes I$ agit sur $A$ (premier facteur, fil du bas), $I\otimes U$
agit sur $B$, et $U_A\otimes U_B$ agit sur les deux en parallèle (figure~\ref{fig:local}).

\begin{figure}[!ht]
\centering
\begin{tikzpicture}[>=Stealth, font=\small]
  \foreach \xs/\haut/\bas/\tit in {0/{}/{$X$}/{(a) $X\otimes I$},
                                   5.4/{$X$}/{}/{(b) $I\otimes X$},
                                   10.8/{$H$}/{$H$}/{(c) $H\otimes H$}} {
    \begin{scope}[xshift=\xs cm]
      \draw (0,0.75) -- (2.6,0.75);
      \draw (0,0) -- (2.6,0);
      \node[anchor=east] at (0,0.75) {$B$};
      \node[anchor=east] at (0,0) {$A$};
      \ifx\haut\empty\else \node[porte] at (1.3,0.75) {\haut}; \fi
      \ifx\bas\empty\else \node[porte] at (1.3,0) {\bas}; \fi
      \node[font=\footnotesize] at (1.3,-0.85) {\tit};
    \end{scope}
  }
\end{tikzpicture}
\caption{Portes en parallèle sur deux qubits. Le qubit $A$, premier facteur du produit tensoriel,
est sur le fil du bas. (a) $X$ sur $A$ seulement. (b) $X$ sur $B$ seulement. (c) $H$ sur les deux.}
\label{fig:local}
\end{figure}

\subsubsection{Portes agissant sur un seul des deux qubits}

On note $I=\mathrm{Id}$ ($2\times2$). $X\otimes I$ agit sur $A$, $I\otimes X$ agit sur $B$ :
\begin{align*}
  I\otimes X&=\begin{pmatrix}1\cdot X & 0\cdot X\\ 0\cdot X & 1\cdot X\end{pmatrix}
    =\begin{pmatrix}0 & 1 & 0 & 0\\ 1 & 0 & 0 & 0\\ 0 & 0 & 0 & 1\\ 0 & 0 & 1 & 0\end{pmatrix},\\[1ex]
  X\otimes I&=\begin{pmatrix}0\cdot I & 1\cdot I\\ 1\cdot I & 0\cdot I\end{pmatrix}
    =\begin{pmatrix}0 & 0 & 1 & 0\\ 0 & 0 & 0 & 1\\ 1 & 0 & 0 & 0\\ 0 & 1 & 0 & 0\end{pmatrix}.
\end{align*}
Sur $\ket{00}=(1,0,0,0)^{\mathsf T}$ (première colonne de la matrice) :
$(I\otimes X)\ket{00}=(0,1,0,0)^{\mathsf T}=\ket{01}$ et
$(X\otimes I)\ket{00}=(0,0,1,0)^{\mathsf T}=\ket{10}$. Avec la règle :
$(I\otimes X)\ket{00}=I\ket{0}\otimes X\ket{0}=\ket{0}\otimes\ket{1}=\ket{01}$.

Pour la porte de Hadamard, avec $H=\frac{1}{\sqrt2}\begin{pmatrix}1 & 1\\ 1 & -1\end{pmatrix}$ :
\[
  H\otimes I=\frac{1}{\sqrt2}\begin{pmatrix}1\cdot I & 1\cdot I\\ 1\cdot I & -1\cdot I\end{pmatrix}
  =\frac{1}{\sqrt2}\begin{pmatrix}1 & 0 & 1 & 0\\ 0 & 1 & 0 & 1\\ 1 & 0 & -1 & 0\\ 0 & 1 & 0 & -1\end{pmatrix}.
\]
Ses colonnes donnent l'image des quatre états de base :
\begin{align*}
  (H\otimes I)\ket{00}&=\tfrac{1}{\sqrt2}(1,0,1,0)^{\mathsf T}=\ket{+}\otimes\ket{0},&
  (H\otimes I)\ket{01}&=\tfrac{1}{\sqrt2}(0,1,0,1)^{\mathsf T}=\ket{+}\otimes\ket{1},\\
  (H\otimes I)\ket{10}&=\tfrac{1}{\sqrt2}(1,0,-1,0)^{\mathsf T}=\ket{-}\otimes\ket{0},&
  (H\otimes I)\ket{11}&=\tfrac{1}{\sqrt2}(0,1,0,-1)^{\mathsf T}=\ket{-}\otimes\ket{1}.
\end{align*}

\subsubsection{Hadamard sur les deux qubits}

\[
  H\otimes H=\frac12\begin{pmatrix}1\cdot H' & 1\cdot H'\\ 1\cdot H' & -1\cdot H'\end{pmatrix},
  \quad H'=\begin{pmatrix}1 & 1\\ 1 & -1\end{pmatrix},
  \qquad
  H\otimes H=\frac12\begin{pmatrix}1 & 1 & 1 & 1\\ 1 & -1 & 1 & -1\\ 1 & 1 & -1 & -1\\ 1 & -1 & -1 & 1\end{pmatrix}.
\]
La première colonne donne $(H\otimes H)\ket{00}=\frac12(1,1,1,1)^{\mathsf T}=\ket{+}\otimes\ket{+}$
(cohérent avec $H\ket{0}\otimes H\ket{0}$) : les quatre résultats $00,01,10,11$ sont
équiprobables, chacun avec la probabilité $\frac14$. Sur $n$ qubits, $H^{\otimes n}\ket{0\ldots0}$
donne une superposition uniforme de $2^n$ états.

% ---------------------------------------------------------------------
\subsection{Porte CNOT}\label{sec:cnot}
% ---------------------------------------------------------------------

\begin{definition}{Porte CNOT (NOT contrôlé)}
La porte CNOT agit sur deux qubits : le qubit $B$ est le \emph{contrôle}, le qubit $A$ est la
\emph{cible}. La cible est inversée (porte $X$) si le contrôle vaut $1$ et reste inchangée
s'il vaut $0$. Avec $\ket{a,b}=\ket{a}_A\otimes\ket{b}_B$ :
\[
  \mathrm{CNOT}\,\ket{a,b}=\ket{a\oplus b,\,b},
  \qquad
  \mathrm{CNOT}=I\otimes\ket{0}\bra{0}+X\otimes\ket{1}\bra{1},
\]
où $\oplus$ est l'addition modulo $2$ ($0\oplus0=1\oplus1=0$ et $0\oplus1=1\oplus0=1$).
\end{definition}

Sur le circuit de la figure~\ref{fig:cnot}, le point noir posé sur le fil de $B$ (en haut) marque
le contrôle et le symbole $\oplus$ posé sur le fil de $A$ (en bas) marque la cible.

\begin{figure}[!ht]
\centering
\begin{tikzpicture}[>=Stealth, font=\small]
  \draw (0,1.2) -- (6,1.2);
  \draw (0,0) -- (6,0);
  \node[anchor=east] at (0,1.2) {$\ket{b}_B$};
  \node[anchor=east] at (0,0) {$\ket{a}_A$};
  \node[anchor=west] at (6,1.2) {$\ket{b}_B$};
  \node[anchor=west] at (6,0) {$\ket{a\oplus b}_A$};
  \draw[insarouge, thick] (3,1.2) -- (3,0);
  \pic at (3,1.2) {ctl};
  \pic at (3,0) {cible};
  \node[anchor=south, text=insagris] at (3,1.5) {qubit de contrôle};
  \node[anchor=north, text=insagris] at (3,-0.5) {qubit cible};
\end{tikzpicture}
\caption{Porte CNOT : le fil du haut ($B$) est le contrôle, le fil du bas ($A$) est la cible. Le
ket s'écrit $\ket{a,b}$ : le fil du haut donne le dernier chiffre.}
\label{fig:cnot}
\end{figure}

\subsubsection{Table de vérité}

\begin{center}
\renewcommand{\arraystretch}{1.3}
\begin{tabular}{@{}cccc@{}}
\toprule
Entrée $\ket{a,b}$ & Contrôle $b$ & Sortie $\ket{a\oplus b,b}$ & Effet sur la cible $A$ \\
\midrule
$\ket{00}$ & $0$ & $\ket{00}$ & aucun \\
$\ket{01}$ & $1$ & $\ket{11}$ & $0\to1$ \\
$\ket{10}$ & $0$ & $\ket{10}$ & aucun \\
$\ket{11}$ & $1$ & $\ket{01}$ & $1\to0$ \\
\bottomrule
\end{tabular}
\end{center}

\subsubsection{Matrice du CNOT}

La $j$-ième colonne d'une matrice $4\times4$ est l'image du $j$-ième vecteur de base
($\ket{00},\ket{01},\ket{10},\ket{11}$, dans cet ordre). La table donne :
\begin{align*}
  \mathrm{CNOT}\ket{00}&=\ket{00}=\begin{pmatrix}1\\0\\0\\0\end{pmatrix},&
  \mathrm{CNOT}\ket{01}&=\ket{11}=\begin{pmatrix}0\\0\\0\\1\end{pmatrix},\\[1.5ex]
  \mathrm{CNOT}\ket{10}&=\ket{10}=\begin{pmatrix}0\\0\\1\\0\end{pmatrix},&
  \mathrm{CNOT}\ket{11}&=\ket{01}=\begin{pmatrix}0\\1\\0\\0\end{pmatrix},
\end{align*}
d'où, en juxtaposant ces quatre colonnes,
\[
  \mathrm{CNOT}=\begin{pmatrix}
    1 & 0 & 0 & 0\\ 0 & 0 & 0 & 1\\ 0 & 0 & 1 & 0\\ 0 & 1 & 0 & 0
  \end{pmatrix}.
\]

On retrouve cette matrice avec la formule $I\otimes\ket{0}\bra{0}+X\otimes\ket{1}\bra{1}$ et
la règle $A\otimes B=(A_{ij}B)$ (section~\ref{sec:tenseur}), où
$\ket{0}\bra{0}=\begin{pmatrix}1&0\\0&0\end{pmatrix}$ et
$\ket{1}\bra{1}=\begin{pmatrix}0&0\\0&1\end{pmatrix}$ :
\[
  I\otimes\ket{0}\bra{0}=\begin{pmatrix}1&0&0&0\\0&0&0&0\\0&0&1&0\\0&0&0&0\end{pmatrix},
  \qquad
  X\otimes\ket{1}\bra{1}=\begin{pmatrix}0&0&0&0\\0&0&0&1\\0&0&0&0\\0&1&0&0\end{pmatrix},
\]
et la somme de ces deux matrices est la matrice du CNOT.

\subsubsection{Action sur un état quelconque}

Pour $\ket{\Psi}=c_{00}\ket{00}+c_{01}\ket{01}+c_{10}\ket{10}+c_{11}\ket{11}$ :
\[
  \mathrm{CNOT}\begin{pmatrix}c_{00}\\c_{01}\\c_{10}\\c_{11}\end{pmatrix}
  =\begin{pmatrix}
    1\cdot c_{00}+0\cdot c_{01}+0\cdot c_{10}+0\cdot c_{11}\\
    0\cdot c_{00}+0\cdot c_{01}+0\cdot c_{10}+1\cdot c_{11}\\
    0\cdot c_{00}+0\cdot c_{01}+1\cdot c_{10}+0\cdot c_{11}\\
    0\cdot c_{00}+1\cdot c_{01}+0\cdot c_{10}+0\cdot c_{11}
  \end{pmatrix}
  =\begin{pmatrix}c_{00}\\c_{11}\\c_{10}\\c_{01}\end{pmatrix}.
\]
Le CNOT échange les amplitudes de $\ket{01}$ et de $\ket{11}$, c'est-à-dire des deux états
dont le contrôle vaut $1$, et laisse les autres.

\paragraph{Exemple : contrôle en superposition.} Soit $\ket{0}_A\otimes\ket{+}_B
=\frac{1}{\sqrt2}\big(\ket{00}+\ket{01}\big)=\frac{1}{\sqrt2}(1,1,0,0)^{\mathsf T}$,
un état séparable. Alors
\[
  \mathrm{CNOT}\,\frac{1}{\sqrt2}\begin{pmatrix}1\\1\\0\\0\end{pmatrix}
  =\frac{1}{\sqrt2}\begin{pmatrix}1\\0\\0\\1\end{pmatrix}
  =\frac{1}{\sqrt2}\big(\ket{00}+\ket{11}\big).
\]
Le résultat est intriqué ($c_{00}c_{11}-c_{01}c_{10}=\frac12\neq0$ ; section~\ref{sec:separable}) :
un CNOT dont le contrôle est en superposition crée de l'intrication. C'est l'état de Bell
$\ket{\Phi^+}$ du chapitre~\ref{chap:bell}.

\paragraph{Réversibilité.} Deux CNOT successifs redonnent l'état initial :
$\ket{a,b}\to\ket{a\oplus b,b}\to\ket{a\oplus b\oplus b,b}=\ket{a,b}$, donc
$\mathrm{CNOT}^2=\mathrm{Id}$. La matrice est réelle et symétrique, donc
$\mathrm{CNOT}^\dagger\mathrm{CNOT}=\mathrm{CNOT}^2=\mathrm{Id}$ : elle est bien unitaire.

\begin{remarque}[CNOT et copie d'un bit]
Avec la cible dans l'état $\ket{0}$, $\mathrm{CNOT}\ket{0,b}=\ket{b,b}$ pour $b=0$ et $b=1$ : le CNOT
copie le bit classique $b$. En revanche $\mathrm{CNOT}\big(\ket{0}\otimes\ket{+}\big)$ est l'état
ci-dessus, et non $\ket{+}\otimes\ket{+}$ : un état quantique inconnu ne peut pas être copié
(théorème de non-clonage, section~\ref{sec:teleportation}).
\end{remarque}

\begin{remarque}[Contrôle en dernier ou en premier]
La matrice dépend de l'ordre d'écriture du ket. Notons $\mathrm{CNOT}_{21}$ le CNOT dont le
contrôle est le \emph{second} chiffre du ket : c'est la matrice ci-dessus, adoptée dans ce
document (fil du haut). Beaucoup d'ouvrages mettent le contrôle en \emph{premier} ; le CNOT
correspondant est
\[
  \mathrm{CNOT}_{12}=\ket{0}\bra{0}\otimes I+\ket{1}\bra{1}\otimes X
  =\begin{pmatrix}1&0&0&0\\0&1&0&0\\0&0&0&1\\0&0&1&0\end{pmatrix}.
\]
Le même circuit a donc deux matrices, selon que le contrôle est écrit en dernier ou en premier.
Les deux sont reliées par le SWAP (section~\ref{sec:swap}). La section~\ref{sec:controlees} sur les
portes contrôlées met le contrôle en premier, avec la matrice $\mathrm{diag}(I,U)$.
\end{remarque}

% ---------------------------------------------------------------------
\subsection{Porte SWAP}\label{sec:swap}
% ---------------------------------------------------------------------

\begin{definition}{Porte SWAP (échange)}
La porte SWAP échange les états des deux qubits :
\[
  \mathrm{SWAP}\big(\ket{\psi}\otimes\ket{\phi}\big)=\ket{\phi}\otimes\ket{\psi}.
\]
Sur les états de base, $\mathrm{SWAP}\ket{a,b}=\ket{b,a}$ ; on prolonge par linéarité.
\end{definition}

Son symbole est formé de deux croix reliées par un trait vertical
(figure~\ref{fig:swap}).

\begin{figure}[!ht]
\centering
\begin{tikzpicture}[>=Stealth, font=\small]
  \draw (0,1.2) -- (6,1.2);
  \draw (0,0) -- (6,0);
  \node[anchor=east] at (0,1.2) {$\ket{a}$};
  \node[anchor=east] at (0,0) {$\ket{b}$};
  \node[anchor=west] at (6,1.2) {$\ket{b}$};
  \node[anchor=west] at (6,0) {$\ket{a}$};
  \draw[insarouge, thick] (3,1.2) -- (3,0);
  \pic at (3,1.2) {croix};
  \pic at (3,0) {croix};
\end{tikzpicture}
\caption{Porte SWAP : les deux qubits échangent leurs états.}
\label{fig:swap}
\end{figure}

\subsubsection{Table et matrice}

\begin{center}
\renewcommand{\arraystretch}{1.3}
\begin{tabular}{@{}cc@{}}
\toprule
Entrée & Sortie \\
\midrule
$\ket{00}$ & $\ket{00}$ \\
$\ket{01}$ & $\ket{10}$ \\
$\ket{10}$ & $\ket{01}$ \\
$\ket{11}$ & $\ket{11}$ \\
\bottomrule
\end{tabular}
\end{center}

Les images des quatre vecteurs de base sont les colonnes de la matrice :
\[
  \mathrm{SWAP}\ket{00}=\begin{pmatrix}1\\0\\0\\0\end{pmatrix},\quad
  \mathrm{SWAP}\ket{01}=\begin{pmatrix}0\\0\\1\\0\end{pmatrix},\quad
  \mathrm{SWAP}\ket{10}=\begin{pmatrix}0\\1\\0\\0\end{pmatrix},\quad
  \mathrm{SWAP}\ket{11}=\begin{pmatrix}0\\0\\0\\1\end{pmatrix},
\]
\[
  \mathrm{SWAP}=\begin{pmatrix}
    1 & 0 & 0 & 0\\ 0 & 0 & 1 & 0\\ 0 & 1 & 0 & 0\\ 0 & 0 & 0 & 1
  \end{pmatrix}.
\]
Appliquée à $\ket{01}=(0,1,0,0)^{\mathsf T}$, cette matrice sélectionne sa deuxième colonne :
\[
  \begin{pmatrix}
    1 & 0 & 0 & 0\\ 0 & 0 & 1 & 0\\ 0 & 1 & 0 & 0\\ 0 & 0 & 0 & 1
  \end{pmatrix}
  \begin{pmatrix}0\\1\\0\\0\end{pmatrix}
  =\begin{pmatrix}0\\0\\1\\0\end{pmatrix}=\ket{10}.
\]
Sur un état quelconque, $\mathrm{SWAP}\,(c_{00},c_{01},c_{10},c_{11})^{\mathsf T}
=(c_{00},c_{10},c_{01},c_{11})^{\mathsf T}$ : les amplitudes de $\ket{01}$ et $\ket{10}$
sont échangées.

\subsubsection{Accord avec la définition}

Pour deux états $\ket{\psi}=(a_0,a_1)^{\mathsf T}$ et $\ket{\phi}=(b_0,b_1)^{\mathsf T}$ :
\[
  \mathrm{SWAP}\,(\ket{\psi}\otimes\ket{\phi})
  =\mathrm{SWAP}\begin{pmatrix}a_0b_0\\a_0b_1\\a_1b_0\\a_1b_1\end{pmatrix}
  =\begin{pmatrix}a_0b_0\\a_1b_0\\a_0b_1\\a_1b_1\end{pmatrix}
  =\begin{pmatrix}b_0\\b_1\end{pmatrix}\otimes\begin{pmatrix}a_0\\a_1\end{pmatrix}
  =\ket{\phi}\otimes\ket{\psi},
\]
ce qui est bien la définition. La porte s'écrit aussi avec des projecteurs :
\[
  \mathrm{SWAP}=\sum_{a,b\in\{0,1\}}\ket{a}\bra{b}\otimes\ket{b}\bra{a}.
\]
En effet, pour $\ket{c}\otimes\ket{d}$ on trouve
$\sum_{a,b}\ket{a}\braket{b}{c}\otimes\ket{b}\braket{a}{d}$, et seul le terme $b=c$, $a=d$
est non nul : le résultat est $\ket{d}\otimes\ket{c}$.

\paragraph{Exemples.}
\begin{itemize}[nosep]
  \item Deux états différents : $\mathrm{SWAP}\big(\ket{0}\otimes\ket{+}\big)
        =\mathrm{SWAP}\,\frac{1}{\sqrt2}(1,1,0,0)^{\mathsf T}=\frac{1}{\sqrt2}(1,0,1,0)^{\mathsf T}
        =\ket{+}\otimes\ket{0}$.
  \item L'état $\ket{\psi}=\frac{1}{\sqrt2}\ket{00}+\frac{1}{\sqrt2}\ket{11}$ ne change pas :
        $\mathrm{SWAP}\ket{\psi}=\frac{1}{\sqrt2}\ket{00}+\frac{1}{\sqrt2}\ket{11}=\ket{\psi}$,
        car $\ket{00}$ et $\ket{11}$ sont invariants. De même
        $\frac{1}{\sqrt2}(\ket{00}-\ket{11})$ et $\frac{1}{\sqrt2}(\ket{01}+\ket{10})$ sont invariants,
        alors que $\mathrm{SWAP}\,\frac{1}{\sqrt2}(\ket{01}-\ket{10})=\frac{1}{\sqrt2}(\ket{10}-\ket{01})$
        est l'opposé de l'état de départ. Ce sont les quatre états de Bell (chapitre~\ref{chap:bell}).
\end{itemize}

\paragraph{Propriétés.} Échanger deux fois redonne l'état de départ :
$\mathrm{SWAP}^2=\mathrm{Id}$ ; la matrice est réelle et symétrique, donc
$\mathrm{SWAP}^\dagger\mathrm{SWAP}=\mathrm{SWAP}^2=\mathrm{Id}$ : elle est unitaire.

\subsubsection{Lien avec le CNOT}

\begin{resultat}{SWAP échange les rôles de contrôle et de cible}
\[
  \mathrm{CNOT}_{21}=\mathrm{SWAP}\;\mathrm{CNOT}_{12}\;\mathrm{SWAP}.
\]
\end{resultat}

\emph{Sur les états de base.}
$\mathrm{SWAP}\,\mathrm{CNOT}_{12}\,\mathrm{SWAP}\ket{a,b}
=\mathrm{SWAP}\,\mathrm{CNOT}_{12}\ket{b,a}
=\mathrm{SWAP}\ket{b,\,b\oplus a}=\ket{b\oplus a,\,b}=\mathrm{CNOT}_{21}\ket{a,b}$.

\emph{Par les matrices.} Multiplier à gauche par $\mathrm{SWAP}$ échange les lignes $2$ et
$3$, multiplier à droite l'échange les colonnes $2$ et $3$ :
\[
  \mathrm{SWAP}\,\mathrm{CNOT}_{12}
  =\begin{pmatrix}1&0&0&0\\0&0&0&1\\0&1&0&0\\0&0&1&0\end{pmatrix},
  \qquad
  \mathrm{SWAP}\,\mathrm{CNOT}_{12}\,\mathrm{SWAP}
  =\begin{pmatrix}1&0&0&0\\0&0&0&1\\0&0&1&0\\0&1&0&0\end{pmatrix}=\mathrm{CNOT}_{21}.
\]

\begin{remarque}[SWAP avec trois CNOT]
$\mathrm{SWAP}=\mathrm{CNOT}_{12}\,\mathrm{CNOT}_{21}\,\mathrm{CNOT}_{12}$. En effet, sur
$\ket{a,b}$ (on applique les portes de droite à gauche) :
\[
  \ket{a,b}\ \xrightarrow{\ \mathrm{CNOT}_{12}\ }\ \ket{a,\,a\oplus b}
  \ \xrightarrow{\ \mathrm{CNOT}_{21}\ }\ \ket{a\oplus(a\oplus b),\,a\oplus b}=\ket{b,\,a\oplus b}
  \ \xrightarrow{\ \mathrm{CNOT}_{12}\ }\ \ket{b,\,b\oplus(a\oplus b)}=\ket{b,a}.
\]
Les deux membres coïncident sur les états de base, donc sur tous les états par linéarité.
\end{remarque}

% ---------------------------------------------------------------------
\subsection{Portes contrôlées : contrôlé-\texorpdfstring{$U$}{U}, Fredkin et Toffoli}\label{sec:controlees}
% ---------------------------------------------------------------------

Le CNOT est le cas le plus simple d'une idée générale : appliquer une opération à un ou
plusieurs qubits \emph{cibles} seulement si un ou plusieurs qubits de \emph{contrôle} valent $1$.
Dans cette section, le contrôle est le \emph{premier} chiffre
du ket (remarque de la section~\ref{sec:cnot}). Les qubits sont notés $x$, $y$, $z$ (le qubit $x$ est le contrôle) et
$\ket{x,y,z}=\ket{x}\otimes\ket{y}\otimes\ket{z}$ ; d'après la convention de la section~\ref{sec:circuits}, le
contrôle est dessiné en bas (figure~\ref{fig:controlees}). Pour éviter la confusion avec les
matrices de Pauli, on écrit ici $\sigma_x=X$ pour la porte NOT.

\begin{figure}[!ht]
\centering
\begin{adjustbox}{max width=\linewidth}
\begin{tikzpicture}[>=Stealth, font=\small]
  % (a) contrôlé-U : contrôle x en bas, cible y en haut
  \draw (0,1.2) -- (3.4,1.2);
  \draw (0,0) -- (3.4,0);
  \node[anchor=east] at (0,1.2) {$y$};
  \node[anchor=east] at (0,0) {$x$};
  \draw[insarouge, thick] (1.7,0) -- (1.7,0.85);
  \pic at (1.7,0) {ctl};
  \node[porte] at (1.7,1.2) {$U$};
  \node[anchor=north, font=\footnotesize] at (1.7,-0.5) {(a) contrôlé-$U$};
  % (b) contrôlé-SWAP (Fredkin) : x contrôle, y et z échangés
  \draw (5,2.4) -- (8.4,2.4);
  \draw (5,1.2) -- (8.4,1.2);
  \draw (5,0) -- (8.4,0);
  \node[anchor=east] at (5,2.4) {$z$};
  \node[anchor=east] at (5,1.2) {$y$};
  \node[anchor=east] at (5,0) {$x$};
  \draw[insarouge, thick] (6.7,0) -- (6.7,2.4);
  \pic at (6.7,0) {ctl};
  \pic at (6.7,1.2) {croix};
  \pic at (6.7,2.4) {croix};
  \node[anchor=north, font=\footnotesize] at (6.7,-0.5) {(b) Fredkin (C-SWAP)};
  % (c) Toffoli (CCNOT) : x et y contrôlent, z cible
  \draw (10,2.4) -- (13.4,2.4);
  \draw (10,1.2) -- (13.4,1.2);
  \draw (10,0) -- (13.4,0);
  \node[anchor=east] at (10,2.4) {$z$};
  \node[anchor=east] at (10,1.2) {$y$};
  \node[anchor=east] at (10,0) {$x$};
  \draw[insarouge, thick] (11.7,0) -- (11.7,2.4);
  \pic at (11.7,0) {ctl};
  \pic at (11.7,1.2) {ctl};
  \pic at (11.7,2.4) {cible};
  \node[anchor=north, font=\footnotesize] at (11.7,-0.5) {(c) Toffoli (CCNOT)};
\end{tikzpicture}
\end{adjustbox}
\caption{Portes contrôlées. Le contrôle $x$ est le premier chiffre du ket, donc le fil du bas.
(a) Le qubit $y$ subit $U$ si $x=1$ ; pour $U=X$ on retrouve le CNOT. (b) Si $x=1$, les qubits
$y$ et $z$ sont échangés. (c) Si $x=y=1$, le qubit $z$ est inversé.}
\label{fig:controlees}
\end{figure}

\subsubsection{Contrôlé-\texorpdfstring{$U$}{U}}

\begin{definition}{Porte contrôlée-$U$}
Soit $U$ une matrice unitaire $2\times2$. Sur le système $(x,y)$, de contrôle $x$ et de cible $y$,
\[
  \mathrm{C}U=\ket{0}\bra{0}_x\otimes I_y+\ket{1}\bra{1}_x\otimes U_y .
\]
Le premier terme agit quand le contrôle est $\ket{0}$ : la cible ne change pas ($I$). Le second
agit quand le contrôle est $\ket{1}$ : la cible subit $U$.
\end{definition}

Comme $\ket{0}\bra{0}\ket{x}=\delta_{x0}\ket{0}$ et $\ket{1}\bra{1}\ket{x}=\delta_{x1}\ket{1}$,
\[
  \mathrm{C}U\big(\ket{x}\otimes\ket{y}\big)=\ket{x}\otimes U^x\ket{y},\qquad U^0=I,\ U^1=U .
\]
Par la règle $A\otimes B=(A_{ij}B)$, $\ket{0}\bra{0}\otimes I=\begin{pmatrix}I&0\\0&0\end{pmatrix}$
et $\ket{1}\bra{1}\otimes U=\begin{pmatrix}0&0\\0&U\end{pmatrix}$ (blocs $2\times2$), donc
\[
  \mathrm{C}U=\begin{pmatrix}I&0\\0&U\end{pmatrix}
  =\begin{pmatrix}1&0&0&0\\0&1&0&0\\0&0&u_{00}&u_{01}\\0&0&u_{10}&u_{11}\end{pmatrix}
  \qquad\text{pour }U=\begin{pmatrix}u_{00}&u_{01}\\u_{10}&u_{11}\end{pmatrix}.
\]
Cette matrice est unitaire car $I^\dagger I=I$ et $U^\dagger U=I$.

\paragraph{Le CNOT est un cas particulier.} Pour $U=\sigma_x$,
$\mathrm{C}\sigma_x=\ket{0}\bra{0}_x\otimes I_y+\ket{1}\bra{1}_x\otimes\sigma_x$ est le CNOT de
contrôle $x$ (premier chiffre), c'est-à-dire $\mathrm{CNOT}_{12}$ de la section~\ref{sec:cnot}.
On calcule l'action sur les quatre états de base, avec $\braket{0}{0}=\braket{1}{1}=1$,
$\braket{0}{1}=\braket{1}{0}=0$, $\sigma_x\ket{0}=\ket{1}$ et $\sigma_x\ket{1}=\ket{0}$ :
\begin{align*}
  \mathrm{C}\sigma_x\ket{00}&=\big[\ket{0}\bra{0}_x\otimes I_y\big]\ket{0}_x\ket{0}_y
      +\big[\ket{1}\bra{1}_x\otimes\sigma_x\big]\ket{0}_x\ket{0}_y\\
    &=\ket{0}\braket{0}{0}\otimes I\ket{0}+\ket{1}\braket{1}{0}\otimes\sigma_x\ket{0}
     =\ket{0}\otimes\ket{0}+0=\ket{00},\\[1.2ex]
  \mathrm{C}\sigma_x\ket{01}&=\ket{0}\braket{0}{0}\otimes I\ket{1}+\ket{1}\braket{1}{0}\otimes\sigma_x\ket{1}
     =\ket{0}\otimes\ket{1}+0=\ket{01},\\[1.2ex]
  \mathrm{C}\sigma_x\ket{10}&=\ket{0}\braket{0}{1}\otimes I\ket{0}+\ket{1}\braket{1}{1}\otimes\sigma_x\ket{0}
     =0+\ket{1}\otimes\ket{1}=\ket{11},\\[1.2ex]
  \mathrm{C}\sigma_x\ket{11}&=\ket{0}\braket{0}{1}\otimes I\ket{1}+\ket{1}\braket{1}{1}\otimes\sigma_x\ket{1}
     =0+\ket{1}\otimes\ket{0}=\ket{10}.
\end{align*}
On retrouve la table du CNOT : la cible n'est inversée que si le contrôle vaut $1$.

\paragraph{Un autre exemple : $U=Z$.} Ici
\[
  \mathrm{C}Z=\begin{pmatrix}I&0\\0&Z\end{pmatrix}=\mathrm{diag}(1,1,1,-1) :
\]
la porte multiplie $\ket{11}$ par $-1$ et ne touche pas les autres états de base. Elle est
symétrique entre les deux qubits. Elle se déduit du CNOT par
$\mathrm{C}Z=(I\otimes H)\,\mathrm{C}X\,(I\otimes H)$, où $\mathrm{C}X$ est le contrôlé-$X$.
En effet, $I\otimes H=\mathrm{diag}(H,H)$ (deux blocs $2\times2$ égaux à $H$), et le produit de
matrices diagonales par blocs se fait bloc par bloc :
\[
  (I\otimes H)\begin{pmatrix}I&0\\0&X\end{pmatrix}(I\otimes H)
  =\begin{pmatrix}H\,I\,H&0\\0&H\,X\,H\end{pmatrix}
  =\begin{pmatrix}I&0\\0&HXH\end{pmatrix},
\]
car $H^2=I$, et
\[
  HXH=\frac12\begin{pmatrix}1&1\\1&-1\end{pmatrix}\begin{pmatrix}0&1\\1&0\end{pmatrix}\begin{pmatrix}1&1\\1&-1\end{pmatrix}
  =\frac12\begin{pmatrix}1&1\\1&-1\end{pmatrix}\begin{pmatrix}1&-1\\1&1\end{pmatrix}
  =\frac12\begin{pmatrix}2&0\\0&-2\end{pmatrix}=Z
\]
(section~\ref{sec:hadamard} : $H$ échange les axes $x$ et $z$).

\subsubsection{Contrôlé-SWAP : la porte de Fredkin}

\begin{definition}{Porte de Fredkin (C-SWAP)}
Sur trois qubits $(x,y,z)$ (espace de dimension $2^3=8$), avec $x$ pour contrôle,
\[
  \mathrm{C\text{-}SWAP}=\ket{0}\bra{0}_x\otimes I_y\otimes I_z+\ket{1}\bra{1}_x\otimes\mathrm{SWAP}_{yz}.
\]
Si $x=0$, rien ne change ; si $x=1$, les qubits $y$ et $z$ sont échangés.
\end{definition}

Calcul sur quatre états de base (avec $I\ket{ab}=\ket{ab}$ et $\mathrm{SWAP}\ket{ab}=\ket{ba}$) :
\begin{align*}
  \mathrm{C\text{-}SWAP}\ket{000}&=\ket{0}\braket{0}{0}\otimes I\ket{00}+\ket{1}\braket{1}{0}\otimes\mathrm{SWAP}\ket{00}
     =\ket{0}\otimes\ket{00}+0=\ket{000},\\[1.2ex]
  \mathrm{C\text{-}SWAP}\ket{101}&=\ket{0}\braket{0}{1}\otimes I\ket{01}+\ket{1}\braket{1}{1}\otimes\mathrm{SWAP}\ket{01}
     =0+\ket{1}\otimes\ket{10}=\ket{110},\\[1.2ex]
  \mathrm{C\text{-}SWAP}\ket{110}&=\ket{0}\braket{0}{1}\otimes I\ket{10}+\ket{1}\braket{1}{1}\otimes\mathrm{SWAP}\ket{10}
     =0+\ket{1}\otimes\ket{01}=\ket{101},\\[1.2ex]
  \mathrm{C\text{-}SWAP}\ket{010}&=\ket{0}\braket{0}{0}\otimes I\ket{10}+\ket{1}\braket{1}{0}\otimes\mathrm{SWAP}\ket{10}
     =\ket{0}\otimes\ket{10}+0=\ket{010}.
\end{align*}
Les huit états de base donnent :
\begin{center}
\renewcommand{\arraystretch}{1.3}
\begin{tabular}{@{}ccccccccc@{}}
\toprule
Entrée & $\ket{000}$ & $\ket{001}$ & $\ket{010}$ & $\ket{011}$ & $\ket{100}$ & $\ket{101}$ & $\ket{110}$ & $\ket{111}$ \\
\midrule
Sortie & $\ket{000}$ & $\ket{001}$ & $\ket{010}$ & $\ket{011}$ & $\ket{100}$ & $\ket{110}$ & $\ket{101}$ & $\ket{111}$ \\
\bottomrule
\end{tabular}
\end{center}
Seuls $\ket{101}$ et $\ket{110}$ changent (contrôle à $1$ et $y\neq z$) ; pour $\ket{100}$ et
$\ket{111}$ l'échange de deux qubits égaux ne change rien. Dans la base
$\ket{000},\ket{001},\dots,\ket{111}$ (ordre binaire), la matrice est
$\mathrm{diag}(I_4,\mathrm{SWAP})$ (deux blocs $4\times4$) :
\[
  \mathrm{C\text{-}SWAP}=\begingroup\setlength{\arraycolsep}{3pt}\begin{pmatrix}
    1&0&0&0&0&0&0&0\\
    0&1&0&0&0&0&0&0\\
    0&0&1&0&0&0&0&0\\
    0&0&0&1&0&0&0&0\\
    0&0&0&0&1&0&0&0\\
    0&0&0&0&0&0&1&0\\
    0&0&0&0&0&1&0&0\\
    0&0&0&0&0&0&0&1
  \end{pmatrix}\endgroup .
\]
C'est l'identité dont les lignes $6$ et $7$ sont échangées.

\paragraph{Contrôle en superposition.} Avec $\ket{+}_x\otimes\ket{0}_y\otimes\ket{1}_z
=\frac{1}{\sqrt2}\big(\ket{001}+\ket{101}\big)$ :
\[
  \mathrm{C\text{-}SWAP}\,\frac{1}{\sqrt2}\big(\ket{001}+\ket{101}\big)
  =\frac{1}{\sqrt2}\big(\ket{001}+\ket{110}\big).
\]
Le résultat n'est pas un produit : les coefficients de $\ket{x}\otimes\ket{yz}$ forment la
matrice (lignes $x=0,1$, colonnes $yz=00,01,10,11$)
\[
  \frac{1}{\sqrt2}\begin{pmatrix}0&1&0&0\\0&0&1&0\end{pmatrix},
\]
dont les deux lignes ne sont pas proportionnelles. Le qubit $x$ est intriqué avec la paire
$(y,z)$.

\subsubsection{Contrôlé-contrôlé-NOT : la porte de Toffoli}

\begin{definition}{Porte de Toffoli (CCNOT)}
Sur trois qubits $(x,y,z)$, avec $x$ et $y$ pour contrôles et $z$ pour cible,
\[
  \mathrm{CCNOT}=\ket{0}\bra{0}_x\otimes I_y\otimes I_z
   +\ket{1}\bra{1}_x\otimes\Big[\ket{0}\bra{0}_y\otimes I_z+\ket{1}\bra{1}_y\otimes\sigma_x\Big].
\]
Le crochet est un CNOT (contrôle $y$, cible $z$), appliqué seulement si $x=1$ : si $x=0$, rien ;
si $x=1$ et $y=0$, rien ; si $x=1$ et $y=1$, le qubit $z$ est inversé.
\end{definition}

Calcul de $\mathrm{CCNOT}\ket{110}$ :
\begin{align*}
  \mathrm{CCNOT}\ket{110}&=\ket{0}\braket{0}{1}\otimes I\ket{1}\otimes I\ket{0}
     +\ket{1}\braket{1}{1}\otimes\Big[\ket{0}\braket{0}{1}\otimes I\ket{0}
       +\ket{1}\braket{1}{1}\otimes\sigma_x\ket{0}\Big]\\
  &=0+\ket{1}\otimes\Big[0+\ket{1}\otimes\ket{1}\Big]=\ket{111}.
\end{align*}
Pour les huit états de base, $\mathrm{CCNOT}\ket{x,y,z}=\ket{x,\,y,\,z\oplus xy}$ :
\begin{center}
\renewcommand{\arraystretch}{1.3}
\begin{tabular}{@{}ccccccccc@{}}
\toprule
Entrée & $\ket{000}$ & $\ket{001}$ & $\ket{010}$ & $\ket{011}$ & $\ket{100}$ & $\ket{101}$ & $\ket{110}$ & $\ket{111}$ \\
\midrule
Sortie & $\ket{000}$ & $\ket{001}$ & $\ket{010}$ & $\ket{011}$ & $\ket{100}$ & $\ket{101}$ & $\ket{111}$ & $\ket{110}$ \\
\bottomrule
\end{tabular}
\end{center}
Seuls $\ket{110}$ et $\ket{111}$ sont échangés : la matrice est $\mathrm{diag}(I_6,X)$
(un bloc $6\times6$ égal à l'identité, puis $X$),
\[
  \mathrm{CCNOT}=\begingroup\setlength{\arraycolsep}{3pt}\begin{pmatrix}
    1&0&0&0&0&0&0&0\\
    0&1&0&0&0&0&0&0\\
    0&0&1&0&0&0&0&0\\
    0&0&0&1&0&0&0&0\\
    0&0&0&0&1&0&0&0\\
    0&0&0&0&0&1&0&0\\
    0&0&0&0&0&0&0&1\\
    0&0&0&0&0&0&1&0
  \end{pmatrix}\endgroup .
\]

\begin{remarque}[Portes réversibles]
Les portes de Fredkin et de Toffoli sont des permutations des états de base qui échangent
seulement deux d'entre eux : elles sont leurs propres inverses ($\mathrm{C\text{-}SWAP}^2=\mathrm{CCNOT}^2=\mathrm{Id}$)
et unitaires. Avec $z=0$, la Toffoli calcule le ET logique $x\wedge y$ sans rien effacer :
$\ket{x,y,0}\to\ket{x,y,xy}$ ; avec $z=1$, elle calcule le NON-ET. Elle sert à écrire tout calcul
classique de façon réversible.
\end{remarque}

\clearpage
% ---------------------------------------------------------------------
\subsection{Récapitulatif du chapitre}\label{sec:recap4}
% ---------------------------------------------------------------------

\begin{center}
\adjustbox{max width=\linewidth}{%
\begin{tikzpicture}[>=Stealth, font=\footnotesize,
    boite/.style={draw=insarouge, fill=insarouge!6, rounded corners=2pt, thick,
                  minimum width=4.4cm, minimum height=1.5cm, align=center, inner sep=3pt},
    lien/.style={->, very thick, insagris},
    etiq/.style={font=\scriptsize\itshape, text=insagris, inner sep=2pt}]
  \node[boite] (uni) at (0,0)
    {\textbf{Porte unitaire}\\ $U^\dagger U=\mathrm{Id}$, \ $\ket{\Psi}\to U\ket{\Psi}$\\ {\scriptsize\S\,\ref{sec:portes}}};
  \node[boite] (un) at (5.9,0)
    {\textbf{Portes à un qubit}\\ $X$, $Y$, $Z$, $H$, $S$, $T$\\ {\scriptsize\S\,\ref{sec:pauli}--\ref{sec:phase}}};
  \node[boite] (blo) at (11.8,0)
    {\textbf{Rotations de Bloch}\\ $X,Y,Z,H$ : angle $\pi$ ; $P_\theta$ autour de $z$\\
     {\scriptsize\S\,\ref{sec:portes-bloch}}};
  \node[boite] (cir) at (11.8,-2.5)
    {\textbf{Circuits}\\ fils, boîtes, $U_2U_1$\\ {\scriptsize\S\,\ref{sec:circuits}, \ref{sec:op2}}};
  \node[boite] (cn) at (5.9,-2.5)
    {\textbf{CNOT et SWAP}\\ $\ket{a,b}\to\ket{a\oplus b,b}$\\ {\scriptsize\S\,\ref{sec:cnot}, \ref{sec:swap}}};
  \node[boite] (cu) at (0,-2.5)
    {\textbf{Portes contrôlées}\\ $\mathrm{C}U=\mathrm{diag}(I,U)$, Fredkin, Toffoli\\
     {\scriptsize\S\,\ref{sec:controlees}}};
  \draw[lien] (uni) -- (un);
  \draw[lien] (un) -- (blo);
  \draw[lien] (blo) -- (cir);
  \draw[lien] (cir) -- (cn) node[etiq, midway, above=1pt] {2 qubits};
  \draw[lien] (cn) -- (cu) node[etiq, midway, above=1pt] {contrôle};
\end{tikzpicture}}
\end{center}

\begin{aretenir}
\begin{itemize}
  \item Une porte est un opérateur unitaire : elle conserve la norme, elle est réversible
        ($U^{-1}=U^\dagger$) et agit sur la sphère de Bloch comme une rotation. C'est l'évolution de
        Schrödinger, $U=\ee^{-\ii\Hop t/\hbar}$.
  \item $X$ échange $\ket{0}$ et $\ket{1}$, $Z$ change le signe de $\ket{1}$, $Y=\ii XZ$ fait les deux ;
        $H$ échange les bases $Z$ et $X$ et crée des superpositions ; $P_\theta$, $S=P_{\pi/2}$ et
        $T=P_{\pi/4}$ font tourner autour de $z$ ($T^2=S$, $S^2=Z$).
  \item Dans un circuit, le temps va de gauche à droite et la matrice est le produit dans l'ordre
        inverse, $U_2U_1$. Le fil du haut porte le dernier chiffre du ket. Des portes en parallèle
        forment un produit tensoriel $U_A\otimes U_B$.
  \item Le CNOT inverse la cible si le contrôle vaut $1$ : $\ket{a,b}\to\ket{a\oplus b,b}$. Le SWAP
        échange deux qubits ; $\mathrm{CNOT}_{21}=\mathrm{SWAP}\,\mathrm{CNOT}_{12}\,\mathrm{SWAP}$.
  \item Une porte contrôlée a pour matrice $\mathrm{diag}(I,U)$. Fredkin (contrôlé-SWAP) et
        Toffoli (CCNOT) sont des permutations des états de base : réversibles et universelles
        pour le calcul classique.
\end{itemize}
\end{aretenir}

\begin{center}
\renewcommand{\arraystretch}{1.45}
\begin{tabular}{@{}l l r@{}}
\toprule
\textbf{Notion} & \textbf{Formule} & \textbf{§} \\
\midrule
Porte unitaire & $U^\dagger U=\mathrm{Id}$, \quad $\bra{\Psi}U^\dagger U\ket{\Psi}=1$ & \ref{sec:portes} \\
Pauli & $X\ket{0}=\ket{1}$, \quad $Z\ket{1}=-\ket{1}$, \quad $X^2=Y^2=Z^2=\mathrm{Id}$ & \ref{sec:pauli} \\
Hadamard & $H\ket{0}=\ket{+}$, \quad $H\ket{1}=\ket{-}$, \quad $H^2=\mathrm{Id}$ & \ref{sec:hadamard} \\
Phase & $P_\theta=\mathrm{diag}(1,\ee^{\ii\theta})$, \quad $S=P_{\pi/2}$, \quad $T=P_{\pi/4}$ & \ref{sec:phase} \\
Circuit & $U_1$ puis $U_2$ : \quad $U_2U_1$ & \ref{sec:circuits} \\
Portes locales & $(U\otimes I)(\ket{u}\otimes\ket{v})=U\ket{u}\otimes\ket{v}$ & \ref{sec:op2} \\
CNOT & $\ket{a,b}\to\ket{a\oplus b,b}$, \quad $\mathrm{CNOT}=I\otimes\ket{0}\bra{0}+X\otimes\ket{1}\bra{1}$
  & \ref{sec:cnot} \\
SWAP & $\mathrm{SWAP}\ket{a,b}=\ket{b,a}$, \quad $\mathrm{SWAP}=\mathrm{CNOT}_{12}\mathrm{CNOT}_{21}\mathrm{CNOT}_{12}$
  & \ref{sec:swap} \\
Contrôlé-$U$ & $\mathrm{C}U=\ket{0}\bra{0}\otimes I+\ket{1}\bra{1}\otimes U=\mathrm{diag}(I,U)$ & \ref{sec:controlees} \\
Toffoli & $\mathrm{CCNOT}\ket{x,y,z}=\ket{x,y,z\oplus xy}$ & \ref{sec:controlees} \\
\bottomrule
\end{tabular}
\end{center}

\begin{savoirfaire}
\begin{itemize}
  \item Appliquer une porte à un état (par la matrice ou par les ket-bra) et identifier sa rotation
        de Bloch (exercices~\ref{ex:conj} et~\ref{ex:phases}).
  \item Lire un circuit, écrire sa matrice dans le bon ordre et suivre l'état à chaque étape
        (exercice~\ref{ex:kickback}).
  \item Écrire la matrice de $U\otimes I$, du CNOT, du SWAP ou d'une porte contrôlée, et calculer
        leur action sur les états de base.
  \item Reconnaître une porte réversible qui calcule une fonction booléenne
        (exercice~\ref{ex:toffoli}).
\end{itemize}
\end{savoirfaire}
.
%  Il ne contient ni préambule ni \begin{document} : les environnements
%  (definition, resultat, remarque, aretenir, savoirfaire, exercice, correction),
%  les macros (\ii, \ee, \dd, \pd, \pdd, \Hop, \ket, \bra, \braket), les symboles
%  de circuits (ctl, cible, croix, mesure, porte) et les couleurs (insarouge,
%  insagris) sont définis dans main.tex. Un chapitre = une \section.
% =====================================================================

% Macros de Dirac (déjà définies dans main.tex ; ce bloc évite l'erreur
% « Undefined control sequence \ket » si main.tex est une ancienne version).
\providecommand{\ket}[1]{|#1\rangle}
\providecommand{\bra}[1]{\langle #1|}
\providecommand{\braket}[2]{\langle #1|#2\rangle}

% =====================================================================
\section{Portes quantiques et systèmes multi-qubits}\label{chap:portes}
% =====================================================================

Un calcul quantique est une suite de transformations unitaires appliquées à des qubits. Ce
chapitre décrit les portes à un qubit (Pauli, Hadamard, phase) et leur effet sur la sphère de
Bloch, la lecture d'un circuit, puis les portes qui agissent sur plusieurs qubits : CNOT, SWAP et
portes contrôlées.

% ---------------------------------------------------------------------
\subsection{Portes quantiques}\label{sec:portes}
% ---------------------------------------------------------------------

Un calcul quantique transforme des états. Les opérations sur un seul qubit sont appelées
\emph{portes} ; les sections suivantes passent à plusieurs qubits.

\begin{resultat}{Postulat 4 : évolution d'un système isolé}
L'évolution d'un système isolé est décrite par un opérateur unitaire : $\ket{\Psi}\to U\ket{\Psi}$
avec $U^\dagger U=\mathrm{Id}$.
\end{resultat}

C'est la forme de l'équation de Schrödinger du chapitre~\ref{chap:histoire}. Pour un hamiltonien
$\Hop$ indépendant du temps, $\ii\hbar\,\partial_t\ket{\Psi}=\Hop\ket{\Psi}$ a pour solution
$\ket{\Psi(t)}=U(t)\ket{\Psi(0)}$ avec $U(t)=\ee^{-\ii\Hop t/\hbar}$. Comme $\Hop$ est hermitien,
$U(t)^\dagger=\ee^{+\ii\Hop t/\hbar}$ et $U^\dagger U=\mathrm{Id}$ : l'évolution est unitaire. Une
porte est une telle évolution, obtenue en contrôlant $\Hop$ pendant une durée donnée.

\begin{definition}{Porte quantique}
Une porte à un qubit est un opérateur $U$ (matrice $2\times2$) \emph{unitaire} :
$U^\dagger U=\mathrm{Id}$. Elle transforme $\ket{\Psi}\to U\ket{\Psi}$, conserve la norme
($\bra{\Psi}U^\dagger U\ket{\Psi}=\braket{\Psi}{\Psi}=1$) et est réversible
($U^{-1}=U^\dagger$). Sur la sphère de Bloch (section~\ref{sec:bloch}), c'est une
\emph{rotation}.
\end{definition}

% ---------------------------------------------------------------------
\subsection{Portes de Pauli}\label{sec:pauli}
% ---------------------------------------------------------------------

Les matrices de Pauli de la section~\ref{sec:op-dirac} sont trois portes.

\begin{definition}{Portes $X$, $Y$, $Z$}
\begin{center}
\renewcommand{\arraystretch}{1}
\begin{tabular}{@{}lccl@{}}
\textbf{Porte} & \textbf{Matrice} & \textbf{Ket-bra} & \textbf{Nom} \\[0.8ex]
$X=\sigma_x$ & $\begin{pmatrix}0 & 1\\ 1 & 0\end{pmatrix}$
  & $\ket{0}\bra{1}+\ket{1}\bra{0}$ & bit flip \\[2.2ex]
$Y=\sigma_y$ & $\begin{pmatrix}0 & -\ii\\ \ii & 0\end{pmatrix}$
  & $-\ii\ket{0}\bra{1}+\ii\ket{1}\bra{0}$ &  \\[2.2ex]
$Z=\sigma_z$ & $\begin{pmatrix}1 & 0\\ 0 & -1\end{pmatrix}$
  & $\ket{0}\bra{0}-\ket{1}\bra{1}$ & phase flip \\
\end{tabular}
\end{center}
\end{definition}

\paragraph{Porte $X$ (inverseur).} Avec $\braket{0}{0}=\braket{1}{1}=1$ et
$\braket{0}{1}=\braket{1}{0}=0$ :
\begin{align*}
  X\ket{0}&=\ket{0}\braket{1}{0}+\ket{1}\braket{0}{0}=0\cdot\ket{0}+1\cdot\ket{1}=\ket{1},\\
  X\ket{1}&=\ket{0}\braket{1}{1}+\ket{1}\braket{0}{1}=1\cdot\ket{0}+0\cdot\ket{1}=\ket{0}.
\end{align*}
$X$ échange $\ket{0}$ et $\ket{1}$ : c'est le NOT quantique (\emph{bit flip}). Sur la
sphère de Bloch, c'est une rotation de $\pi$ autour de l'axe $x$ : le pôle nord devient
le pôle sud. Les états de l'axe $x$ ne bougent pas : $X\ket{+}=\ket{+}$ et
$X\ket{-}=-\ket{-}$.

\paragraph{Porte $Z$ (changement de phase).}
\begin{align*}
  Z\ket{0}&=\ket{0}\braket{0}{0}-\ket{1}\braket{1}{0}=\ket{0},\\
  Z\ket{1}&=\ket{0}\braket{0}{1}-\ket{1}\braket{1}{1}=-\ket{1},
\end{align*}
donc $Z$ change le signe de la composante sur $\ket{1}$ (\emph{phase flip}). Sur les
états de la base $X$, par linéarité :
\begin{align*}
  Z\ket{+}&=\tfrac{1}{\sqrt2}\big(Z\ket{0}+Z\ket{1}\big)=\tfrac{1}{\sqrt2}\big(\ket{0}-\ket{1}\big)=\ket{-},\\
  Z\ket{-}&=\tfrac{1}{\sqrt2}\big(Z\ket{0}-Z\ket{1}\big)=\tfrac{1}{\sqrt2}\big(\ket{0}+\ket{1}\big)=\ket{+}.
\end{align*}
Sous forme matricielle, $Z\ket{+}=\begin{pmatrix}1 & 0\\ 0 & -1\end{pmatrix}
\frac{1}{\sqrt2}\begin{pmatrix}1\\ 1\end{pmatrix}=\frac{1}{\sqrt2}\begin{pmatrix}1\\ -1\end{pmatrix}=\ket{-}$.
$Z$ est une rotation de $\pi$ autour de l'axe $z$ : elle laisse $\ket{0}$ et $\ket{1}$ en
place et échange $\ket{+}$ et $\ket{-}$.

\paragraph{Porte $Y$.} Avec $Y=-\ii\ket{0}\bra{1}+\ii\ket{1}\bra{0}$ :
\begin{align*}
  Y\ket{0}&=-\ii\ket{0}\braket{1}{0}+\ii\ket{1}\braket{0}{0}=\ii\ket{1},\\
  Y\ket{1}&=-\ii\ket{0}\braket{1}{1}+\ii\ket{1}\braket{0}{1}=-\ii\ket{0}.
\end{align*}
Sur les états de la base $X$ :
\begin{align*}
  Y\ket{+}&=\tfrac{1}{\sqrt2}\big(\ii\ket{1}-\ii\ket{0}\big)=-\ii\,\tfrac{1}{\sqrt2}\big(\ket{0}-\ket{1}\big)=-\ii\ket{-},\\
  Y\ket{-}&=\tfrac{1}{\sqrt2}\big(\ii\ket{1}+\ii\ket{0}\big)=\ii\,\tfrac{1}{\sqrt2}\big(\ket{0}+\ket{1}\big)=\ii\ket{+}.
\end{align*}
Un facteur global ($-\ii$ ou $\ii$) ne change pas l'état physique : $Y$ échange bien
$\ket{+}$ et $\ket{-}$. Sur les états de l'axe $y$ :
\begin{align*}
  Y\ket{{+}\ii}&=\tfrac{1}{\sqrt2}\big(Y\ket{0}+\ii Y\ket{1}\big)
   =\tfrac{1}{\sqrt2}\big(\ii\ket{1}+\ket{0}\big)=\ket{{+}\ii},\\
  Y\ket{{-}\ii}&=\tfrac{1}{\sqrt2}\big(Y\ket{0}-\ii Y\ket{1}\big)
   =\tfrac{1}{\sqrt2}\big(\ii\ket{1}-\ket{0}\big)=-\ket{{-}\ii}.
\end{align*}
$Y$ est une rotation de $\pi$ autour de l'axe $y$.

\begin{remarque}[Chaque Pauli fixe son axe]
$Z\ket{0}=\ket{0}$, $Z\ket{1}=-\ket{1}$ ; $X\ket{\pm}=\pm\ket{\pm}$ ; $Y\ket{{\pm}\ii}=\pm\ket{{\pm}\ii}$.
Les bases $Z$, $X$, $Y$ de la section~\ref{sec:bloch} sont les bases propres de
$Z$, $X$, $Y$. Une porte de Pauli laisse donc son propre axe invariant et retourne les
deux autres ; c'est pourquoi les exemples utilisent les états des deux autres axes.
\end{remarque}

% ---------------------------------------------------------------------
\subsection{Porte de Hadamard}\label{sec:hadamard}
% ---------------------------------------------------------------------

\begin{definition}{Porte de Hadamard}
\[
  H=\frac{1}{\sqrt2}\begin{pmatrix}1 & 1\\ 1 & -1\end{pmatrix}
  =\ket{+}\bra{0}+\ket{-}\bra{1}.
\]
\end{definition}

\begin{align*}
  H\ket{0}&=\frac{1}{\sqrt2}\begin{pmatrix}1 & 1\\ 1 & -1\end{pmatrix}\begin{pmatrix}1\\ 0\end{pmatrix}
    =\frac{1}{\sqrt2}\begin{pmatrix}1\\ 1\end{pmatrix}=\ket{+},\\[0.6ex]
  H\ket{1}&=\frac{1}{\sqrt2}\begin{pmatrix}1 & 1\\ 1 & -1\end{pmatrix}\begin{pmatrix}0\\ 1\end{pmatrix}
    =\frac{1}{\sqrt2}\begin{pmatrix}1\\ -1\end{pmatrix}=\ket{-},\\[0.6ex]
  H\ket{+}&=\frac{1}{2}\begin{pmatrix}1 & 1\\ 1 & -1\end{pmatrix}\begin{pmatrix}1\\ 1\end{pmatrix}
    =\frac{1}{2}\begin{pmatrix}2\\ 0\end{pmatrix}=\ket{0},\qquad
  H\ket{-}=\frac{1}{2}\begin{pmatrix}0\\ 2\end{pmatrix}=\ket{1}.
\end{align*}
Comme $H\ket{+}=\ket{0}$ et $H\ket{-}=\ket{1}$, on a $H^2=\mathrm{Id}$ ($H$ est son propre
inverse). Sur la sphère de Bloch, $H$ est une rotation de $\pi$ autour de l'axe
$(x+z)/\sqrt2$ : elle échange l'axe $x$ et l'axe $z$.

\begin{remarque}[À quoi sert $H$ ?]
\begin{itemize}[nosep]
  \item Elle \emph{crée une superposition} à partir d'un état de base :
        $H\ket{0}=\ket{+}$ donne $P(0)=P(1)=\frac12$. C'est la première étape de
        la plupart des algorithmes quantiques.
  \item Elle \emph{change de base} : elle envoie la base $Z$ sur la base $X$ et
        inversement. Mesurer dans la base $X$ revient à appliquer $H$ puis à mesurer
        dans la base $Z$ (figure~\ref{fig:mesurex}).
  \item Elle est réversible ($H^2=\mathrm{Id}$).
\end{itemize}
\end{remarque}

% ---------------------------------------------------------------------
\subsection{Portes de phase}\label{sec:phase}
% ---------------------------------------------------------------------

\begin{definition}{Portes de phase $P_\theta$, $S$ et $T$}
\[
  P_\theta=\begin{pmatrix}1 & 0\\ 0 & \ee^{\ii\theta}\end{pmatrix}
  =\ket{0}\bra{0}+\ee^{\ii\theta}\ket{1}\bra{1},
  \qquad
  S=P_{\pi/2}=\begin{pmatrix}1 & 0\\ 0 & \ii\end{pmatrix},
  \qquad
  T=P_{\pi/4}=\begin{pmatrix}1 & 0\\ 0 & \dfrac{1+\ii}{\sqrt2}\end{pmatrix}.
\]
\end{definition}

Ici $\ee^{\ii\pi/2}=\ii$ et $\ee^{\ii\pi/4}=\cos\frac\pi4+\ii\sin\frac\pi4=\frac{1+\ii}{\sqrt2}$.
$P_\theta$ laisse $\ket{0}$ inchangé et multiplie $\ket{1}$ par $\ee^{\ii\theta}$ :
\[
  P_\theta\big(\alpha\ket{0}+\beta\ket{1}\big)=\alpha\ket{0}+\ee^{\ii\theta}\beta\ket{1}.
\]
Dans la paramétrisation de Bloch, $\varphi\to\varphi+\theta$ : $P_\theta$ est une rotation
d'angle $\theta$ autour de l'axe $z$. On retrouve $P_\pi=Z$, et $S^2=Z$, $T^2=S$.

\paragraph{Exemples.}
\begin{align*}
  S\ket{+}&=\tfrac{1}{\sqrt2}\big(\ket{0}+\ii\ket{1}\big)=\ket{{+}\ii},\\
  S\ket{{+}\ii}&=\tfrac{1}{\sqrt2}\big(\ket{0}+\ii\cdot\ii\ket{1}\big)=\tfrac{1}{\sqrt2}\big(\ket{0}-\ket{1}\big)=\ket{-},\\
  T\ket{+}&=\tfrac{1}{\sqrt2}\begin{pmatrix}1 & 0\\ 0 & \frac{1+\ii}{\sqrt2}\end{pmatrix}\begin{pmatrix}1\\ 1\end{pmatrix}
    =\tfrac{1}{\sqrt2}\begin{pmatrix}1\\ \frac{1+\ii}{\sqrt2}\end{pmatrix}
    =\tfrac{1}{\sqrt2}\big(\ket{0}+\ee^{\ii\pi/4}\ket{1}\big)\quad\big(\text{angle polaire }\tfrac\pi2,\ \text{azimut }\tfrac\pi4\big).
\end{align*}
$S$ tourne donc $\ket{+}$ d'un quart de tour autour de $z$ (de l'axe $x$ à l'axe $y$),
puis d'un second quart de tour jusqu'à $\ket{-}$ ; $T$ fait un huitième de tour.

\begin{remarque}[Une porte de phase est une évolution libre]
Pour $\Hop=\frac{\hbar\omega}{2}Z$, $U(t)=\ee^{-\ii\Hop t/\hbar}=\mathrm{diag}\big(\ee^{-\ii\omega t/2},\ee^{\ii\omega t/2}\big)
=\ee^{-\ii\omega t/2}\,P_{\omega t}$. Au facteur global près, laisser évoluer le qubit pendant la
durée $t$ applique la porte de phase $P_{\omega t}$ : une rotation d'angle $\omega t$ autour de
l'axe $z$.
\end{remarque}

% ---------------------------------------------------------------------
\subsection{Portes et sphère de Bloch}\label{sec:portes-bloch}
% ---------------------------------------------------------------------

Toute porte à un qubit est une rotation de la sphère de Bloch. La figure~\ref{fig:portes} montre
les axes des portes usuelles et le tableau qui suit les résume.

\begin{figure}[!ht]
\centering
\begin{tikzpicture}[x={(-1.63cm,-0.56cm)}, y={(1.63cm,-0.56cm)}, z={(0cm,2.16cm)}, >=Stealth, scale=0.9]
  % sphère et grands cercles
  \shade[ball color=insagris!15, opacity=0.55] (0,0,0) circle[radius=2.3cm];
  \begin{scope}[insagris!75, thin]
    \draw[domain=-45:135, samples=60, variable=\t] plot ({cos(\t)},{sin(\t)},0);
    \draw[domain=135:315, samples=60, variable=\t, dashed] plot ({cos(\t)},{sin(\t)},0);
    \draw[domain=-62.7:117.3, samples=60, variable=\t] plot ({cos(\t)},0,{sin(\t)});
    \draw[domain=117.3:297.3, samples=60, variable=\t, dashed] plot ({cos(\t)},0,{sin(\t)});
    \draw[domain=-62.7:117.3, samples=60, variable=\t] plot (0,{cos(\t)},{sin(\t)});
    \draw[domain=117.3:297.3, samples=60, variable=\t, dashed] plot (0,{cos(\t)},{sin(\t)});
  \end{scope}
  % axes x, y, z (portes X, Y, Z)
  \draw[insagris, dashed] (0,0,0) -- (-1,0,0);
  \draw[insagris, dashed] (0,0,0) -- (0,-1,0);
  \draw[insagris, dashed] (0,0,0) -- (0,0,-1);
  \draw[insagris] (-1,0,0) -- (-1.3,0,0);
  \draw[insagris] (0,-1,0) -- (0,-1.3,0);
  \draw[insagris] (0,0,-1) -- (0,0,-1.3);
  \draw[insagris, -{Stealth[length=2.2mm]}] (0,0,0) -- (1.3,0,0);
  \draw[insagris, -{Stealth[length=2.2mm]}] (0,0,0) -- (0,1.3,0);
  \draw[insagris, -{Stealth[length=2.2mm]}] (0,0,0) -- (0,0,1.3);
  \node[anchor=east]  at (1.33,0,0) {$x$};
  \node[anchor=west]  at (0,1.33,0) {$y$};
  \node[anchor=south] at (0,0,1.33) {$z$};
  % axe de la porte H : (x+z)/sqrt(2)
  \draw[insarouge, dashed] (0,0,0) -- (-0.707,0,-0.707);
  \draw[insarouge, very thick, -{Stealth[length=2.6mm]}] (0,0,0) -- (0.98,0,0.98);
  \node[insarouge, anchor=south east] at (0.98,0,0.98) {axe de $H$};
  % rotation autour de z (portes de phase)
  \draw[insarouge, thick, -{Stealth[length=2.2mm]}, domain=15:105, samples=30, variable=\t]
    plot ({1.15*cos(\t)},{1.15*sin(\t)},0);
  \node[insarouge, anchor=north west, font=\small] at ({1.15*cos(60)},{1.15*sin(60)},0) {$P_\theta$};
\end{tikzpicture}
\caption{Portes et rotations sur la sphère de Bloch. Les portes $X$, $Y$, $Z$ sont des
rotations de $\pi$ autour des axes $x$, $y$, $z$ ; $H$ est une rotation de $\pi$ autour de
l'axe $(x+z)/\sqrt2$ (rouge) ; $P_\theta$ (donc $Z$, $S$, $T$) tourne d'un angle $\theta$
autour de $z$.}
\label{fig:portes}
\end{figure}

\begin{center}
\renewcommand{\arraystretch}{1}
\begin{tabular}{@{}lccc@{}}
\toprule
\textbf{Porte} & \textbf{Matrice} & \textbf{Effet sur la base} & \textbf{Rotation de Bloch} \\
\midrule
$X$ & $\begin{pmatrix}0 & 1\\ 1 & 0\end{pmatrix}$ & $\ket{0}\leftrightarrow\ket{1}$ & $\pi$ autour de $x$ \\[2.2ex]
$Y$ & $\begin{pmatrix}0 & -\ii\\ \ii & 0\end{pmatrix}$ & $\ket{0}\to\ii\ket{1},\ \ket{1}\to-\ii\ket{0}$ & $\pi$ autour de $y$ \\[2.2ex]
$Z$ & $\begin{pmatrix}1 & 0\\ 0 & -1\end{pmatrix}$ & $\ket{1}\to-\ket{1},\ \ket{\pm}\leftrightarrow\ket{\mp}$ & $\pi$ autour de $z$ \\[2.2ex]
$H$ & $\dfrac{1}{\sqrt2}\begin{pmatrix}1 & 1\\ 1 & -1\end{pmatrix}$ & $\ket{0}\leftrightarrow\ket{+},\ \ket{1}\leftrightarrow\ket{-}$ & $\pi$ autour de $\frac{x+z}{\sqrt2}$ \\[2.2ex]
$S$ & $\begin{pmatrix}1 & 0\\ 0 & \ii\end{pmatrix}$ & $\ket{1}\to\ii\ket{1}$ & $\frac\pi2$ autour de $z$ \\[2.2ex]
$T$ & $\begin{pmatrix}1 & 0\\ 0 & \frac{1+\ii}{\sqrt2}\end{pmatrix}$ & $\ket{1}\to\ee^{\ii\pi/4}\ket{1}$ & $\frac\pi4$ autour de $z$ \\
\bottomrule
\end{tabular}
\end{center}

\begin{figure}[!ht]
\centering
\begin{tikzpicture}[x={(-1.63cm,-0.56cm)}, y={(1.63cm,-0.56cm)}, z={(0cm,2.16cm)}, >=Stealth, scale=0.9]
  % sphère et grands cercles
  \shade[ball color=insagris!15, opacity=0.55] (0,0,0) circle[radius=2.3cm];
  \begin{scope}[insagris!75, thin]
    \draw[domain=-45:135, samples=60, variable=\t] plot ({cos(\t)},{sin(\t)},0);
    \draw[domain=135:315, samples=60, variable=\t, dashed] plot ({cos(\t)},{sin(\t)},0);
    \draw[domain=-62.7:117.3, samples=60, variable=\t] plot ({cos(\t)},0,{sin(\t)});
    \draw[domain=117.3:297.3, samples=60, variable=\t, dashed] plot ({cos(\t)},0,{sin(\t)});
    \draw[domain=-62.7:117.3, samples=60, variable=\t] plot (0,{cos(\t)},{sin(\t)});
    \draw[domain=117.3:297.3, samples=60, variable=\t, dashed] plot (0,{cos(\t)},{sin(\t)});
  \end{scope}
  % axes x, y, z (portes X, Y, Z)
  \draw[insagris, dashed] (0,0,0) -- (-1,0,0);
  \draw[insagris, dashed] (0,0,0) -- (0,-1,0);
  \draw[insagris, dashed] (0,0,0) -- (0,0,-1);
  \draw[insagris] (-1,0,0) -- (-1.3,0,0);
  \draw[insagris] (0,-1,0) -- (0,-1.3,0);
  \draw[insagris] (0,0,-1) -- (0,0,-1.3);
  \draw[insagris, -{Stealth[length=2.2mm]}] (0,0,0) -- (1.3,0,0);
  \draw[insagris, -{Stealth[length=2.2mm]}] (0,0,0) -- (0,1.3,0);
  \draw[insagris, -{Stealth[length=2.2mm]}] (0,0,0) -- (0,0,1.3);
  % trajet : |0> --H--> |+> --S--> |+i>
  \draw[insarouge, very thick, -{Stealth[length=2.6mm]}, domain=3:87, samples=30, variable=\t]
    plot ({sin(\t)},0,{cos(\t)});
  \draw[insarouge, very thick, -{Stealth[length=2.6mm]}, domain=3:87, samples=30, variable=\t]
    plot ({cos(\t)},{sin(\t)},0);
  \foreach \ax/\ay/\az in {0/0/1, 1/0/0, 0/1/0} \fill[insarouge] (\ax,\ay,\az) circle[radius=2pt];
  \node[insarouge, anchor=south] at (0,0,1.33) {$\ket{0}$};
  \node[insarouge, anchor=east] at (1.33,0,0) {$\ket{+}$};
  \node[insarouge, anchor=west] at (0,1.33,0) {$\ket{{+}\ii}$};
  \node[insarouge, font=\small, anchor=south west] at ({sin(45)},0,{cos(45)}) {$H$};
  \node[insarouge, font=\small, anchor=north west] at ({cos(45)},{sin(45)},0) {$S$};
\end{tikzpicture}
\caption{Trajet d'un état sous le circuit $H$ puis $S$ (flèches schématiques). $H$ envoie
$\ket{0}$ sur $\ket{+}$ ; $S$, rotation d'un quart de tour autour de $z$, envoie $\ket{+}$ sur
$\ket{{+}\ii}$. Les trois états sont alignés sur les axes $z$, $x$ et $y$.}
\label{fig:trajet}
\end{figure}

% ---------------------------------------------------------------------
\subsection{Circuits quantiques}\label{sec:circuits}
% ---------------------------------------------------------------------

Un \emph{circuit} représente un calcul : chaque qubit est un fil horizontal, le temps va de gauche
à droite, et chaque porte est un symbole posé sur les fils qu'elle concerne. La
figure~\ref{fig:symboles} donne les symboles utilisés dans ce document.

\begin{figure}[!ht]
\centering
\begin{adjustbox}{max width=\linewidth}
\begin{tikzpicture}[>=Stealth, font=\small,
    nom/.style={font=\footnotesize, align=center, anchor=north}]
  % (a) porte à un qubit
  \draw (0,0) -- (1.8,0);
  \node[porte] at (0.9,0) {$U$};
  \node[nom] at (0.9,-1.15) {porte $U$\\ à un qubit};
  % (b) CNOT
  \begin{scope}[xshift=3.2cm]
    \draw (0,0.6) -- (1.8,0.6);
    \draw (0,-0.6) -- (1.8,-0.6);
    \draw[insarouge, thick] (0.9,0.6) -- (0.9,-0.6);
    \pic at (0.9,0.6) {ctl};
    \pic at (0.9,-0.6) {cible};
    \node[nom] at (0.9,-1.15) {CNOT\\ (contrôle $\bullet$, cible $\oplus$)};
  \end{scope}
  % (c) SWAP
  \begin{scope}[xshift=6.4cm]
    \draw (0,0.6) -- (1.8,0.6);
    \draw (0,-0.6) -- (1.8,-0.6);
    \draw[insarouge, thick] (0.9,0.6) -- (0.9,-0.6);
    \pic at (0.9,0.6) {croix};
    \pic at (0.9,-0.6) {croix};
    \node[nom] at (0.9,-1.15) {SWAP\\ (échange)};
  \end{scope}
  % (d) contrôlé-U
  \begin{scope}[xshift=9.6cm]
    \draw (0,0.6) -- (1.8,0.6);
    \draw (0,-0.6) -- (1.8,-0.6);
    \draw[insarouge, thick] (0.9,-0.6) -- (0.9,0.3);
    \pic at (0.9,-0.6) {ctl};
    \node[porte] at (0.9,0.6) {$U$};
    \node[nom] at (0.9,-1.15) {contrôlé-$U$\\ (contrôle en bas)};
  \end{scope}
  % (e) mesure
  \begin{scope}[xshift=12.8cm]
    \draw (0,0) -- (0.6,0);
    \pic at (0.9,0) {mesure};
    \draw (1.2,0.03) -- (2.0,0.03);
    \draw (1.2,-0.03) -- (2.0,-0.03);
    \node[nom] at (1.0,-1.15) {mesure\\ (fil double : bit)};
  \end{scope}
\end{tikzpicture}
\end{adjustbox}
\caption{Symboles des circuits quantiques. Le point noir $\bullet$ marque un qubit de contrôle,
le symbole $\oplus$ une cible qui subit $X$, les deux croix un échange. Après une mesure, le fil
double transporte un bit classique.}
\label{fig:symboles}
\end{figure}

\begin{remarque}[Conventions de lecture]
\begin{itemize}[nosep]
  \item Le temps va de gauche à droite. Deux portes successives $U_1$ puis $U_2$ se multiplient
        dans l'ordre inverse : la matrice du circuit est $U_2U_1$, car $U_1$ agit d'abord sur le ket.
  \item Avec plusieurs fils, le fil du \emph{haut} porte le \emph{dernier} chiffre du ket et le
        fil du \emph{bas} le premier : si $A$ est le premier qubit et $B$ le second,
        $\ket{a,b}=\ket{a}_A\otimes\ket{b}_B$ et $B$ est dessiné au-dessus de $A$.
  \item Deux portes placées dans la même colonne, sur deux fils, agissent en parallèle : c'est un
        produit tensoriel $U_A\otimes U_B$ (section~\ref{sec:op2}).
\end{itemize}
\end{remarque}

\paragraph{Exemple.} Le circuit de la figure~\ref{fig:circuit-hs} applique $H$ puis $S$ à
$\ket{0}$. Sa matrice est $SH$ :
\[
  SH=\begin{pmatrix}1&0\\0&\ii\end{pmatrix}\frac{1}{\sqrt2}\begin{pmatrix}1&1\\1&-1\end{pmatrix}
  =\frac{1}{\sqrt2}\begin{pmatrix}1&1\\ \ii&-\ii\end{pmatrix},
  \qquad
  SH\ket{0}=\frac{1}{\sqrt2}\begin{pmatrix}1\\ \ii\end{pmatrix}=\ket{{+}\ii}.
\]
Le circuit prépare donc l'état $\ket{{+}\ii}$ de la base $Y$ à partir de $\ket{0}$
(figure~\ref{fig:trajet}). L'ordre compte : $HS\ket{0}=H\ket{0}=\ket{+}$, car $S\ket{0}=\ket{0}$.

\begin{figure}[!ht]
\centering
\begin{tikzpicture}[>=Stealth, font=\small]
  \draw (0,0) -- (6.6,0);
  \node[anchor=east] at (0,0) {$\ket{0}$};
  \node[porte] at (2.0,0) {$H$};
  \node[porte] at (4.3,0) {$S$};
  \node[anchor=west] at (6.6,0) {$\ket{{+}\ii}$};
  \foreach \x in {0.7, 3.15, 5.5} \draw[dashed, insagris!70] (\x,0.6) -- (\x,-0.95);
  \node[font=\footnotesize, anchor=north] at (0.7,-0.95) {$\ket{0}$};
  \node[font=\footnotesize, anchor=north] at (3.15,-0.95) {$\ket{+}$};
  \node[font=\footnotesize, anchor=north] at (5.5,-0.95) {$\ket{{+}\ii}$};
\end{tikzpicture}
\caption{Circuit $H$ puis $S$ appliqué à $\ket{0}$. Les états intermédiaires sont indiqués sous
les traits pointillés ; la matrice du circuit est $SH$.}
\label{fig:circuit-hs}
\end{figure}

\paragraph{Mesure.} Le symbole de mesure (compteur) donne le résultat $0$ ou $1$ dans la base $Z$.
Mesurer dans la base $X$ revient à appliquer $H$ puis à mesurer dans la base $Z$
(figure~\ref{fig:mesurex}) : comme $H\ket{\pm}=\ket{0}$ ou $\ket{1}$, le résultat $0$ correspond à
$\ket{+}$ et le résultat $1$ à $\ket{-}$, avec $P(\pm)=|\braket{\pm}{\Psi}|^2$.

\begin{figure}[!ht]
\centering
\begin{tikzpicture}[>=Stealth, font=\small]
  % (a) base Z
  \draw (0,0) -- (1.2,0);
  \node[anchor=east] at (0,0) {$\ket{\Psi}$};
  \pic at (1.5,0) {mesure};
  \draw (1.8,0.03) -- (2.6,0.03);
  \draw (1.8,-0.03) -- (2.6,-0.03);
  \node[anchor=west] at (2.7,0) {$0$ ou $1$};
  \node[font=\footnotesize] at (1.5,-0.95) {(a) mesure en base $Z$};
  % (b) base X
  \begin{scope}[xshift=7.5cm]
    \draw (0,0) -- (2.6,0);
    \node[anchor=east] at (0,0) {$\ket{\Psi}$};
    \node[porte] at (1.0,0) {$H$};
    \pic at (2.3,0) {mesure};
    \draw (2.6,0.03) -- (3.4,0.03);
    \draw (2.6,-0.03) -- (3.4,-0.03);
    \node[anchor=west] at (3.5,0) {$+$ ou $-$};
    \node[font=\footnotesize] at (1.9,-0.95) {(b) mesure en base $X$};
  \end{scope}
\end{tikzpicture}
\caption{(a) Mesure dans la base $Z$. (b) Mesure dans la base $X$ : porte de Hadamard, puis
mesure dans la base $Z$ ; les résultats $0$ et $1$ se lisent $+$ et $-$.}
\label{fig:mesurex}
\end{figure}

% ---------------------------------------------------------------------
\subsection{Portes locales sur deux qubits}\label{sec:op2}
% ---------------------------------------------------------------------

Une porte appliquée à un seul qubit d'un système de deux s'écrit avec un produit tensoriel
(section~\ref{sec:tenseur}) : $U\otimes I$ agit sur $A$ (premier facteur, fil du bas), $I\otimes U$
agit sur $B$, et $U_A\otimes U_B$ agit sur les deux en parallèle (figure~\ref{fig:local}).

\begin{figure}[!ht]
\centering
\begin{tikzpicture}[>=Stealth, font=\small]
  \foreach \xs/\haut/\bas/\tit in {0/{}/{$X$}/{(a) $X\otimes I$},
                                   5.4/{$X$}/{}/{(b) $I\otimes X$},
                                   10.8/{$H$}/{$H$}/{(c) $H\otimes H$}} {
    \begin{scope}[xshift=\xs cm]
      \draw (0,0.75) -- (2.6,0.75);
      \draw (0,0) -- (2.6,0);
      \node[anchor=east] at (0,0.75) {$B$};
      \node[anchor=east] at (0,0) {$A$};
      \ifx\haut\empty\else \node[porte] at (1.3,0.75) {\haut}; \fi
      \ifx\bas\empty\else \node[porte] at (1.3,0) {\bas}; \fi
      \node[font=\footnotesize] at (1.3,-0.85) {\tit};
    \end{scope}
  }
\end{tikzpicture}
\caption{Portes en parallèle sur deux qubits. Le qubit $A$, premier facteur du produit tensoriel,
est sur le fil du bas. (a) $X$ sur $A$ seulement. (b) $X$ sur $B$ seulement. (c) $H$ sur les deux.}
\label{fig:local}
\end{figure}

\subsubsection{Portes agissant sur un seul des deux qubits}

On note $I=\mathrm{Id}$ ($2\times2$). $X\otimes I$ agit sur $A$, $I\otimes X$ agit sur $B$ :
\begin{align*}
  I\otimes X&=\begin{pmatrix}1\cdot X & 0\cdot X\\ 0\cdot X & 1\cdot X\end{pmatrix}
    =\begin{pmatrix}0 & 1 & 0 & 0\\ 1 & 0 & 0 & 0\\ 0 & 0 & 0 & 1\\ 0 & 0 & 1 & 0\end{pmatrix},\\[1ex]
  X\otimes I&=\begin{pmatrix}0\cdot I & 1\cdot I\\ 1\cdot I & 0\cdot I\end{pmatrix}
    =\begin{pmatrix}0 & 0 & 1 & 0\\ 0 & 0 & 0 & 1\\ 1 & 0 & 0 & 0\\ 0 & 1 & 0 & 0\end{pmatrix}.
\end{align*}
Sur $\ket{00}=(1,0,0,0)^{\mathsf T}$ (première colonne de la matrice) :
$(I\otimes X)\ket{00}=(0,1,0,0)^{\mathsf T}=\ket{01}$ et
$(X\otimes I)\ket{00}=(0,0,1,0)^{\mathsf T}=\ket{10}$. Avec la règle :
$(I\otimes X)\ket{00}=I\ket{0}\otimes X\ket{0}=\ket{0}\otimes\ket{1}=\ket{01}$.

Pour la porte de Hadamard, avec $H=\frac{1}{\sqrt2}\begin{pmatrix}1 & 1\\ 1 & -1\end{pmatrix}$ :
\[
  H\otimes I=\frac{1}{\sqrt2}\begin{pmatrix}1\cdot I & 1\cdot I\\ 1\cdot I & -1\cdot I\end{pmatrix}
  =\frac{1}{\sqrt2}\begin{pmatrix}1 & 0 & 1 & 0\\ 0 & 1 & 0 & 1\\ 1 & 0 & -1 & 0\\ 0 & 1 & 0 & -1\end{pmatrix}.
\]
Ses colonnes donnent l'image des quatre états de base :
\begin{align*}
  (H\otimes I)\ket{00}&=\tfrac{1}{\sqrt2}(1,0,1,0)^{\mathsf T}=\ket{+}\otimes\ket{0},&
  (H\otimes I)\ket{01}&=\tfrac{1}{\sqrt2}(0,1,0,1)^{\mathsf T}=\ket{+}\otimes\ket{1},\\
  (H\otimes I)\ket{10}&=\tfrac{1}{\sqrt2}(1,0,-1,0)^{\mathsf T}=\ket{-}\otimes\ket{0},&
  (H\otimes I)\ket{11}&=\tfrac{1}{\sqrt2}(0,1,0,-1)^{\mathsf T}=\ket{-}\otimes\ket{1}.
\end{align*}

\subsubsection{Hadamard sur les deux qubits}

\[
  H\otimes H=\frac12\begin{pmatrix}1\cdot H' & 1\cdot H'\\ 1\cdot H' & -1\cdot H'\end{pmatrix},
  \quad H'=\begin{pmatrix}1 & 1\\ 1 & -1\end{pmatrix},
  \qquad
  H\otimes H=\frac12\begin{pmatrix}1 & 1 & 1 & 1\\ 1 & -1 & 1 & -1\\ 1 & 1 & -1 & -1\\ 1 & -1 & -1 & 1\end{pmatrix}.
\]
La première colonne donne $(H\otimes H)\ket{00}=\frac12(1,1,1,1)^{\mathsf T}=\ket{+}\otimes\ket{+}$
(cohérent avec $H\ket{0}\otimes H\ket{0}$) : les quatre résultats $00,01,10,11$ sont
équiprobables, chacun avec la probabilité $\frac14$. Sur $n$ qubits, $H^{\otimes n}\ket{0\ldots0}$
donne une superposition uniforme de $2^n$ états.

% ---------------------------------------------------------------------
\subsection{Porte CNOT}\label{sec:cnot}
% ---------------------------------------------------------------------

\begin{definition}{Porte CNOT (NOT contrôlé)}
La porte CNOT agit sur deux qubits : le qubit $B$ est le \emph{contrôle}, le qubit $A$ est la
\emph{cible}. La cible est inversée (porte $X$) si le contrôle vaut $1$ et reste inchangée
s'il vaut $0$. Avec $\ket{a,b}=\ket{a}_A\otimes\ket{b}_B$ :
\[
  \mathrm{CNOT}\,\ket{a,b}=\ket{a\oplus b,\,b},
  \qquad
  \mathrm{CNOT}=I\otimes\ket{0}\bra{0}+X\otimes\ket{1}\bra{1},
\]
où $\oplus$ est l'addition modulo $2$ ($0\oplus0=1\oplus1=0$ et $0\oplus1=1\oplus0=1$).
\end{definition}

Sur le circuit de la figure~\ref{fig:cnot}, le point noir posé sur le fil de $B$ (en haut) marque
le contrôle et le symbole $\oplus$ posé sur le fil de $A$ (en bas) marque la cible.

\begin{figure}[!ht]
\centering
\begin{tikzpicture}[>=Stealth, font=\small]
  \draw (0,1.2) -- (6,1.2);
  \draw (0,0) -- (6,0);
  \node[anchor=east] at (0,1.2) {$\ket{b}_B$};
  \node[anchor=east] at (0,0) {$\ket{a}_A$};
  \node[anchor=west] at (6,1.2) {$\ket{b}_B$};
  \node[anchor=west] at (6,0) {$\ket{a\oplus b}_A$};
  \draw[insarouge, thick] (3,1.2) -- (3,0);
  \pic at (3,1.2) {ctl};
  \pic at (3,0) {cible};
  \node[anchor=south, text=insagris] at (3,1.5) {qubit de contrôle};
  \node[anchor=north, text=insagris] at (3,-0.5) {qubit cible};
\end{tikzpicture}
\caption{Porte CNOT : le fil du haut ($B$) est le contrôle, le fil du bas ($A$) est la cible. Le
ket s'écrit $\ket{a,b}$ : le fil du haut donne le dernier chiffre.}
\label{fig:cnot}
\end{figure}

\subsubsection{Table de vérité}

\begin{center}
\renewcommand{\arraystretch}{1.3}
\begin{tabular}{@{}cccc@{}}
\toprule
Entrée $\ket{a,b}$ & Contrôle $b$ & Sortie $\ket{a\oplus b,b}$ & Effet sur la cible $A$ \\
\midrule
$\ket{00}$ & $0$ & $\ket{00}$ & aucun \\
$\ket{01}$ & $1$ & $\ket{11}$ & $0\to1$ \\
$\ket{10}$ & $0$ & $\ket{10}$ & aucun \\
$\ket{11}$ & $1$ & $\ket{01}$ & $1\to0$ \\
\bottomrule
\end{tabular}
\end{center}

\subsubsection{Matrice du CNOT}

La $j$-ième colonne d'une matrice $4\times4$ est l'image du $j$-ième vecteur de base
($\ket{00},\ket{01},\ket{10},\ket{11}$, dans cet ordre). La table donne :
\begin{align*}
  \mathrm{CNOT}\ket{00}&=\ket{00}=\begin{pmatrix}1\\0\\0\\0\end{pmatrix},&
  \mathrm{CNOT}\ket{01}&=\ket{11}=\begin{pmatrix}0\\0\\0\\1\end{pmatrix},\\[1.5ex]
  \mathrm{CNOT}\ket{10}&=\ket{10}=\begin{pmatrix}0\\0\\1\\0\end{pmatrix},&
  \mathrm{CNOT}\ket{11}&=\ket{01}=\begin{pmatrix}0\\1\\0\\0\end{pmatrix},
\end{align*}
d'où, en juxtaposant ces quatre colonnes,
\[
  \mathrm{CNOT}=\begin{pmatrix}
    1 & 0 & 0 & 0\\ 0 & 0 & 0 & 1\\ 0 & 0 & 1 & 0\\ 0 & 1 & 0 & 0
  \end{pmatrix}.
\]

On retrouve cette matrice avec la formule $I\otimes\ket{0}\bra{0}+X\otimes\ket{1}\bra{1}$ et
la règle $A\otimes B=(A_{ij}B)$ (section~\ref{sec:tenseur}), où
$\ket{0}\bra{0}=\begin{pmatrix}1&0\\0&0\end{pmatrix}$ et
$\ket{1}\bra{1}=\begin{pmatrix}0&0\\0&1\end{pmatrix}$ :
\[
  I\otimes\ket{0}\bra{0}=\begin{pmatrix}1&0&0&0\\0&0&0&0\\0&0&1&0\\0&0&0&0\end{pmatrix},
  \qquad
  X\otimes\ket{1}\bra{1}=\begin{pmatrix}0&0&0&0\\0&0&0&1\\0&0&0&0\\0&1&0&0\end{pmatrix},
\]
et la somme de ces deux matrices est la matrice du CNOT.

\subsubsection{Action sur un état quelconque}

Pour $\ket{\Psi}=c_{00}\ket{00}+c_{01}\ket{01}+c_{10}\ket{10}+c_{11}\ket{11}$ :
\[
  \mathrm{CNOT}\begin{pmatrix}c_{00}\\c_{01}\\c_{10}\\c_{11}\end{pmatrix}
  =\begin{pmatrix}
    1\cdot c_{00}+0\cdot c_{01}+0\cdot c_{10}+0\cdot c_{11}\\
    0\cdot c_{00}+0\cdot c_{01}+0\cdot c_{10}+1\cdot c_{11}\\
    0\cdot c_{00}+0\cdot c_{01}+1\cdot c_{10}+0\cdot c_{11}\\
    0\cdot c_{00}+1\cdot c_{01}+0\cdot c_{10}+0\cdot c_{11}
  \end{pmatrix}
  =\begin{pmatrix}c_{00}\\c_{11}\\c_{10}\\c_{01}\end{pmatrix}.
\]
Le CNOT échange les amplitudes de $\ket{01}$ et de $\ket{11}$, c'est-à-dire des deux états
dont le contrôle vaut $1$, et laisse les autres.

\paragraph{Exemple : contrôle en superposition.} Soit $\ket{0}_A\otimes\ket{+}_B
=\frac{1}{\sqrt2}\big(\ket{00}+\ket{01}\big)=\frac{1}{\sqrt2}(1,1,0,0)^{\mathsf T}$,
un état séparable. Alors
\[
  \mathrm{CNOT}\,\frac{1}{\sqrt2}\begin{pmatrix}1\\1\\0\\0\end{pmatrix}
  =\frac{1}{\sqrt2}\begin{pmatrix}1\\0\\0\\1\end{pmatrix}
  =\frac{1}{\sqrt2}\big(\ket{00}+\ket{11}\big).
\]
Le résultat est intriqué ($c_{00}c_{11}-c_{01}c_{10}=\frac12\neq0$ ; section~\ref{sec:separable}) :
un CNOT dont le contrôle est en superposition crée de l'intrication. C'est l'état de Bell
$\ket{\Phi^+}$ du chapitre~\ref{chap:bell}.

\paragraph{Réversibilité.} Deux CNOT successifs redonnent l'état initial :
$\ket{a,b}\to\ket{a\oplus b,b}\to\ket{a\oplus b\oplus b,b}=\ket{a,b}$, donc
$\mathrm{CNOT}^2=\mathrm{Id}$. La matrice est réelle et symétrique, donc
$\mathrm{CNOT}^\dagger\mathrm{CNOT}=\mathrm{CNOT}^2=\mathrm{Id}$ : elle est bien unitaire.

\begin{remarque}[CNOT et copie d'un bit]
Avec la cible dans l'état $\ket{0}$, $\mathrm{CNOT}\ket{0,b}=\ket{b,b}$ pour $b=0$ et $b=1$ : le CNOT
copie le bit classique $b$. En revanche $\mathrm{CNOT}\big(\ket{0}\otimes\ket{+}\big)$ est l'état
ci-dessus, et non $\ket{+}\otimes\ket{+}$ : un état quantique inconnu ne peut pas être copié
(théorème de non-clonage, section~\ref{sec:teleportation}).
\end{remarque}

\begin{remarque}[Contrôle en dernier ou en premier]
La matrice dépend de l'ordre d'écriture du ket. Notons $\mathrm{CNOT}_{21}$ le CNOT dont le
contrôle est le \emph{second} chiffre du ket : c'est la matrice ci-dessus, adoptée dans ce
document (fil du haut). Beaucoup d'ouvrages mettent le contrôle en \emph{premier} ; le CNOT
correspondant est
\[
  \mathrm{CNOT}_{12}=\ket{0}\bra{0}\otimes I+\ket{1}\bra{1}\otimes X
  =\begin{pmatrix}1&0&0&0\\0&1&0&0\\0&0&0&1\\0&0&1&0\end{pmatrix}.
\]
Le même circuit a donc deux matrices, selon que le contrôle est écrit en dernier ou en premier.
Les deux sont reliées par le SWAP (section~\ref{sec:swap}). La section~\ref{sec:controlees} sur les
portes contrôlées met le contrôle en premier, avec la matrice $\mathrm{diag}(I,U)$.
\end{remarque}

% ---------------------------------------------------------------------
\subsection{Porte SWAP}\label{sec:swap}
% ---------------------------------------------------------------------

\begin{definition}{Porte SWAP (échange)}
La porte SWAP échange les états des deux qubits :
\[
  \mathrm{SWAP}\big(\ket{\psi}\otimes\ket{\phi}\big)=\ket{\phi}\otimes\ket{\psi}.
\]
Sur les états de base, $\mathrm{SWAP}\ket{a,b}=\ket{b,a}$ ; on prolonge par linéarité.
\end{definition}

Son symbole est formé de deux croix reliées par un trait vertical
(figure~\ref{fig:swap}).

\begin{figure}[!ht]
\centering
\begin{tikzpicture}[>=Stealth, font=\small]
  \draw (0,1.2) -- (6,1.2);
  \draw (0,0) -- (6,0);
  \node[anchor=east] at (0,1.2) {$\ket{a}$};
  \node[anchor=east] at (0,0) {$\ket{b}$};
  \node[anchor=west] at (6,1.2) {$\ket{b}$};
  \node[anchor=west] at (6,0) {$\ket{a}$};
  \draw[insarouge, thick] (3,1.2) -- (3,0);
  \pic at (3,1.2) {croix};
  \pic at (3,0) {croix};
\end{tikzpicture}
\caption{Porte SWAP : les deux qubits échangent leurs états.}
\label{fig:swap}
\end{figure}

\subsubsection{Table et matrice}

\begin{center}
\renewcommand{\arraystretch}{1.3}
\begin{tabular}{@{}cc@{}}
\toprule
Entrée & Sortie \\
\midrule
$\ket{00}$ & $\ket{00}$ \\
$\ket{01}$ & $\ket{10}$ \\
$\ket{10}$ & $\ket{01}$ \\
$\ket{11}$ & $\ket{11}$ \\
\bottomrule
\end{tabular}
\end{center}

Les images des quatre vecteurs de base sont les colonnes de la matrice :
\[
  \mathrm{SWAP}\ket{00}=\begin{pmatrix}1\\0\\0\\0\end{pmatrix},\quad
  \mathrm{SWAP}\ket{01}=\begin{pmatrix}0\\0\\1\\0\end{pmatrix},\quad
  \mathrm{SWAP}\ket{10}=\begin{pmatrix}0\\1\\0\\0\end{pmatrix},\quad
  \mathrm{SWAP}\ket{11}=\begin{pmatrix}0\\0\\0\\1\end{pmatrix},
\]
\[
  \mathrm{SWAP}=\begin{pmatrix}
    1 & 0 & 0 & 0\\ 0 & 0 & 1 & 0\\ 0 & 1 & 0 & 0\\ 0 & 0 & 0 & 1
  \end{pmatrix}.
\]
Appliquée à $\ket{01}=(0,1,0,0)^{\mathsf T}$, cette matrice sélectionne sa deuxième colonne :
\[
  \begin{pmatrix}
    1 & 0 & 0 & 0\\ 0 & 0 & 1 & 0\\ 0 & 1 & 0 & 0\\ 0 & 0 & 0 & 1
  \end{pmatrix}
  \begin{pmatrix}0\\1\\0\\0\end{pmatrix}
  =\begin{pmatrix}0\\0\\1\\0\end{pmatrix}=\ket{10}.
\]
Sur un état quelconque, $\mathrm{SWAP}\,(c_{00},c_{01},c_{10},c_{11})^{\mathsf T}
=(c_{00},c_{10},c_{01},c_{11})^{\mathsf T}$ : les amplitudes de $\ket{01}$ et $\ket{10}$
sont échangées.

\subsubsection{Accord avec la définition}

Pour deux états $\ket{\psi}=(a_0,a_1)^{\mathsf T}$ et $\ket{\phi}=(b_0,b_1)^{\mathsf T}$ :
\[
  \mathrm{SWAP}\,(\ket{\psi}\otimes\ket{\phi})
  =\mathrm{SWAP}\begin{pmatrix}a_0b_0\\a_0b_1\\a_1b_0\\a_1b_1\end{pmatrix}
  =\begin{pmatrix}a_0b_0\\a_1b_0\\a_0b_1\\a_1b_1\end{pmatrix}
  =\begin{pmatrix}b_0\\b_1\end{pmatrix}\otimes\begin{pmatrix}a_0\\a_1\end{pmatrix}
  =\ket{\phi}\otimes\ket{\psi},
\]
ce qui est bien la définition. La porte s'écrit aussi avec des projecteurs :
\[
  \mathrm{SWAP}=\sum_{a,b\in\{0,1\}}\ket{a}\bra{b}\otimes\ket{b}\bra{a}.
\]
En effet, pour $\ket{c}\otimes\ket{d}$ on trouve
$\sum_{a,b}\ket{a}\braket{b}{c}\otimes\ket{b}\braket{a}{d}$, et seul le terme $b=c$, $a=d$
est non nul : le résultat est $\ket{d}\otimes\ket{c}$.

\paragraph{Exemples.}
\begin{itemize}[nosep]
  \item Deux états différents : $\mathrm{SWAP}\big(\ket{0}\otimes\ket{+}\big)
        =\mathrm{SWAP}\,\frac{1}{\sqrt2}(1,1,0,0)^{\mathsf T}=\frac{1}{\sqrt2}(1,0,1,0)^{\mathsf T}
        =\ket{+}\otimes\ket{0}$.
  \item L'état $\ket{\psi}=\frac{1}{\sqrt2}\ket{00}+\frac{1}{\sqrt2}\ket{11}$ ne change pas :
        $\mathrm{SWAP}\ket{\psi}=\frac{1}{\sqrt2}\ket{00}+\frac{1}{\sqrt2}\ket{11}=\ket{\psi}$,
        car $\ket{00}$ et $\ket{11}$ sont invariants. De même
        $\frac{1}{\sqrt2}(\ket{00}-\ket{11})$ et $\frac{1}{\sqrt2}(\ket{01}+\ket{10})$ sont invariants,
        alors que $\mathrm{SWAP}\,\frac{1}{\sqrt2}(\ket{01}-\ket{10})=\frac{1}{\sqrt2}(\ket{10}-\ket{01})$
        est l'opposé de l'état de départ. Ce sont les quatre états de Bell (chapitre~\ref{chap:bell}).
\end{itemize}

\paragraph{Propriétés.} Échanger deux fois redonne l'état de départ :
$\mathrm{SWAP}^2=\mathrm{Id}$ ; la matrice est réelle et symétrique, donc
$\mathrm{SWAP}^\dagger\mathrm{SWAP}=\mathrm{SWAP}^2=\mathrm{Id}$ : elle est unitaire.

\subsubsection{Lien avec le CNOT}

\begin{resultat}{SWAP échange les rôles de contrôle et de cible}
\[
  \mathrm{CNOT}_{21}=\mathrm{SWAP}\;\mathrm{CNOT}_{12}\;\mathrm{SWAP}.
\]
\end{resultat}

\emph{Sur les états de base.}
$\mathrm{SWAP}\,\mathrm{CNOT}_{12}\,\mathrm{SWAP}\ket{a,b}
=\mathrm{SWAP}\,\mathrm{CNOT}_{12}\ket{b,a}
=\mathrm{SWAP}\ket{b,\,b\oplus a}=\ket{b\oplus a,\,b}=\mathrm{CNOT}_{21}\ket{a,b}$.

\emph{Par les matrices.} Multiplier à gauche par $\mathrm{SWAP}$ échange les lignes $2$ et
$3$, multiplier à droite l'échange les colonnes $2$ et $3$ :
\[
  \mathrm{SWAP}\,\mathrm{CNOT}_{12}
  =\begin{pmatrix}1&0&0&0\\0&0&0&1\\0&1&0&0\\0&0&1&0\end{pmatrix},
  \qquad
  \mathrm{SWAP}\,\mathrm{CNOT}_{12}\,\mathrm{SWAP}
  =\begin{pmatrix}1&0&0&0\\0&0&0&1\\0&0&1&0\\0&1&0&0\end{pmatrix}=\mathrm{CNOT}_{21}.
\]

\begin{remarque}[SWAP avec trois CNOT]
$\mathrm{SWAP}=\mathrm{CNOT}_{12}\,\mathrm{CNOT}_{21}\,\mathrm{CNOT}_{12}$. En effet, sur
$\ket{a,b}$ (on applique les portes de droite à gauche) :
\[
  \ket{a,b}\ \xrightarrow{\ \mathrm{CNOT}_{12}\ }\ \ket{a,\,a\oplus b}
  \ \xrightarrow{\ \mathrm{CNOT}_{21}\ }\ \ket{a\oplus(a\oplus b),\,a\oplus b}=\ket{b,\,a\oplus b}
  \ \xrightarrow{\ \mathrm{CNOT}_{12}\ }\ \ket{b,\,b\oplus(a\oplus b)}=\ket{b,a}.
\]
Les deux membres coïncident sur les états de base, donc sur tous les états par linéarité.
\end{remarque}

% ---------------------------------------------------------------------
\subsection{Portes contrôlées : contrôlé-\texorpdfstring{$U$}{U}, Fredkin et Toffoli}\label{sec:controlees}
% ---------------------------------------------------------------------

Le CNOT est le cas le plus simple d'une idée générale : appliquer une opération à un ou
plusieurs qubits \emph{cibles} seulement si un ou plusieurs qubits de \emph{contrôle} valent $1$.
Dans cette section, le contrôle est le \emph{premier} chiffre
du ket (remarque de la section~\ref{sec:cnot}). Les qubits sont notés $x$, $y$, $z$ (le qubit $x$ est le contrôle) et
$\ket{x,y,z}=\ket{x}\otimes\ket{y}\otimes\ket{z}$ ; d'après la convention de la section~\ref{sec:circuits}, le
contrôle est dessiné en bas (figure~\ref{fig:controlees}). Pour éviter la confusion avec les
matrices de Pauli, on écrit ici $\sigma_x=X$ pour la porte NOT.

\begin{figure}[!ht]
\centering
\begin{adjustbox}{max width=\linewidth}
\begin{tikzpicture}[>=Stealth, font=\small]
  % (a) contrôlé-U : contrôle x en bas, cible y en haut
  \draw (0,1.2) -- (3.4,1.2);
  \draw (0,0) -- (3.4,0);
  \node[anchor=east] at (0,1.2) {$y$};
  \node[anchor=east] at (0,0) {$x$};
  \draw[insarouge, thick] (1.7,0) -- (1.7,0.85);
  \pic at (1.7,0) {ctl};
  \node[porte] at (1.7,1.2) {$U$};
  \node[anchor=north, font=\footnotesize] at (1.7,-0.5) {(a) contrôlé-$U$};
  % (b) contrôlé-SWAP (Fredkin) : x contrôle, y et z échangés
  \draw (5,2.4) -- (8.4,2.4);
  \draw (5,1.2) -- (8.4,1.2);
  \draw (5,0) -- (8.4,0);
  \node[anchor=east] at (5,2.4) {$z$};
  \node[anchor=east] at (5,1.2) {$y$};
  \node[anchor=east] at (5,0) {$x$};
  \draw[insarouge, thick] (6.7,0) -- (6.7,2.4);
  \pic at (6.7,0) {ctl};
  \pic at (6.7,1.2) {croix};
  \pic at (6.7,2.4) {croix};
  \node[anchor=north, font=\footnotesize] at (6.7,-0.5) {(b) Fredkin (C-SWAP)};
  % (c) Toffoli (CCNOT) : x et y contrôlent, z cible
  \draw (10,2.4) -- (13.4,2.4);
  \draw (10,1.2) -- (13.4,1.2);
  \draw (10,0) -- (13.4,0);
  \node[anchor=east] at (10,2.4) {$z$};
  \node[anchor=east] at (10,1.2) {$y$};
  \node[anchor=east] at (10,0) {$x$};
  \draw[insarouge, thick] (11.7,0) -- (11.7,2.4);
  \pic at (11.7,0) {ctl};
  \pic at (11.7,1.2) {ctl};
  \pic at (11.7,2.4) {cible};
  \node[anchor=north, font=\footnotesize] at (11.7,-0.5) {(c) Toffoli (CCNOT)};
\end{tikzpicture}
\end{adjustbox}
\caption{Portes contrôlées. Le contrôle $x$ est le premier chiffre du ket, donc le fil du bas.
(a) Le qubit $y$ subit $U$ si $x=1$ ; pour $U=X$ on retrouve le CNOT. (b) Si $x=1$, les qubits
$y$ et $z$ sont échangés. (c) Si $x=y=1$, le qubit $z$ est inversé.}
\label{fig:controlees}
\end{figure}

\subsubsection{Contrôlé-\texorpdfstring{$U$}{U}}

\begin{definition}{Porte contrôlée-$U$}
Soit $U$ une matrice unitaire $2\times2$. Sur le système $(x,y)$, de contrôle $x$ et de cible $y$,
\[
  \mathrm{C}U=\ket{0}\bra{0}_x\otimes I_y+\ket{1}\bra{1}_x\otimes U_y .
\]
Le premier terme agit quand le contrôle est $\ket{0}$ : la cible ne change pas ($I$). Le second
agit quand le contrôle est $\ket{1}$ : la cible subit $U$.
\end{definition}

Comme $\ket{0}\bra{0}\ket{x}=\delta_{x0}\ket{0}$ et $\ket{1}\bra{1}\ket{x}=\delta_{x1}\ket{1}$,
\[
  \mathrm{C}U\big(\ket{x}\otimes\ket{y}\big)=\ket{x}\otimes U^x\ket{y},\qquad U^0=I,\ U^1=U .
\]
Par la règle $A\otimes B=(A_{ij}B)$, $\ket{0}\bra{0}\otimes I=\begin{pmatrix}I&0\\0&0\end{pmatrix}$
et $\ket{1}\bra{1}\otimes U=\begin{pmatrix}0&0\\0&U\end{pmatrix}$ (blocs $2\times2$), donc
\[
  \mathrm{C}U=\begin{pmatrix}I&0\\0&U\end{pmatrix}
  =\begin{pmatrix}1&0&0&0\\0&1&0&0\\0&0&u_{00}&u_{01}\\0&0&u_{10}&u_{11}\end{pmatrix}
  \qquad\text{pour }U=\begin{pmatrix}u_{00}&u_{01}\\u_{10}&u_{11}\end{pmatrix}.
\]
Cette matrice est unitaire car $I^\dagger I=I$ et $U^\dagger U=I$.

\paragraph{Le CNOT est un cas particulier.} Pour $U=\sigma_x$,
$\mathrm{C}\sigma_x=\ket{0}\bra{0}_x\otimes I_y+\ket{1}\bra{1}_x\otimes\sigma_x$ est le CNOT de
contrôle $x$ (premier chiffre), c'est-à-dire $\mathrm{CNOT}_{12}$ de la section~\ref{sec:cnot}.
On calcule l'action sur les quatre états de base, avec $\braket{0}{0}=\braket{1}{1}=1$,
$\braket{0}{1}=\braket{1}{0}=0$, $\sigma_x\ket{0}=\ket{1}$ et $\sigma_x\ket{1}=\ket{0}$ :
\begin{align*}
  \mathrm{C}\sigma_x\ket{00}&=\big[\ket{0}\bra{0}_x\otimes I_y\big]\ket{0}_x\ket{0}_y
      +\big[\ket{1}\bra{1}_x\otimes\sigma_x\big]\ket{0}_x\ket{0}_y\\
    &=\ket{0}\braket{0}{0}\otimes I\ket{0}+\ket{1}\braket{1}{0}\otimes\sigma_x\ket{0}
     =\ket{0}\otimes\ket{0}+0=\ket{00},\\[1.2ex]
  \mathrm{C}\sigma_x\ket{01}&=\ket{0}\braket{0}{0}\otimes I\ket{1}+\ket{1}\braket{1}{0}\otimes\sigma_x\ket{1}
     =\ket{0}\otimes\ket{1}+0=\ket{01},\\[1.2ex]
  \mathrm{C}\sigma_x\ket{10}&=\ket{0}\braket{0}{1}\otimes I\ket{0}+\ket{1}\braket{1}{1}\otimes\sigma_x\ket{0}
     =0+\ket{1}\otimes\ket{1}=\ket{11},\\[1.2ex]
  \mathrm{C}\sigma_x\ket{11}&=\ket{0}\braket{0}{1}\otimes I\ket{1}+\ket{1}\braket{1}{1}\otimes\sigma_x\ket{1}
     =0+\ket{1}\otimes\ket{0}=\ket{10}.
\end{align*}
On retrouve la table du CNOT : la cible n'est inversée que si le contrôle vaut $1$.

\paragraph{Un autre exemple : $U=Z$.} Ici
\[
  \mathrm{C}Z=\begin{pmatrix}I&0\\0&Z\end{pmatrix}=\mathrm{diag}(1,1,1,-1) :
\]
la porte multiplie $\ket{11}$ par $-1$ et ne touche pas les autres états de base. Elle est
symétrique entre les deux qubits. Elle se déduit du CNOT par
$\mathrm{C}Z=(I\otimes H)\,\mathrm{C}X\,(I\otimes H)$, où $\mathrm{C}X$ est le contrôlé-$X$.
En effet, $I\otimes H=\mathrm{diag}(H,H)$ (deux blocs $2\times2$ égaux à $H$), et le produit de
matrices diagonales par blocs se fait bloc par bloc :
\[
  (I\otimes H)\begin{pmatrix}I&0\\0&X\end{pmatrix}(I\otimes H)
  =\begin{pmatrix}H\,I\,H&0\\0&H\,X\,H\end{pmatrix}
  =\begin{pmatrix}I&0\\0&HXH\end{pmatrix},
\]
car $H^2=I$, et
\[
  HXH=\frac12\begin{pmatrix}1&1\\1&-1\end{pmatrix}\begin{pmatrix}0&1\\1&0\end{pmatrix}\begin{pmatrix}1&1\\1&-1\end{pmatrix}
  =\frac12\begin{pmatrix}1&1\\1&-1\end{pmatrix}\begin{pmatrix}1&-1\\1&1\end{pmatrix}
  =\frac12\begin{pmatrix}2&0\\0&-2\end{pmatrix}=Z
\]
(section~\ref{sec:hadamard} : $H$ échange les axes $x$ et $z$).

\subsubsection{Contrôlé-SWAP : la porte de Fredkin}

\begin{definition}{Porte de Fredkin (C-SWAP)}
Sur trois qubits $(x,y,z)$ (espace de dimension $2^3=8$), avec $x$ pour contrôle,
\[
  \mathrm{C\text{-}SWAP}=\ket{0}\bra{0}_x\otimes I_y\otimes I_z+\ket{1}\bra{1}_x\otimes\mathrm{SWAP}_{yz}.
\]
Si $x=0$, rien ne change ; si $x=1$, les qubits $y$ et $z$ sont échangés.
\end{definition}

Calcul sur quatre états de base (avec $I\ket{ab}=\ket{ab}$ et $\mathrm{SWAP}\ket{ab}=\ket{ba}$) :
\begin{align*}
  \mathrm{C\text{-}SWAP}\ket{000}&=\ket{0}\braket{0}{0}\otimes I\ket{00}+\ket{1}\braket{1}{0}\otimes\mathrm{SWAP}\ket{00}
     =\ket{0}\otimes\ket{00}+0=\ket{000},\\[1.2ex]
  \mathrm{C\text{-}SWAP}\ket{101}&=\ket{0}\braket{0}{1}\otimes I\ket{01}+\ket{1}\braket{1}{1}\otimes\mathrm{SWAP}\ket{01}
     =0+\ket{1}\otimes\ket{10}=\ket{110},\\[1.2ex]
  \mathrm{C\text{-}SWAP}\ket{110}&=\ket{0}\braket{0}{1}\otimes I\ket{10}+\ket{1}\braket{1}{1}\otimes\mathrm{SWAP}\ket{10}
     =0+\ket{1}\otimes\ket{01}=\ket{101},\\[1.2ex]
  \mathrm{C\text{-}SWAP}\ket{010}&=\ket{0}\braket{0}{0}\otimes I\ket{10}+\ket{1}\braket{1}{0}\otimes\mathrm{SWAP}\ket{10}
     =\ket{0}\otimes\ket{10}+0=\ket{010}.
\end{align*}
Les huit états de base donnent :
\begin{center}
\renewcommand{\arraystretch}{1.3}
\begin{tabular}{@{}ccccccccc@{}}
\toprule
Entrée & $\ket{000}$ & $\ket{001}$ & $\ket{010}$ & $\ket{011}$ & $\ket{100}$ & $\ket{101}$ & $\ket{110}$ & $\ket{111}$ \\
\midrule
Sortie & $\ket{000}$ & $\ket{001}$ & $\ket{010}$ & $\ket{011}$ & $\ket{100}$ & $\ket{110}$ & $\ket{101}$ & $\ket{111}$ \\
\bottomrule
\end{tabular}
\end{center}
Seuls $\ket{101}$ et $\ket{110}$ changent (contrôle à $1$ et $y\neq z$) ; pour $\ket{100}$ et
$\ket{111}$ l'échange de deux qubits égaux ne change rien. Dans la base
$\ket{000},\ket{001},\dots,\ket{111}$ (ordre binaire), la matrice est
$\mathrm{diag}(I_4,\mathrm{SWAP})$ (deux blocs $4\times4$) :
\[
  \mathrm{C\text{-}SWAP}=\begingroup\setlength{\arraycolsep}{3pt}\begin{pmatrix}
    1&0&0&0&0&0&0&0\\
    0&1&0&0&0&0&0&0\\
    0&0&1&0&0&0&0&0\\
    0&0&0&1&0&0&0&0\\
    0&0&0&0&1&0&0&0\\
    0&0&0&0&0&0&1&0\\
    0&0&0&0&0&1&0&0\\
    0&0&0&0&0&0&0&1
  \end{pmatrix}\endgroup .
\]
C'est l'identité dont les lignes $6$ et $7$ sont échangées.

\paragraph{Contrôle en superposition.} Avec $\ket{+}_x\otimes\ket{0}_y\otimes\ket{1}_z
=\frac{1}{\sqrt2}\big(\ket{001}+\ket{101}\big)$ :
\[
  \mathrm{C\text{-}SWAP}\,\frac{1}{\sqrt2}\big(\ket{001}+\ket{101}\big)
  =\frac{1}{\sqrt2}\big(\ket{001}+\ket{110}\big).
\]
Le résultat n'est pas un produit : les coefficients de $\ket{x}\otimes\ket{yz}$ forment la
matrice (lignes $x=0,1$, colonnes $yz=00,01,10,11$)
\[
  \frac{1}{\sqrt2}\begin{pmatrix}0&1&0&0\\0&0&1&0\end{pmatrix},
\]
dont les deux lignes ne sont pas proportionnelles. Le qubit $x$ est intriqué avec la paire
$(y,z)$.

\subsubsection{Contrôlé-contrôlé-NOT : la porte de Toffoli}

\begin{definition}{Porte de Toffoli (CCNOT)}
Sur trois qubits $(x,y,z)$, avec $x$ et $y$ pour contrôles et $z$ pour cible,
\[
  \mathrm{CCNOT}=\ket{0}\bra{0}_x\otimes I_y\otimes I_z
   +\ket{1}\bra{1}_x\otimes\Big[\ket{0}\bra{0}_y\otimes I_z+\ket{1}\bra{1}_y\otimes\sigma_x\Big].
\]
Le crochet est un CNOT (contrôle $y$, cible $z$), appliqué seulement si $x=1$ : si $x=0$, rien ;
si $x=1$ et $y=0$, rien ; si $x=1$ et $y=1$, le qubit $z$ est inversé.
\end{definition}

Calcul de $\mathrm{CCNOT}\ket{110}$ :
\begin{align*}
  \mathrm{CCNOT}\ket{110}&=\ket{0}\braket{0}{1}\otimes I\ket{1}\otimes I\ket{0}
     +\ket{1}\braket{1}{1}\otimes\Big[\ket{0}\braket{0}{1}\otimes I\ket{0}
       +\ket{1}\braket{1}{1}\otimes\sigma_x\ket{0}\Big]\\
  &=0+\ket{1}\otimes\Big[0+\ket{1}\otimes\ket{1}\Big]=\ket{111}.
\end{align*}
Pour les huit états de base, $\mathrm{CCNOT}\ket{x,y,z}=\ket{x,\,y,\,z\oplus xy}$ :
\begin{center}
\renewcommand{\arraystretch}{1.3}
\begin{tabular}{@{}ccccccccc@{}}
\toprule
Entrée & $\ket{000}$ & $\ket{001}$ & $\ket{010}$ & $\ket{011}$ & $\ket{100}$ & $\ket{101}$ & $\ket{110}$ & $\ket{111}$ \\
\midrule
Sortie & $\ket{000}$ & $\ket{001}$ & $\ket{010}$ & $\ket{011}$ & $\ket{100}$ & $\ket{101}$ & $\ket{111}$ & $\ket{110}$ \\
\bottomrule
\end{tabular}
\end{center}
Seuls $\ket{110}$ et $\ket{111}$ sont échangés : la matrice est $\mathrm{diag}(I_6,X)$
(un bloc $6\times6$ égal à l'identité, puis $X$),
\[
  \mathrm{CCNOT}=\begingroup\setlength{\arraycolsep}{3pt}\begin{pmatrix}
    1&0&0&0&0&0&0&0\\
    0&1&0&0&0&0&0&0\\
    0&0&1&0&0&0&0&0\\
    0&0&0&1&0&0&0&0\\
    0&0&0&0&1&0&0&0\\
    0&0&0&0&0&1&0&0\\
    0&0&0&0&0&0&0&1\\
    0&0&0&0&0&0&1&0
  \end{pmatrix}\endgroup .
\]

\begin{remarque}[Portes réversibles]
Les portes de Fredkin et de Toffoli sont des permutations des états de base qui échangent
seulement deux d'entre eux : elles sont leurs propres inverses ($\mathrm{C\text{-}SWAP}^2=\mathrm{CCNOT}^2=\mathrm{Id}$)
et unitaires. Avec $z=0$, la Toffoli calcule le ET logique $x\wedge y$ sans rien effacer :
$\ket{x,y,0}\to\ket{x,y,xy}$ ; avec $z=1$, elle calcule le NON-ET. Elle sert à écrire tout calcul
classique de façon réversible.
\end{remarque}

\clearpage
% ---------------------------------------------------------------------
\subsection{Récapitulatif du chapitre}\label{sec:recap4}
% ---------------------------------------------------------------------

\begin{center}
\adjustbox{max width=\linewidth}{%
\begin{tikzpicture}[>=Stealth, font=\footnotesize,
    boite/.style={draw=insarouge, fill=insarouge!6, rounded corners=2pt, thick,
                  minimum width=4.4cm, minimum height=1.5cm, align=center, inner sep=3pt},
    lien/.style={->, very thick, insagris},
    etiq/.style={font=\scriptsize\itshape, text=insagris, inner sep=2pt}]
  \node[boite] (uni) at (0,0)
    {\textbf{Porte unitaire}\\ $U^\dagger U=\mathrm{Id}$, \ $\ket{\Psi}\to U\ket{\Psi}$\\ {\scriptsize\S\,\ref{sec:portes}}};
  \node[boite] (un) at (5.9,0)
    {\textbf{Portes à un qubit}\\ $X$, $Y$, $Z$, $H$, $S$, $T$\\ {\scriptsize\S\,\ref{sec:pauli}--\ref{sec:phase}}};
  \node[boite] (blo) at (11.8,0)
    {\textbf{Rotations de Bloch}\\ $X,Y,Z,H$ : angle $\pi$ ; $P_\theta$ autour de $z$\\
     {\scriptsize\S\,\ref{sec:portes-bloch}}};
  \node[boite] (cir) at (11.8,-2.5)
    {\textbf{Circuits}\\ fils, boîtes, $U_2U_1$\\ {\scriptsize\S\,\ref{sec:circuits}, \ref{sec:op2}}};
  \node[boite] (cn) at (5.9,-2.5)
    {\textbf{CNOT et SWAP}\\ $\ket{a,b}\to\ket{a\oplus b,b}$\\ {\scriptsize\S\,\ref{sec:cnot}, \ref{sec:swap}}};
  \node[boite] (cu) at (0,-2.5)
    {\textbf{Portes contrôlées}\\ $\mathrm{C}U=\mathrm{diag}(I,U)$, Fredkin, Toffoli\\
     {\scriptsize\S\,\ref{sec:controlees}}};
  \draw[lien] (uni) -- (un);
  \draw[lien] (un) -- (blo);
  \draw[lien] (blo) -- (cir);
  \draw[lien] (cir) -- (cn) node[etiq, midway, above=1pt] {2 qubits};
  \draw[lien] (cn) -- (cu) node[etiq, midway, above=1pt] {contrôle};
\end{tikzpicture}}
\end{center}

\begin{aretenir}
\begin{itemize}
  \item Une porte est un opérateur unitaire : elle conserve la norme, elle est réversible
        ($U^{-1}=U^\dagger$) et agit sur la sphère de Bloch comme une rotation. C'est l'évolution de
        Schrödinger, $U=\ee^{-\ii\Hop t/\hbar}$.
  \item $X$ échange $\ket{0}$ et $\ket{1}$, $Z$ change le signe de $\ket{1}$, $Y=\ii XZ$ fait les deux ;
        $H$ échange les bases $Z$ et $X$ et crée des superpositions ; $P_\theta$, $S=P_{\pi/2}$ et
        $T=P_{\pi/4}$ font tourner autour de $z$ ($T^2=S$, $S^2=Z$).
  \item Dans un circuit, le temps va de gauche à droite et la matrice est le produit dans l'ordre
        inverse, $U_2U_1$. Le fil du haut porte le dernier chiffre du ket. Des portes en parallèle
        forment un produit tensoriel $U_A\otimes U_B$.
  \item Le CNOT inverse la cible si le contrôle vaut $1$ : $\ket{a,b}\to\ket{a\oplus b,b}$. Le SWAP
        échange deux qubits ; $\mathrm{CNOT}_{21}=\mathrm{SWAP}\,\mathrm{CNOT}_{12}\,\mathrm{SWAP}$.
  \item Une porte contrôlée a pour matrice $\mathrm{diag}(I,U)$. Fredkin (contrôlé-SWAP) et
        Toffoli (CCNOT) sont des permutations des états de base : réversibles et universelles
        pour le calcul classique.
\end{itemize}
\end{aretenir}

\begin{center}
\renewcommand{\arraystretch}{1.45}
\begin{tabular}{@{}l l r@{}}
\toprule
\textbf{Notion} & \textbf{Formule} & \textbf{§} \\
\midrule
Porte unitaire & $U^\dagger U=\mathrm{Id}$, \quad $\bra{\Psi}U^\dagger U\ket{\Psi}=1$ & \ref{sec:portes} \\
Pauli & $X\ket{0}=\ket{1}$, \quad $Z\ket{1}=-\ket{1}$, \quad $X^2=Y^2=Z^2=\mathrm{Id}$ & \ref{sec:pauli} \\
Hadamard & $H\ket{0}=\ket{+}$, \quad $H\ket{1}=\ket{-}$, \quad $H^2=\mathrm{Id}$ & \ref{sec:hadamard} \\
Phase & $P_\theta=\mathrm{diag}(1,\ee^{\ii\theta})$, \quad $S=P_{\pi/2}$, \quad $T=P_{\pi/4}$ & \ref{sec:phase} \\
Circuit & $U_1$ puis $U_2$ : \quad $U_2U_1$ & \ref{sec:circuits} \\
Portes locales & $(U\otimes I)(\ket{u}\otimes\ket{v})=U\ket{u}\otimes\ket{v}$ & \ref{sec:op2} \\
CNOT & $\ket{a,b}\to\ket{a\oplus b,b}$, \quad $\mathrm{CNOT}=I\otimes\ket{0}\bra{0}+X\otimes\ket{1}\bra{1}$
  & \ref{sec:cnot} \\
SWAP & $\mathrm{SWAP}\ket{a,b}=\ket{b,a}$, \quad $\mathrm{SWAP}=\mathrm{CNOT}_{12}\mathrm{CNOT}_{21}\mathrm{CNOT}_{12}$
  & \ref{sec:swap} \\
Contrôlé-$U$ & $\mathrm{C}U=\ket{0}\bra{0}\otimes I+\ket{1}\bra{1}\otimes U=\mathrm{diag}(I,U)$ & \ref{sec:controlees} \\
Toffoli & $\mathrm{CCNOT}\ket{x,y,z}=\ket{x,y,z\oplus xy}$ & \ref{sec:controlees} \\
\bottomrule
\end{tabular}
\end{center}

\begin{savoirfaire}
\begin{itemize}
  \item Appliquer une porte à un état (par la matrice ou par les ket-bra) et identifier sa rotation
        de Bloch (exercices~\ref{ex:conj} et~\ref{ex:phases}).
  \item Lire un circuit, écrire sa matrice dans le bon ordre et suivre l'état à chaque étape
        (exercice~\ref{ex:kickback}).
  \item Écrire la matrice de $U\otimes I$, du CNOT, du SWAP ou d'une porte contrôlée, et calculer
        leur action sur les états de base.
  \item Reconnaître une porte réversible qui calcule une fonction booléenne
        (exercice~\ref{ex:toffoli}).
\end{itemize}
\end{savoirfaire}
.
%  Il ne contient ni préambule ni \begin{document} : les environnements
%  (definition, resultat, remarque, aretenir, savoirfaire, exercice, correction),
%  les macros (\ii, \ee, \dd, \pd, \pdd, \Hop, \ket, \bra, \braket), les symboles
%  de circuits (ctl, cible, croix, mesure, porte) et les couleurs (insarouge,
%  insagris) sont définis dans main.tex. Un chapitre = une \section.
% =====================================================================

% Macros de Dirac (déjà définies dans main.tex ; ce bloc évite l'erreur
% « Undefined control sequence \ket » si main.tex est une ancienne version).
\providecommand{\ket}[1]{|#1\rangle}
\providecommand{\bra}[1]{\langle #1|}
\providecommand{\braket}[2]{\langle #1|#2\rangle}

% =====================================================================
\section{Portes quantiques et systèmes multi-qubits}\label{chap:portes}
% =====================================================================

Un calcul quantique est une suite de transformations unitaires appliquées à des qubits. Ce
chapitre décrit les portes à un qubit (Pauli, Hadamard, phase) et leur effet sur la sphère de
Bloch, la lecture d'un circuit, puis les portes qui agissent sur plusieurs qubits : CNOT, SWAP et
portes contrôlées.

% ---------------------------------------------------------------------
\subsection{Portes quantiques}\label{sec:portes}
% ---------------------------------------------------------------------

Un calcul quantique transforme des états. Les opérations sur un seul qubit sont appelées
\emph{portes} ; les sections suivantes passent à plusieurs qubits.

\begin{resultat}{Postulat 4 : évolution d'un système isolé}
L'évolution d'un système isolé est décrite par un opérateur unitaire : $\ket{\Psi}\to U\ket{\Psi}$
avec $U^\dagger U=\mathrm{Id}$.
\end{resultat}

C'est la forme de l'équation de Schrödinger du chapitre~\ref{chap:histoire}. Pour un hamiltonien
$\Hop$ indépendant du temps, $\ii\hbar\,\partial_t\ket{\Psi}=\Hop\ket{\Psi}$ a pour solution
$\ket{\Psi(t)}=U(t)\ket{\Psi(0)}$ avec $U(t)=\ee^{-\ii\Hop t/\hbar}$. Comme $\Hop$ est hermitien,
$U(t)^\dagger=\ee^{+\ii\Hop t/\hbar}$ et $U^\dagger U=\mathrm{Id}$ : l'évolution est unitaire. Une
porte est une telle évolution, obtenue en contrôlant $\Hop$ pendant une durée donnée.

\begin{definition}{Porte quantique}
Une porte à un qubit est un opérateur $U$ (matrice $2\times2$) \emph{unitaire} :
$U^\dagger U=\mathrm{Id}$. Elle transforme $\ket{\Psi}\to U\ket{\Psi}$, conserve la norme
($\bra{\Psi}U^\dagger U\ket{\Psi}=\braket{\Psi}{\Psi}=1$) et est réversible
($U^{-1}=U^\dagger$). Sur la sphère de Bloch (section~\ref{sec:bloch}), c'est une
\emph{rotation}.
\end{definition}

% ---------------------------------------------------------------------
\subsection{Portes de Pauli}\label{sec:pauli}
% ---------------------------------------------------------------------

Les matrices de Pauli de la section~\ref{sec:op-dirac} sont trois portes.

\begin{definition}{Portes $X$, $Y$, $Z$}
\begin{center}
\renewcommand{\arraystretch}{1}
\begin{tabular}{@{}lccl@{}}
\textbf{Porte} & \textbf{Matrice} & \textbf{Ket-bra} & \textbf{Nom} \\[0.8ex]
$X=\sigma_x$ & $\begin{pmatrix}0 & 1\\ 1 & 0\end{pmatrix}$
  & $\ket{0}\bra{1}+\ket{1}\bra{0}$ & bit flip \\[2.2ex]
$Y=\sigma_y$ & $\begin{pmatrix}0 & -\ii\\ \ii & 0\end{pmatrix}$
  & $-\ii\ket{0}\bra{1}+\ii\ket{1}\bra{0}$ &  \\[2.2ex]
$Z=\sigma_z$ & $\begin{pmatrix}1 & 0\\ 0 & -1\end{pmatrix}$
  & $\ket{0}\bra{0}-\ket{1}\bra{1}$ & phase flip \\
\end{tabular}
\end{center}
\end{definition}

\paragraph{Porte $X$ (inverseur).} Avec $\braket{0}{0}=\braket{1}{1}=1$ et
$\braket{0}{1}=\braket{1}{0}=0$ :
\begin{align*}
  X\ket{0}&=\ket{0}\braket{1}{0}+\ket{1}\braket{0}{0}=0\cdot\ket{0}+1\cdot\ket{1}=\ket{1},\\
  X\ket{1}&=\ket{0}\braket{1}{1}+\ket{1}\braket{0}{1}=1\cdot\ket{0}+0\cdot\ket{1}=\ket{0}.
\end{align*}
$X$ échange $\ket{0}$ et $\ket{1}$ : c'est le NOT quantique (\emph{bit flip}). Sur la
sphère de Bloch, c'est une rotation de $\pi$ autour de l'axe $x$ : le pôle nord devient
le pôle sud. Les états de l'axe $x$ ne bougent pas : $X\ket{+}=\ket{+}$ et
$X\ket{-}=-\ket{-}$.

\paragraph{Porte $Z$ (changement de phase).}
\begin{align*}
  Z\ket{0}&=\ket{0}\braket{0}{0}-\ket{1}\braket{1}{0}=\ket{0},\\
  Z\ket{1}&=\ket{0}\braket{0}{1}-\ket{1}\braket{1}{1}=-\ket{1},
\end{align*}
donc $Z$ change le signe de la composante sur $\ket{1}$ (\emph{phase flip}). Sur les
états de la base $X$, par linéarité :
\begin{align*}
  Z\ket{+}&=\tfrac{1}{\sqrt2}\big(Z\ket{0}+Z\ket{1}\big)=\tfrac{1}{\sqrt2}\big(\ket{0}-\ket{1}\big)=\ket{-},\\
  Z\ket{-}&=\tfrac{1}{\sqrt2}\big(Z\ket{0}-Z\ket{1}\big)=\tfrac{1}{\sqrt2}\big(\ket{0}+\ket{1}\big)=\ket{+}.
\end{align*}
Sous forme matricielle, $Z\ket{+}=\begin{pmatrix}1 & 0\\ 0 & -1\end{pmatrix}
\frac{1}{\sqrt2}\begin{pmatrix}1\\ 1\end{pmatrix}=\frac{1}{\sqrt2}\begin{pmatrix}1\\ -1\end{pmatrix}=\ket{-}$.
$Z$ est une rotation de $\pi$ autour de l'axe $z$ : elle laisse $\ket{0}$ et $\ket{1}$ en
place et échange $\ket{+}$ et $\ket{-}$.

\paragraph{Porte $Y$.} Avec $Y=-\ii\ket{0}\bra{1}+\ii\ket{1}\bra{0}$ :
\begin{align*}
  Y\ket{0}&=-\ii\ket{0}\braket{1}{0}+\ii\ket{1}\braket{0}{0}=\ii\ket{1},\\
  Y\ket{1}&=-\ii\ket{0}\braket{1}{1}+\ii\ket{1}\braket{0}{1}=-\ii\ket{0}.
\end{align*}
Sur les états de la base $X$ :
\begin{align*}
  Y\ket{+}&=\tfrac{1}{\sqrt2}\big(\ii\ket{1}-\ii\ket{0}\big)=-\ii\,\tfrac{1}{\sqrt2}\big(\ket{0}-\ket{1}\big)=-\ii\ket{-},\\
  Y\ket{-}&=\tfrac{1}{\sqrt2}\big(\ii\ket{1}+\ii\ket{0}\big)=\ii\,\tfrac{1}{\sqrt2}\big(\ket{0}+\ket{1}\big)=\ii\ket{+}.
\end{align*}
Un facteur global ($-\ii$ ou $\ii$) ne change pas l'état physique : $Y$ échange bien
$\ket{+}$ et $\ket{-}$. Sur les états de l'axe $y$ :
\begin{align*}
  Y\ket{{+}\ii}&=\tfrac{1}{\sqrt2}\big(Y\ket{0}+\ii Y\ket{1}\big)
   =\tfrac{1}{\sqrt2}\big(\ii\ket{1}+\ket{0}\big)=\ket{{+}\ii},\\
  Y\ket{{-}\ii}&=\tfrac{1}{\sqrt2}\big(Y\ket{0}-\ii Y\ket{1}\big)
   =\tfrac{1}{\sqrt2}\big(\ii\ket{1}-\ket{0}\big)=-\ket{{-}\ii}.
\end{align*}
$Y$ est une rotation de $\pi$ autour de l'axe $y$.

\begin{remarque}[Chaque Pauli fixe son axe]
$Z\ket{0}=\ket{0}$, $Z\ket{1}=-\ket{1}$ ; $X\ket{\pm}=\pm\ket{\pm}$ ; $Y\ket{{\pm}\ii}=\pm\ket{{\pm}\ii}$.
Les bases $Z$, $X$, $Y$ de la section~\ref{sec:bloch} sont les bases propres de
$Z$, $X$, $Y$. Une porte de Pauli laisse donc son propre axe invariant et retourne les
deux autres ; c'est pourquoi les exemples utilisent les états des deux autres axes.
\end{remarque}

% ---------------------------------------------------------------------
\subsection{Porte de Hadamard}\label{sec:hadamard}
% ---------------------------------------------------------------------

\begin{definition}{Porte de Hadamard}
\[
  H=\frac{1}{\sqrt2}\begin{pmatrix}1 & 1\\ 1 & -1\end{pmatrix}
  =\ket{+}\bra{0}+\ket{-}\bra{1}.
\]
\end{definition}

\begin{align*}
  H\ket{0}&=\frac{1}{\sqrt2}\begin{pmatrix}1 & 1\\ 1 & -1\end{pmatrix}\begin{pmatrix}1\\ 0\end{pmatrix}
    =\frac{1}{\sqrt2}\begin{pmatrix}1\\ 1\end{pmatrix}=\ket{+},\\[0.6ex]
  H\ket{1}&=\frac{1}{\sqrt2}\begin{pmatrix}1 & 1\\ 1 & -1\end{pmatrix}\begin{pmatrix}0\\ 1\end{pmatrix}
    =\frac{1}{\sqrt2}\begin{pmatrix}1\\ -1\end{pmatrix}=\ket{-},\\[0.6ex]
  H\ket{+}&=\frac{1}{2}\begin{pmatrix}1 & 1\\ 1 & -1\end{pmatrix}\begin{pmatrix}1\\ 1\end{pmatrix}
    =\frac{1}{2}\begin{pmatrix}2\\ 0\end{pmatrix}=\ket{0},\qquad
  H\ket{-}=\frac{1}{2}\begin{pmatrix}0\\ 2\end{pmatrix}=\ket{1}.
\end{align*}
Comme $H\ket{+}=\ket{0}$ et $H\ket{-}=\ket{1}$, on a $H^2=\mathrm{Id}$ ($H$ est son propre
inverse). Sur la sphère de Bloch, $H$ est une rotation de $\pi$ autour de l'axe
$(x+z)/\sqrt2$ : elle échange l'axe $x$ et l'axe $z$.

\begin{remarque}[À quoi sert $H$ ?]
\begin{itemize}[nosep]
  \item Elle \emph{crée une superposition} à partir d'un état de base :
        $H\ket{0}=\ket{+}$ donne $P(0)=P(1)=\frac12$. C'est la première étape de
        la plupart des algorithmes quantiques.
  \item Elle \emph{change de base} : elle envoie la base $Z$ sur la base $X$ et
        inversement. Mesurer dans la base $X$ revient à appliquer $H$ puis à mesurer
        dans la base $Z$ (figure~\ref{fig:mesurex}).
  \item Elle est réversible ($H^2=\mathrm{Id}$).
\end{itemize}
\end{remarque}

% ---------------------------------------------------------------------
\subsection{Portes de phase}\label{sec:phase}
% ---------------------------------------------------------------------

\begin{definition}{Portes de phase $P_\theta$, $S$ et $T$}
\[
  P_\theta=\begin{pmatrix}1 & 0\\ 0 & \ee^{\ii\theta}\end{pmatrix}
  =\ket{0}\bra{0}+\ee^{\ii\theta}\ket{1}\bra{1},
  \qquad
  S=P_{\pi/2}=\begin{pmatrix}1 & 0\\ 0 & \ii\end{pmatrix},
  \qquad
  T=P_{\pi/4}=\begin{pmatrix}1 & 0\\ 0 & \dfrac{1+\ii}{\sqrt2}\end{pmatrix}.
\]
\end{definition}

Ici $\ee^{\ii\pi/2}=\ii$ et $\ee^{\ii\pi/4}=\cos\frac\pi4+\ii\sin\frac\pi4=\frac{1+\ii}{\sqrt2}$.
$P_\theta$ laisse $\ket{0}$ inchangé et multiplie $\ket{1}$ par $\ee^{\ii\theta}$ :
\[
  P_\theta\big(\alpha\ket{0}+\beta\ket{1}\big)=\alpha\ket{0}+\ee^{\ii\theta}\beta\ket{1}.
\]
Dans la paramétrisation de Bloch, $\varphi\to\varphi+\theta$ : $P_\theta$ est une rotation
d'angle $\theta$ autour de l'axe $z$. On retrouve $P_\pi=Z$, et $S^2=Z$, $T^2=S$.

\paragraph{Exemples.}
\begin{align*}
  S\ket{+}&=\tfrac{1}{\sqrt2}\big(\ket{0}+\ii\ket{1}\big)=\ket{{+}\ii},\\
  S\ket{{+}\ii}&=\tfrac{1}{\sqrt2}\big(\ket{0}+\ii\cdot\ii\ket{1}\big)=\tfrac{1}{\sqrt2}\big(\ket{0}-\ket{1}\big)=\ket{-},\\
  T\ket{+}&=\tfrac{1}{\sqrt2}\begin{pmatrix}1 & 0\\ 0 & \frac{1+\ii}{\sqrt2}\end{pmatrix}\begin{pmatrix}1\\ 1\end{pmatrix}
    =\tfrac{1}{\sqrt2}\begin{pmatrix}1\\ \frac{1+\ii}{\sqrt2}\end{pmatrix}
    =\tfrac{1}{\sqrt2}\big(\ket{0}+\ee^{\ii\pi/4}\ket{1}\big)\quad\big(\text{angle polaire }\tfrac\pi2,\ \text{azimut }\tfrac\pi4\big).
\end{align*}
$S$ tourne donc $\ket{+}$ d'un quart de tour autour de $z$ (de l'axe $x$ à l'axe $y$),
puis d'un second quart de tour jusqu'à $\ket{-}$ ; $T$ fait un huitième de tour.

\begin{remarque}[Une porte de phase est une évolution libre]
Pour $\Hop=\frac{\hbar\omega}{2}Z$, $U(t)=\ee^{-\ii\Hop t/\hbar}=\mathrm{diag}\big(\ee^{-\ii\omega t/2},\ee^{\ii\omega t/2}\big)
=\ee^{-\ii\omega t/2}\,P_{\omega t}$. Au facteur global près, laisser évoluer le qubit pendant la
durée $t$ applique la porte de phase $P_{\omega t}$ : une rotation d'angle $\omega t$ autour de
l'axe $z$.
\end{remarque}

% ---------------------------------------------------------------------
\subsection{Portes et sphère de Bloch}\label{sec:portes-bloch}
% ---------------------------------------------------------------------

Toute porte à un qubit est une rotation de la sphère de Bloch. La figure~\ref{fig:portes} montre
les axes des portes usuelles et le tableau qui suit les résume.

\begin{figure}[!ht]
\centering
\begin{tikzpicture}[x={(-1.63cm,-0.56cm)}, y={(1.63cm,-0.56cm)}, z={(0cm,2.16cm)}, >=Stealth, scale=0.9]
  % sphère et grands cercles
  \shade[ball color=insagris!15, opacity=0.55] (0,0,0) circle[radius=2.3cm];
  \begin{scope}[insagris!75, thin]
    \draw[domain=-45:135, samples=60, variable=\t] plot ({cos(\t)},{sin(\t)},0);
    \draw[domain=135:315, samples=60, variable=\t, dashed] plot ({cos(\t)},{sin(\t)},0);
    \draw[domain=-62.7:117.3, samples=60, variable=\t] plot ({cos(\t)},0,{sin(\t)});
    \draw[domain=117.3:297.3, samples=60, variable=\t, dashed] plot ({cos(\t)},0,{sin(\t)});
    \draw[domain=-62.7:117.3, samples=60, variable=\t] plot (0,{cos(\t)},{sin(\t)});
    \draw[domain=117.3:297.3, samples=60, variable=\t, dashed] plot (0,{cos(\t)},{sin(\t)});
  \end{scope}
  % axes x, y, z (portes X, Y, Z)
  \draw[insagris, dashed] (0,0,0) -- (-1,0,0);
  \draw[insagris, dashed] (0,0,0) -- (0,-1,0);
  \draw[insagris, dashed] (0,0,0) -- (0,0,-1);
  \draw[insagris] (-1,0,0) -- (-1.3,0,0);
  \draw[insagris] (0,-1,0) -- (0,-1.3,0);
  \draw[insagris] (0,0,-1) -- (0,0,-1.3);
  \draw[insagris, -{Stealth[length=2.2mm]}] (0,0,0) -- (1.3,0,0);
  \draw[insagris, -{Stealth[length=2.2mm]}] (0,0,0) -- (0,1.3,0);
  \draw[insagris, -{Stealth[length=2.2mm]}] (0,0,0) -- (0,0,1.3);
  \node[anchor=east]  at (1.33,0,0) {$x$};
  \node[anchor=west]  at (0,1.33,0) {$y$};
  \node[anchor=south] at (0,0,1.33) {$z$};
  % axe de la porte H : (x+z)/sqrt(2)
  \draw[insarouge, dashed] (0,0,0) -- (-0.707,0,-0.707);
  \draw[insarouge, very thick, -{Stealth[length=2.6mm]}] (0,0,0) -- (0.98,0,0.98);
  \node[insarouge, anchor=south east] at (0.98,0,0.98) {axe de $H$};
  % rotation autour de z (portes de phase)
  \draw[insarouge, thick, -{Stealth[length=2.2mm]}, domain=15:105, samples=30, variable=\t]
    plot ({1.15*cos(\t)},{1.15*sin(\t)},0);
  \node[insarouge, anchor=north west, font=\small] at ({1.15*cos(60)},{1.15*sin(60)},0) {$P_\theta$};
\end{tikzpicture}
\caption{Portes et rotations sur la sphère de Bloch. Les portes $X$, $Y$, $Z$ sont des
rotations de $\pi$ autour des axes $x$, $y$, $z$ ; $H$ est une rotation de $\pi$ autour de
l'axe $(x+z)/\sqrt2$ (rouge) ; $P_\theta$ (donc $Z$, $S$, $T$) tourne d'un angle $\theta$
autour de $z$.}
\label{fig:portes}
\end{figure}

\begin{center}
\renewcommand{\arraystretch}{1}
\begin{tabular}{@{}lccc@{}}
\toprule
\textbf{Porte} & \textbf{Matrice} & \textbf{Effet sur la base} & \textbf{Rotation de Bloch} \\
\midrule
$X$ & $\begin{pmatrix}0 & 1\\ 1 & 0\end{pmatrix}$ & $\ket{0}\leftrightarrow\ket{1}$ & $\pi$ autour de $x$ \\[2.2ex]
$Y$ & $\begin{pmatrix}0 & -\ii\\ \ii & 0\end{pmatrix}$ & $\ket{0}\to\ii\ket{1},\ \ket{1}\to-\ii\ket{0}$ & $\pi$ autour de $y$ \\[2.2ex]
$Z$ & $\begin{pmatrix}1 & 0\\ 0 & -1\end{pmatrix}$ & $\ket{1}\to-\ket{1},\ \ket{\pm}\leftrightarrow\ket{\mp}$ & $\pi$ autour de $z$ \\[2.2ex]
$H$ & $\dfrac{1}{\sqrt2}\begin{pmatrix}1 & 1\\ 1 & -1\end{pmatrix}$ & $\ket{0}\leftrightarrow\ket{+},\ \ket{1}\leftrightarrow\ket{-}$ & $\pi$ autour de $\frac{x+z}{\sqrt2}$ \\[2.2ex]
$S$ & $\begin{pmatrix}1 & 0\\ 0 & \ii\end{pmatrix}$ & $\ket{1}\to\ii\ket{1}$ & $\frac\pi2$ autour de $z$ \\[2.2ex]
$T$ & $\begin{pmatrix}1 & 0\\ 0 & \frac{1+\ii}{\sqrt2}\end{pmatrix}$ & $\ket{1}\to\ee^{\ii\pi/4}\ket{1}$ & $\frac\pi4$ autour de $z$ \\
\bottomrule
\end{tabular}
\end{center}

\begin{figure}[!ht]
\centering
\begin{tikzpicture}[x={(-1.63cm,-0.56cm)}, y={(1.63cm,-0.56cm)}, z={(0cm,2.16cm)}, >=Stealth, scale=0.9]
  % sphère et grands cercles
  \shade[ball color=insagris!15, opacity=0.55] (0,0,0) circle[radius=2.3cm];
  \begin{scope}[insagris!75, thin]
    \draw[domain=-45:135, samples=60, variable=\t] plot ({cos(\t)},{sin(\t)},0);
    \draw[domain=135:315, samples=60, variable=\t, dashed] plot ({cos(\t)},{sin(\t)},0);
    \draw[domain=-62.7:117.3, samples=60, variable=\t] plot ({cos(\t)},0,{sin(\t)});
    \draw[domain=117.3:297.3, samples=60, variable=\t, dashed] plot ({cos(\t)},0,{sin(\t)});
    \draw[domain=-62.7:117.3, samples=60, variable=\t] plot (0,{cos(\t)},{sin(\t)});
    \draw[domain=117.3:297.3, samples=60, variable=\t, dashed] plot (0,{cos(\t)},{sin(\t)});
  \end{scope}
  % axes x, y, z (portes X, Y, Z)
  \draw[insagris, dashed] (0,0,0) -- (-1,0,0);
  \draw[insagris, dashed] (0,0,0) -- (0,-1,0);
  \draw[insagris, dashed] (0,0,0) -- (0,0,-1);
  \draw[insagris] (-1,0,0) -- (-1.3,0,0);
  \draw[insagris] (0,-1,0) -- (0,-1.3,0);
  \draw[insagris] (0,0,-1) -- (0,0,-1.3);
  \draw[insagris, -{Stealth[length=2.2mm]}] (0,0,0) -- (1.3,0,0);
  \draw[insagris, -{Stealth[length=2.2mm]}] (0,0,0) -- (0,1.3,0);
  \draw[insagris, -{Stealth[length=2.2mm]}] (0,0,0) -- (0,0,1.3);
  % trajet : |0> --H--> |+> --S--> |+i>
  \draw[insarouge, very thick, -{Stealth[length=2.6mm]}, domain=3:87, samples=30, variable=\t]
    plot ({sin(\t)},0,{cos(\t)});
  \draw[insarouge, very thick, -{Stealth[length=2.6mm]}, domain=3:87, samples=30, variable=\t]
    plot ({cos(\t)},{sin(\t)},0);
  \foreach \ax/\ay/\az in {0/0/1, 1/0/0, 0/1/0} \fill[insarouge] (\ax,\ay,\az) circle[radius=2pt];
  \node[insarouge, anchor=south] at (0,0,1.33) {$\ket{0}$};
  \node[insarouge, anchor=east] at (1.33,0,0) {$\ket{+}$};
  \node[insarouge, anchor=west] at (0,1.33,0) {$\ket{{+}\ii}$};
  \node[insarouge, font=\small, anchor=south west] at ({sin(45)},0,{cos(45)}) {$H$};
  \node[insarouge, font=\small, anchor=north west] at ({cos(45)},{sin(45)},0) {$S$};
\end{tikzpicture}
\caption{Trajet d'un état sous le circuit $H$ puis $S$ (flèches schématiques). $H$ envoie
$\ket{0}$ sur $\ket{+}$ ; $S$, rotation d'un quart de tour autour de $z$, envoie $\ket{+}$ sur
$\ket{{+}\ii}$. Les trois états sont alignés sur les axes $z$, $x$ et $y$.}
\label{fig:trajet}
\end{figure}

% ---------------------------------------------------------------------
\subsection{Circuits quantiques}\label{sec:circuits}
% ---------------------------------------------------------------------

Un \emph{circuit} représente un calcul : chaque qubit est un fil horizontal, le temps va de gauche
à droite, et chaque porte est un symbole posé sur les fils qu'elle concerne. La
figure~\ref{fig:symboles} donne les symboles utilisés dans ce document.

\begin{figure}[!ht]
\centering
\begin{adjustbox}{max width=\linewidth}
\begin{tikzpicture}[>=Stealth, font=\small,
    nom/.style={font=\footnotesize, align=center, anchor=north}]
  % (a) porte à un qubit
  \draw (0,0) -- (1.8,0);
  \node[porte] at (0.9,0) {$U$};
  \node[nom] at (0.9,-1.15) {porte $U$\\ à un qubit};
  % (b) CNOT
  \begin{scope}[xshift=3.2cm]
    \draw (0,0.6) -- (1.8,0.6);
    \draw (0,-0.6) -- (1.8,-0.6);
    \draw[insarouge, thick] (0.9,0.6) -- (0.9,-0.6);
    \pic at (0.9,0.6) {ctl};
    \pic at (0.9,-0.6) {cible};
    \node[nom] at (0.9,-1.15) {CNOT\\ (contrôle $\bullet$, cible $\oplus$)};
  \end{scope}
  % (c) SWAP
  \begin{scope}[xshift=6.4cm]
    \draw (0,0.6) -- (1.8,0.6);
    \draw (0,-0.6) -- (1.8,-0.6);
    \draw[insarouge, thick] (0.9,0.6) -- (0.9,-0.6);
    \pic at (0.9,0.6) {croix};
    \pic at (0.9,-0.6) {croix};
    \node[nom] at (0.9,-1.15) {SWAP\\ (échange)};
  \end{scope}
  % (d) contrôlé-U
  \begin{scope}[xshift=9.6cm]
    \draw (0,0.6) -- (1.8,0.6);
    \draw (0,-0.6) -- (1.8,-0.6);
    \draw[insarouge, thick] (0.9,-0.6) -- (0.9,0.3);
    \pic at (0.9,-0.6) {ctl};
    \node[porte] at (0.9,0.6) {$U$};
    \node[nom] at (0.9,-1.15) {contrôlé-$U$\\ (contrôle en bas)};
  \end{scope}
  % (e) mesure
  \begin{scope}[xshift=12.8cm]
    \draw (0,0) -- (0.6,0);
    \pic at (0.9,0) {mesure};
    \draw (1.2,0.03) -- (2.0,0.03);
    \draw (1.2,-0.03) -- (2.0,-0.03);
    \node[nom] at (1.0,-1.15) {mesure\\ (fil double : bit)};
  \end{scope}
\end{tikzpicture}
\end{adjustbox}
\caption{Symboles des circuits quantiques. Le point noir $\bullet$ marque un qubit de contrôle,
le symbole $\oplus$ une cible qui subit $X$, les deux croix un échange. Après une mesure, le fil
double transporte un bit classique.}
\label{fig:symboles}
\end{figure}

\begin{remarque}[Conventions de lecture]
\begin{itemize}[nosep]
  \item Le temps va de gauche à droite. Deux portes successives $U_1$ puis $U_2$ se multiplient
        dans l'ordre inverse : la matrice du circuit est $U_2U_1$, car $U_1$ agit d'abord sur le ket.
  \item Avec plusieurs fils, le fil du \emph{haut} porte le \emph{dernier} chiffre du ket et le
        fil du \emph{bas} le premier : si $A$ est le premier qubit et $B$ le second,
        $\ket{a,b}=\ket{a}_A\otimes\ket{b}_B$ et $B$ est dessiné au-dessus de $A$.
  \item Deux portes placées dans la même colonne, sur deux fils, agissent en parallèle : c'est un
        produit tensoriel $U_A\otimes U_B$ (section~\ref{sec:op2}).
\end{itemize}
\end{remarque}

\paragraph{Exemple.} Le circuit de la figure~\ref{fig:circuit-hs} applique $H$ puis $S$ à
$\ket{0}$. Sa matrice est $SH$ :
\[
  SH=\begin{pmatrix}1&0\\0&\ii\end{pmatrix}\frac{1}{\sqrt2}\begin{pmatrix}1&1\\1&-1\end{pmatrix}
  =\frac{1}{\sqrt2}\begin{pmatrix}1&1\\ \ii&-\ii\end{pmatrix},
  \qquad
  SH\ket{0}=\frac{1}{\sqrt2}\begin{pmatrix}1\\ \ii\end{pmatrix}=\ket{{+}\ii}.
\]
Le circuit prépare donc l'état $\ket{{+}\ii}$ de la base $Y$ à partir de $\ket{0}$
(figure~\ref{fig:trajet}). L'ordre compte : $HS\ket{0}=H\ket{0}=\ket{+}$, car $S\ket{0}=\ket{0}$.

\begin{figure}[!ht]
\centering
\begin{tikzpicture}[>=Stealth, font=\small]
  \draw (0,0) -- (6.6,0);
  \node[anchor=east] at (0,0) {$\ket{0}$};
  \node[porte] at (2.0,0) {$H$};
  \node[porte] at (4.3,0) {$S$};
  \node[anchor=west] at (6.6,0) {$\ket{{+}\ii}$};
  \foreach \x in {0.7, 3.15, 5.5} \draw[dashed, insagris!70] (\x,0.6) -- (\x,-0.95);
  \node[font=\footnotesize, anchor=north] at (0.7,-0.95) {$\ket{0}$};
  \node[font=\footnotesize, anchor=north] at (3.15,-0.95) {$\ket{+}$};
  \node[font=\footnotesize, anchor=north] at (5.5,-0.95) {$\ket{{+}\ii}$};
\end{tikzpicture}
\caption{Circuit $H$ puis $S$ appliqué à $\ket{0}$. Les états intermédiaires sont indiqués sous
les traits pointillés ; la matrice du circuit est $SH$.}
\label{fig:circuit-hs}
\end{figure}

\paragraph{Mesure.} Le symbole de mesure (compteur) donne le résultat $0$ ou $1$ dans la base $Z$.
Mesurer dans la base $X$ revient à appliquer $H$ puis à mesurer dans la base $Z$
(figure~\ref{fig:mesurex}) : comme $H\ket{\pm}=\ket{0}$ ou $\ket{1}$, le résultat $0$ correspond à
$\ket{+}$ et le résultat $1$ à $\ket{-}$, avec $P(\pm)=|\braket{\pm}{\Psi}|^2$.

\begin{figure}[!ht]
\centering
\begin{tikzpicture}[>=Stealth, font=\small]
  % (a) base Z
  \draw (0,0) -- (1.2,0);
  \node[anchor=east] at (0,0) {$\ket{\Psi}$};
  \pic at (1.5,0) {mesure};
  \draw (1.8,0.03) -- (2.6,0.03);
  \draw (1.8,-0.03) -- (2.6,-0.03);
  \node[anchor=west] at (2.7,0) {$0$ ou $1$};
  \node[font=\footnotesize] at (1.5,-0.95) {(a) mesure en base $Z$};
  % (b) base X
  \begin{scope}[xshift=7.5cm]
    \draw (0,0) -- (2.6,0);
    \node[anchor=east] at (0,0) {$\ket{\Psi}$};
    \node[porte] at (1.0,0) {$H$};
    \pic at (2.3,0) {mesure};
    \draw (2.6,0.03) -- (3.4,0.03);
    \draw (2.6,-0.03) -- (3.4,-0.03);
    \node[anchor=west] at (3.5,0) {$+$ ou $-$};
    \node[font=\footnotesize] at (1.9,-0.95) {(b) mesure en base $X$};
  \end{scope}
\end{tikzpicture}
\caption{(a) Mesure dans la base $Z$. (b) Mesure dans la base $X$ : porte de Hadamard, puis
mesure dans la base $Z$ ; les résultats $0$ et $1$ se lisent $+$ et $-$.}
\label{fig:mesurex}
\end{figure}

% ---------------------------------------------------------------------
\subsection{Portes locales sur deux qubits}\label{sec:op2}
% ---------------------------------------------------------------------

Une porte appliquée à un seul qubit d'un système de deux s'écrit avec un produit tensoriel
(section~\ref{sec:tenseur}) : $U\otimes I$ agit sur $A$ (premier facteur, fil du bas), $I\otimes U$
agit sur $B$, et $U_A\otimes U_B$ agit sur les deux en parallèle (figure~\ref{fig:local}).

\begin{figure}[!ht]
\centering
\begin{tikzpicture}[>=Stealth, font=\small]
  \foreach \xs/\haut/\bas/\tit in {0/{}/{$X$}/{(a) $X\otimes I$},
                                   5.4/{$X$}/{}/{(b) $I\otimes X$},
                                   10.8/{$H$}/{$H$}/{(c) $H\otimes H$}} {
    \begin{scope}[xshift=\xs cm]
      \draw (0,0.75) -- (2.6,0.75);
      \draw (0,0) -- (2.6,0);
      \node[anchor=east] at (0,0.75) {$B$};
      \node[anchor=east] at (0,0) {$A$};
      \ifx\haut\empty\else \node[porte] at (1.3,0.75) {\haut}; \fi
      \ifx\bas\empty\else \node[porte] at (1.3,0) {\bas}; \fi
      \node[font=\footnotesize] at (1.3,-0.85) {\tit};
    \end{scope}
  }
\end{tikzpicture}
\caption{Portes en parallèle sur deux qubits. Le qubit $A$, premier facteur du produit tensoriel,
est sur le fil du bas. (a) $X$ sur $A$ seulement. (b) $X$ sur $B$ seulement. (c) $H$ sur les deux.}
\label{fig:local}
\end{figure}

\subsubsection{Portes agissant sur un seul des deux qubits}

On note $I=\mathrm{Id}$ ($2\times2$). $X\otimes I$ agit sur $A$, $I\otimes X$ agit sur $B$ :
\begin{align*}
  I\otimes X&=\begin{pmatrix}1\cdot X & 0\cdot X\\ 0\cdot X & 1\cdot X\end{pmatrix}
    =\begin{pmatrix}0 & 1 & 0 & 0\\ 1 & 0 & 0 & 0\\ 0 & 0 & 0 & 1\\ 0 & 0 & 1 & 0\end{pmatrix},\\[1ex]
  X\otimes I&=\begin{pmatrix}0\cdot I & 1\cdot I\\ 1\cdot I & 0\cdot I\end{pmatrix}
    =\begin{pmatrix}0 & 0 & 1 & 0\\ 0 & 0 & 0 & 1\\ 1 & 0 & 0 & 0\\ 0 & 1 & 0 & 0\end{pmatrix}.
\end{align*}
Sur $\ket{00}=(1,0,0,0)^{\mathsf T}$ (première colonne de la matrice) :
$(I\otimes X)\ket{00}=(0,1,0,0)^{\mathsf T}=\ket{01}$ et
$(X\otimes I)\ket{00}=(0,0,1,0)^{\mathsf T}=\ket{10}$. Avec la règle :
$(I\otimes X)\ket{00}=I\ket{0}\otimes X\ket{0}=\ket{0}\otimes\ket{1}=\ket{01}$.

Pour la porte de Hadamard, avec $H=\frac{1}{\sqrt2}\begin{pmatrix}1 & 1\\ 1 & -1\end{pmatrix}$ :
\[
  H\otimes I=\frac{1}{\sqrt2}\begin{pmatrix}1\cdot I & 1\cdot I\\ 1\cdot I & -1\cdot I\end{pmatrix}
  =\frac{1}{\sqrt2}\begin{pmatrix}1 & 0 & 1 & 0\\ 0 & 1 & 0 & 1\\ 1 & 0 & -1 & 0\\ 0 & 1 & 0 & -1\end{pmatrix}.
\]
Ses colonnes donnent l'image des quatre états de base :
\begin{align*}
  (H\otimes I)\ket{00}&=\tfrac{1}{\sqrt2}(1,0,1,0)^{\mathsf T}=\ket{+}\otimes\ket{0},&
  (H\otimes I)\ket{01}&=\tfrac{1}{\sqrt2}(0,1,0,1)^{\mathsf T}=\ket{+}\otimes\ket{1},\\
  (H\otimes I)\ket{10}&=\tfrac{1}{\sqrt2}(1,0,-1,0)^{\mathsf T}=\ket{-}\otimes\ket{0},&
  (H\otimes I)\ket{11}&=\tfrac{1}{\sqrt2}(0,1,0,-1)^{\mathsf T}=\ket{-}\otimes\ket{1}.
\end{align*}

\subsubsection{Hadamard sur les deux qubits}

\[
  H\otimes H=\frac12\begin{pmatrix}1\cdot H' & 1\cdot H'\\ 1\cdot H' & -1\cdot H'\end{pmatrix},
  \quad H'=\begin{pmatrix}1 & 1\\ 1 & -1\end{pmatrix},
  \qquad
  H\otimes H=\frac12\begin{pmatrix}1 & 1 & 1 & 1\\ 1 & -1 & 1 & -1\\ 1 & 1 & -1 & -1\\ 1 & -1 & -1 & 1\end{pmatrix}.
\]
La première colonne donne $(H\otimes H)\ket{00}=\frac12(1,1,1,1)^{\mathsf T}=\ket{+}\otimes\ket{+}$
(cohérent avec $H\ket{0}\otimes H\ket{0}$) : les quatre résultats $00,01,10,11$ sont
équiprobables, chacun avec la probabilité $\frac14$. Sur $n$ qubits, $H^{\otimes n}\ket{0\ldots0}$
donne une superposition uniforme de $2^n$ états.

% ---------------------------------------------------------------------
\subsection{Porte CNOT}\label{sec:cnot}
% ---------------------------------------------------------------------

\begin{definition}{Porte CNOT (NOT contrôlé)}
La porte CNOT agit sur deux qubits : le qubit $B$ est le \emph{contrôle}, le qubit $A$ est la
\emph{cible}. La cible est inversée (porte $X$) si le contrôle vaut $1$ et reste inchangée
s'il vaut $0$. Avec $\ket{a,b}=\ket{a}_A\otimes\ket{b}_B$ :
\[
  \mathrm{CNOT}\,\ket{a,b}=\ket{a\oplus b,\,b},
  \qquad
  \mathrm{CNOT}=I\otimes\ket{0}\bra{0}+X\otimes\ket{1}\bra{1},
\]
où $\oplus$ est l'addition modulo $2$ ($0\oplus0=1\oplus1=0$ et $0\oplus1=1\oplus0=1$).
\end{definition}

Sur le circuit de la figure~\ref{fig:cnot}, le point noir posé sur le fil de $B$ (en haut) marque
le contrôle et le symbole $\oplus$ posé sur le fil de $A$ (en bas) marque la cible.

\begin{figure}[!ht]
\centering
\begin{tikzpicture}[>=Stealth, font=\small]
  \draw (0,1.2) -- (6,1.2);
  \draw (0,0) -- (6,0);
  \node[anchor=east] at (0,1.2) {$\ket{b}_B$};
  \node[anchor=east] at (0,0) {$\ket{a}_A$};
  \node[anchor=west] at (6,1.2) {$\ket{b}_B$};
  \node[anchor=west] at (6,0) {$\ket{a\oplus b}_A$};
  \draw[insarouge, thick] (3,1.2) -- (3,0);
  \pic at (3,1.2) {ctl};
  \pic at (3,0) {cible};
  \node[anchor=south, text=insagris] at (3,1.5) {qubit de contrôle};
  \node[anchor=north, text=insagris] at (3,-0.5) {qubit cible};
\end{tikzpicture}
\caption{Porte CNOT : le fil du haut ($B$) est le contrôle, le fil du bas ($A$) est la cible. Le
ket s'écrit $\ket{a,b}$ : le fil du haut donne le dernier chiffre.}
\label{fig:cnot}
\end{figure}

\subsubsection{Table de vérité}

\begin{center}
\renewcommand{\arraystretch}{1.3}
\begin{tabular}{@{}cccc@{}}
\toprule
Entrée $\ket{a,b}$ & Contrôle $b$ & Sortie $\ket{a\oplus b,b}$ & Effet sur la cible $A$ \\
\midrule
$\ket{00}$ & $0$ & $\ket{00}$ & aucun \\
$\ket{01}$ & $1$ & $\ket{11}$ & $0\to1$ \\
$\ket{10}$ & $0$ & $\ket{10}$ & aucun \\
$\ket{11}$ & $1$ & $\ket{01}$ & $1\to0$ \\
\bottomrule
\end{tabular}
\end{center}

\subsubsection{Matrice du CNOT}

La $j$-ième colonne d'une matrice $4\times4$ est l'image du $j$-ième vecteur de base
($\ket{00},\ket{01},\ket{10},\ket{11}$, dans cet ordre). La table donne :
\begin{align*}
  \mathrm{CNOT}\ket{00}&=\ket{00}=\begin{pmatrix}1\\0\\0\\0\end{pmatrix},&
  \mathrm{CNOT}\ket{01}&=\ket{11}=\begin{pmatrix}0\\0\\0\\1\end{pmatrix},\\[1.5ex]
  \mathrm{CNOT}\ket{10}&=\ket{10}=\begin{pmatrix}0\\0\\1\\0\end{pmatrix},&
  \mathrm{CNOT}\ket{11}&=\ket{01}=\begin{pmatrix}0\\1\\0\\0\end{pmatrix},
\end{align*}
d'où, en juxtaposant ces quatre colonnes,
\[
  \mathrm{CNOT}=\begin{pmatrix}
    1 & 0 & 0 & 0\\ 0 & 0 & 0 & 1\\ 0 & 0 & 1 & 0\\ 0 & 1 & 0 & 0
  \end{pmatrix}.
\]

On retrouve cette matrice avec la formule $I\otimes\ket{0}\bra{0}+X\otimes\ket{1}\bra{1}$ et
la règle $A\otimes B=(A_{ij}B)$ (section~\ref{sec:tenseur}), où
$\ket{0}\bra{0}=\begin{pmatrix}1&0\\0&0\end{pmatrix}$ et
$\ket{1}\bra{1}=\begin{pmatrix}0&0\\0&1\end{pmatrix}$ :
\[
  I\otimes\ket{0}\bra{0}=\begin{pmatrix}1&0&0&0\\0&0&0&0\\0&0&1&0\\0&0&0&0\end{pmatrix},
  \qquad
  X\otimes\ket{1}\bra{1}=\begin{pmatrix}0&0&0&0\\0&0&0&1\\0&0&0&0\\0&1&0&0\end{pmatrix},
\]
et la somme de ces deux matrices est la matrice du CNOT.

\subsubsection{Action sur un état quelconque}

Pour $\ket{\Psi}=c_{00}\ket{00}+c_{01}\ket{01}+c_{10}\ket{10}+c_{11}\ket{11}$ :
\[
  \mathrm{CNOT}\begin{pmatrix}c_{00}\\c_{01}\\c_{10}\\c_{11}\end{pmatrix}
  =\begin{pmatrix}
    1\cdot c_{00}+0\cdot c_{01}+0\cdot c_{10}+0\cdot c_{11}\\
    0\cdot c_{00}+0\cdot c_{01}+0\cdot c_{10}+1\cdot c_{11}\\
    0\cdot c_{00}+0\cdot c_{01}+1\cdot c_{10}+0\cdot c_{11}\\
    0\cdot c_{00}+1\cdot c_{01}+0\cdot c_{10}+0\cdot c_{11}
  \end{pmatrix}
  =\begin{pmatrix}c_{00}\\c_{11}\\c_{10}\\c_{01}\end{pmatrix}.
\]
Le CNOT échange les amplitudes de $\ket{01}$ et de $\ket{11}$, c'est-à-dire des deux états
dont le contrôle vaut $1$, et laisse les autres.

\paragraph{Exemple : contrôle en superposition.} Soit $\ket{0}_A\otimes\ket{+}_B
=\frac{1}{\sqrt2}\big(\ket{00}+\ket{01}\big)=\frac{1}{\sqrt2}(1,1,0,0)^{\mathsf T}$,
un état séparable. Alors
\[
  \mathrm{CNOT}\,\frac{1}{\sqrt2}\begin{pmatrix}1\\1\\0\\0\end{pmatrix}
  =\frac{1}{\sqrt2}\begin{pmatrix}1\\0\\0\\1\end{pmatrix}
  =\frac{1}{\sqrt2}\big(\ket{00}+\ket{11}\big).
\]
Le résultat est intriqué ($c_{00}c_{11}-c_{01}c_{10}=\frac12\neq0$ ; section~\ref{sec:separable}) :
un CNOT dont le contrôle est en superposition crée de l'intrication. C'est l'état de Bell
$\ket{\Phi^+}$ du chapitre~\ref{chap:bell}.

\paragraph{Réversibilité.} Deux CNOT successifs redonnent l'état initial :
$\ket{a,b}\to\ket{a\oplus b,b}\to\ket{a\oplus b\oplus b,b}=\ket{a,b}$, donc
$\mathrm{CNOT}^2=\mathrm{Id}$. La matrice est réelle et symétrique, donc
$\mathrm{CNOT}^\dagger\mathrm{CNOT}=\mathrm{CNOT}^2=\mathrm{Id}$ : elle est bien unitaire.

\begin{remarque}[CNOT et copie d'un bit]
Avec la cible dans l'état $\ket{0}$, $\mathrm{CNOT}\ket{0,b}=\ket{b,b}$ pour $b=0$ et $b=1$ : le CNOT
copie le bit classique $b$. En revanche $\mathrm{CNOT}\big(\ket{0}\otimes\ket{+}\big)$ est l'état
ci-dessus, et non $\ket{+}\otimes\ket{+}$ : un état quantique inconnu ne peut pas être copié
(théorème de non-clonage, section~\ref{sec:teleportation}).
\end{remarque}

\begin{remarque}[Contrôle en dernier ou en premier]
La matrice dépend de l'ordre d'écriture du ket. Notons $\mathrm{CNOT}_{21}$ le CNOT dont le
contrôle est le \emph{second} chiffre du ket : c'est la matrice ci-dessus, adoptée dans ce
document (fil du haut). Beaucoup d'ouvrages mettent le contrôle en \emph{premier} ; le CNOT
correspondant est
\[
  \mathrm{CNOT}_{12}=\ket{0}\bra{0}\otimes I+\ket{1}\bra{1}\otimes X
  =\begin{pmatrix}1&0&0&0\\0&1&0&0\\0&0&0&1\\0&0&1&0\end{pmatrix}.
\]
Le même circuit a donc deux matrices, selon que le contrôle est écrit en dernier ou en premier.
Les deux sont reliées par le SWAP (section~\ref{sec:swap}). La section~\ref{sec:controlees} sur les
portes contrôlées met le contrôle en premier, avec la matrice $\mathrm{diag}(I,U)$.
\end{remarque}

% ---------------------------------------------------------------------
\subsection{Porte SWAP}\label{sec:swap}
% ---------------------------------------------------------------------

\begin{definition}{Porte SWAP (échange)}
La porte SWAP échange les états des deux qubits :
\[
  \mathrm{SWAP}\big(\ket{\psi}\otimes\ket{\phi}\big)=\ket{\phi}\otimes\ket{\psi}.
\]
Sur les états de base, $\mathrm{SWAP}\ket{a,b}=\ket{b,a}$ ; on prolonge par linéarité.
\end{definition}

Son symbole est formé de deux croix reliées par un trait vertical
(figure~\ref{fig:swap}).

\begin{figure}[!ht]
\centering
\begin{tikzpicture}[>=Stealth, font=\small]
  \draw (0,1.2) -- (6,1.2);
  \draw (0,0) -- (6,0);
  \node[anchor=east] at (0,1.2) {$\ket{a}$};
  \node[anchor=east] at (0,0) {$\ket{b}$};
  \node[anchor=west] at (6,1.2) {$\ket{b}$};
  \node[anchor=west] at (6,0) {$\ket{a}$};
  \draw[insarouge, thick] (3,1.2) -- (3,0);
  \pic at (3,1.2) {croix};
  \pic at (3,0) {croix};
\end{tikzpicture}
\caption{Porte SWAP : les deux qubits échangent leurs états.}
\label{fig:swap}
\end{figure}

\subsubsection{Table et matrice}

\begin{center}
\renewcommand{\arraystretch}{1.3}
\begin{tabular}{@{}cc@{}}
\toprule
Entrée & Sortie \\
\midrule
$\ket{00}$ & $\ket{00}$ \\
$\ket{01}$ & $\ket{10}$ \\
$\ket{10}$ & $\ket{01}$ \\
$\ket{11}$ & $\ket{11}$ \\
\bottomrule
\end{tabular}
\end{center}

Les images des quatre vecteurs de base sont les colonnes de la matrice :
\[
  \mathrm{SWAP}\ket{00}=\begin{pmatrix}1\\0\\0\\0\end{pmatrix},\quad
  \mathrm{SWAP}\ket{01}=\begin{pmatrix}0\\0\\1\\0\end{pmatrix},\quad
  \mathrm{SWAP}\ket{10}=\begin{pmatrix}0\\1\\0\\0\end{pmatrix},\quad
  \mathrm{SWAP}\ket{11}=\begin{pmatrix}0\\0\\0\\1\end{pmatrix},
\]
\[
  \mathrm{SWAP}=\begin{pmatrix}
    1 & 0 & 0 & 0\\ 0 & 0 & 1 & 0\\ 0 & 1 & 0 & 0\\ 0 & 0 & 0 & 1
  \end{pmatrix}.
\]
Appliquée à $\ket{01}=(0,1,0,0)^{\mathsf T}$, cette matrice sélectionne sa deuxième colonne :
\[
  \begin{pmatrix}
    1 & 0 & 0 & 0\\ 0 & 0 & 1 & 0\\ 0 & 1 & 0 & 0\\ 0 & 0 & 0 & 1
  \end{pmatrix}
  \begin{pmatrix}0\\1\\0\\0\end{pmatrix}
  =\begin{pmatrix}0\\0\\1\\0\end{pmatrix}=\ket{10}.
\]
Sur un état quelconque, $\mathrm{SWAP}\,(c_{00},c_{01},c_{10},c_{11})^{\mathsf T}
=(c_{00},c_{10},c_{01},c_{11})^{\mathsf T}$ : les amplitudes de $\ket{01}$ et $\ket{10}$
sont échangées.

\subsubsection{Accord avec la définition}

Pour deux états $\ket{\psi}=(a_0,a_1)^{\mathsf T}$ et $\ket{\phi}=(b_0,b_1)^{\mathsf T}$ :
\[
  \mathrm{SWAP}\,(\ket{\psi}\otimes\ket{\phi})
  =\mathrm{SWAP}\begin{pmatrix}a_0b_0\\a_0b_1\\a_1b_0\\a_1b_1\end{pmatrix}
  =\begin{pmatrix}a_0b_0\\a_1b_0\\a_0b_1\\a_1b_1\end{pmatrix}
  =\begin{pmatrix}b_0\\b_1\end{pmatrix}\otimes\begin{pmatrix}a_0\\a_1\end{pmatrix}
  =\ket{\phi}\otimes\ket{\psi},
\]
ce qui est bien la définition. La porte s'écrit aussi avec des projecteurs :
\[
  \mathrm{SWAP}=\sum_{a,b\in\{0,1\}}\ket{a}\bra{b}\otimes\ket{b}\bra{a}.
\]
En effet, pour $\ket{c}\otimes\ket{d}$ on trouve
$\sum_{a,b}\ket{a}\braket{b}{c}\otimes\ket{b}\braket{a}{d}$, et seul le terme $b=c$, $a=d$
est non nul : le résultat est $\ket{d}\otimes\ket{c}$.

\paragraph{Exemples.}
\begin{itemize}[nosep]
  \item Deux états différents : $\mathrm{SWAP}\big(\ket{0}\otimes\ket{+}\big)
        =\mathrm{SWAP}\,\frac{1}{\sqrt2}(1,1,0,0)^{\mathsf T}=\frac{1}{\sqrt2}(1,0,1,0)^{\mathsf T}
        =\ket{+}\otimes\ket{0}$.
  \item L'état $\ket{\psi}=\frac{1}{\sqrt2}\ket{00}+\frac{1}{\sqrt2}\ket{11}$ ne change pas :
        $\mathrm{SWAP}\ket{\psi}=\frac{1}{\sqrt2}\ket{00}+\frac{1}{\sqrt2}\ket{11}=\ket{\psi}$,
        car $\ket{00}$ et $\ket{11}$ sont invariants. De même
        $\frac{1}{\sqrt2}(\ket{00}-\ket{11})$ et $\frac{1}{\sqrt2}(\ket{01}+\ket{10})$ sont invariants,
        alors que $\mathrm{SWAP}\,\frac{1}{\sqrt2}(\ket{01}-\ket{10})=\frac{1}{\sqrt2}(\ket{10}-\ket{01})$
        est l'opposé de l'état de départ. Ce sont les quatre états de Bell (chapitre~\ref{chap:bell}).
\end{itemize}

\paragraph{Propriétés.} Échanger deux fois redonne l'état de départ :
$\mathrm{SWAP}^2=\mathrm{Id}$ ; la matrice est réelle et symétrique, donc
$\mathrm{SWAP}^\dagger\mathrm{SWAP}=\mathrm{SWAP}^2=\mathrm{Id}$ : elle est unitaire.

\subsubsection{Lien avec le CNOT}

\begin{resultat}{SWAP échange les rôles de contrôle et de cible}
\[
  \mathrm{CNOT}_{21}=\mathrm{SWAP}\;\mathrm{CNOT}_{12}\;\mathrm{SWAP}.
\]
\end{resultat}

\emph{Sur les états de base.}
$\mathrm{SWAP}\,\mathrm{CNOT}_{12}\,\mathrm{SWAP}\ket{a,b}
=\mathrm{SWAP}\,\mathrm{CNOT}_{12}\ket{b,a}
=\mathrm{SWAP}\ket{b,\,b\oplus a}=\ket{b\oplus a,\,b}=\mathrm{CNOT}_{21}\ket{a,b}$.

\emph{Par les matrices.} Multiplier à gauche par $\mathrm{SWAP}$ échange les lignes $2$ et
$3$, multiplier à droite l'échange les colonnes $2$ et $3$ :
\[
  \mathrm{SWAP}\,\mathrm{CNOT}_{12}
  =\begin{pmatrix}1&0&0&0\\0&0&0&1\\0&1&0&0\\0&0&1&0\end{pmatrix},
  \qquad
  \mathrm{SWAP}\,\mathrm{CNOT}_{12}\,\mathrm{SWAP}
  =\begin{pmatrix}1&0&0&0\\0&0&0&1\\0&0&1&0\\0&1&0&0\end{pmatrix}=\mathrm{CNOT}_{21}.
\]

\begin{remarque}[SWAP avec trois CNOT]
$\mathrm{SWAP}=\mathrm{CNOT}_{12}\,\mathrm{CNOT}_{21}\,\mathrm{CNOT}_{12}$. En effet, sur
$\ket{a,b}$ (on applique les portes de droite à gauche) :
\[
  \ket{a,b}\ \xrightarrow{\ \mathrm{CNOT}_{12}\ }\ \ket{a,\,a\oplus b}
  \ \xrightarrow{\ \mathrm{CNOT}_{21}\ }\ \ket{a\oplus(a\oplus b),\,a\oplus b}=\ket{b,\,a\oplus b}
  \ \xrightarrow{\ \mathrm{CNOT}_{12}\ }\ \ket{b,\,b\oplus(a\oplus b)}=\ket{b,a}.
\]
Les deux membres coïncident sur les états de base, donc sur tous les états par linéarité.
\end{remarque}

% ---------------------------------------------------------------------
\subsection{Portes contrôlées : contrôlé-\texorpdfstring{$U$}{U}, Fredkin et Toffoli}\label{sec:controlees}
% ---------------------------------------------------------------------

Le CNOT est le cas le plus simple d'une idée générale : appliquer une opération à un ou
plusieurs qubits \emph{cibles} seulement si un ou plusieurs qubits de \emph{contrôle} valent $1$.
Dans cette section, le contrôle est le \emph{premier} chiffre
du ket (remarque de la section~\ref{sec:cnot}). Les qubits sont notés $x$, $y$, $z$ (le qubit $x$ est le contrôle) et
$\ket{x,y,z}=\ket{x}\otimes\ket{y}\otimes\ket{z}$ ; d'après la convention de la section~\ref{sec:circuits}, le
contrôle est dessiné en bas (figure~\ref{fig:controlees}). Pour éviter la confusion avec les
matrices de Pauli, on écrit ici $\sigma_x=X$ pour la porte NOT.

\begin{figure}[!ht]
\centering
\begin{adjustbox}{max width=\linewidth}
\begin{tikzpicture}[>=Stealth, font=\small]
  % (a) contrôlé-U : contrôle x en bas, cible y en haut
  \draw (0,1.2) -- (3.4,1.2);
  \draw (0,0) -- (3.4,0);
  \node[anchor=east] at (0,1.2) {$y$};
  \node[anchor=east] at (0,0) {$x$};
  \draw[insarouge, thick] (1.7,0) -- (1.7,0.85);
  \pic at (1.7,0) {ctl};
  \node[porte] at (1.7,1.2) {$U$};
  \node[anchor=north, font=\footnotesize] at (1.7,-0.5) {(a) contrôlé-$U$};
  % (b) contrôlé-SWAP (Fredkin) : x contrôle, y et z échangés
  \draw (5,2.4) -- (8.4,2.4);
  \draw (5,1.2) -- (8.4,1.2);
  \draw (5,0) -- (8.4,0);
  \node[anchor=east] at (5,2.4) {$z$};
  \node[anchor=east] at (5,1.2) {$y$};
  \node[anchor=east] at (5,0) {$x$};
  \draw[insarouge, thick] (6.7,0) -- (6.7,2.4);
  \pic at (6.7,0) {ctl};
  \pic at (6.7,1.2) {croix};
  \pic at (6.7,2.4) {croix};
  \node[anchor=north, font=\footnotesize] at (6.7,-0.5) {(b) Fredkin (C-SWAP)};
  % (c) Toffoli (CCNOT) : x et y contrôlent, z cible
  \draw (10,2.4) -- (13.4,2.4);
  \draw (10,1.2) -- (13.4,1.2);
  \draw (10,0) -- (13.4,0);
  \node[anchor=east] at (10,2.4) {$z$};
  \node[anchor=east] at (10,1.2) {$y$};
  \node[anchor=east] at (10,0) {$x$};
  \draw[insarouge, thick] (11.7,0) -- (11.7,2.4);
  \pic at (11.7,0) {ctl};
  \pic at (11.7,1.2) {ctl};
  \pic at (11.7,2.4) {cible};
  \node[anchor=north, font=\footnotesize] at (11.7,-0.5) {(c) Toffoli (CCNOT)};
\end{tikzpicture}
\end{adjustbox}
\caption{Portes contrôlées. Le contrôle $x$ est le premier chiffre du ket, donc le fil du bas.
(a) Le qubit $y$ subit $U$ si $x=1$ ; pour $U=X$ on retrouve le CNOT. (b) Si $x=1$, les qubits
$y$ et $z$ sont échangés. (c) Si $x=y=1$, le qubit $z$ est inversé.}
\label{fig:controlees}
\end{figure}

\subsubsection{Contrôlé-\texorpdfstring{$U$}{U}}

\begin{definition}{Porte contrôlée-$U$}
Soit $U$ une matrice unitaire $2\times2$. Sur le système $(x,y)$, de contrôle $x$ et de cible $y$,
\[
  \mathrm{C}U=\ket{0}\bra{0}_x\otimes I_y+\ket{1}\bra{1}_x\otimes U_y .
\]
Le premier terme agit quand le contrôle est $\ket{0}$ : la cible ne change pas ($I$). Le second
agit quand le contrôle est $\ket{1}$ : la cible subit $U$.
\end{definition}

Comme $\ket{0}\bra{0}\ket{x}=\delta_{x0}\ket{0}$ et $\ket{1}\bra{1}\ket{x}=\delta_{x1}\ket{1}$,
\[
  \mathrm{C}U\big(\ket{x}\otimes\ket{y}\big)=\ket{x}\otimes U^x\ket{y},\qquad U^0=I,\ U^1=U .
\]
Par la règle $A\otimes B=(A_{ij}B)$, $\ket{0}\bra{0}\otimes I=\begin{pmatrix}I&0\\0&0\end{pmatrix}$
et $\ket{1}\bra{1}\otimes U=\begin{pmatrix}0&0\\0&U\end{pmatrix}$ (blocs $2\times2$), donc
\[
  \mathrm{C}U=\begin{pmatrix}I&0\\0&U\end{pmatrix}
  =\begin{pmatrix}1&0&0&0\\0&1&0&0\\0&0&u_{00}&u_{01}\\0&0&u_{10}&u_{11}\end{pmatrix}
  \qquad\text{pour }U=\begin{pmatrix}u_{00}&u_{01}\\u_{10}&u_{11}\end{pmatrix}.
\]
Cette matrice est unitaire car $I^\dagger I=I$ et $U^\dagger U=I$.

\paragraph{Le CNOT est un cas particulier.} Pour $U=\sigma_x$,
$\mathrm{C}\sigma_x=\ket{0}\bra{0}_x\otimes I_y+\ket{1}\bra{1}_x\otimes\sigma_x$ est le CNOT de
contrôle $x$ (premier chiffre), c'est-à-dire $\mathrm{CNOT}_{12}$ de la section~\ref{sec:cnot}.
On calcule l'action sur les quatre états de base, avec $\braket{0}{0}=\braket{1}{1}=1$,
$\braket{0}{1}=\braket{1}{0}=0$, $\sigma_x\ket{0}=\ket{1}$ et $\sigma_x\ket{1}=\ket{0}$ :
\begin{align*}
  \mathrm{C}\sigma_x\ket{00}&=\big[\ket{0}\bra{0}_x\otimes I_y\big]\ket{0}_x\ket{0}_y
      +\big[\ket{1}\bra{1}_x\otimes\sigma_x\big]\ket{0}_x\ket{0}_y\\
    &=\ket{0}\braket{0}{0}\otimes I\ket{0}+\ket{1}\braket{1}{0}\otimes\sigma_x\ket{0}
     =\ket{0}\otimes\ket{0}+0=\ket{00},\\[1.2ex]
  \mathrm{C}\sigma_x\ket{01}&=\ket{0}\braket{0}{0}\otimes I\ket{1}+\ket{1}\braket{1}{0}\otimes\sigma_x\ket{1}
     =\ket{0}\otimes\ket{1}+0=\ket{01},\\[1.2ex]
  \mathrm{C}\sigma_x\ket{10}&=\ket{0}\braket{0}{1}\otimes I\ket{0}+\ket{1}\braket{1}{1}\otimes\sigma_x\ket{0}
     =0+\ket{1}\otimes\ket{1}=\ket{11},\\[1.2ex]
  \mathrm{C}\sigma_x\ket{11}&=\ket{0}\braket{0}{1}\otimes I\ket{1}+\ket{1}\braket{1}{1}\otimes\sigma_x\ket{1}
     =0+\ket{1}\otimes\ket{0}=\ket{10}.
\end{align*}
On retrouve la table du CNOT : la cible n'est inversée que si le contrôle vaut $1$.

\paragraph{Un autre exemple : $U=Z$.} Ici
\[
  \mathrm{C}Z=\begin{pmatrix}I&0\\0&Z\end{pmatrix}=\mathrm{diag}(1,1,1,-1) :
\]
la porte multiplie $\ket{11}$ par $-1$ et ne touche pas les autres états de base. Elle est
symétrique entre les deux qubits. Elle se déduit du CNOT par
$\mathrm{C}Z=(I\otimes H)\,\mathrm{C}X\,(I\otimes H)$, où $\mathrm{C}X$ est le contrôlé-$X$.
En effet, $I\otimes H=\mathrm{diag}(H,H)$ (deux blocs $2\times2$ égaux à $H$), et le produit de
matrices diagonales par blocs se fait bloc par bloc :
\[
  (I\otimes H)\begin{pmatrix}I&0\\0&X\end{pmatrix}(I\otimes H)
  =\begin{pmatrix}H\,I\,H&0\\0&H\,X\,H\end{pmatrix}
  =\begin{pmatrix}I&0\\0&HXH\end{pmatrix},
\]
car $H^2=I$, et
\[
  HXH=\frac12\begin{pmatrix}1&1\\1&-1\end{pmatrix}\begin{pmatrix}0&1\\1&0\end{pmatrix}\begin{pmatrix}1&1\\1&-1\end{pmatrix}
  =\frac12\begin{pmatrix}1&1\\1&-1\end{pmatrix}\begin{pmatrix}1&-1\\1&1\end{pmatrix}
  =\frac12\begin{pmatrix}2&0\\0&-2\end{pmatrix}=Z
\]
(section~\ref{sec:hadamard} : $H$ échange les axes $x$ et $z$).

\subsubsection{Contrôlé-SWAP : la porte de Fredkin}

\begin{definition}{Porte de Fredkin (C-SWAP)}
Sur trois qubits $(x,y,z)$ (espace de dimension $2^3=8$), avec $x$ pour contrôle,
\[
  \mathrm{C\text{-}SWAP}=\ket{0}\bra{0}_x\otimes I_y\otimes I_z+\ket{1}\bra{1}_x\otimes\mathrm{SWAP}_{yz}.
\]
Si $x=0$, rien ne change ; si $x=1$, les qubits $y$ et $z$ sont échangés.
\end{definition}

Calcul sur quatre états de base (avec $I\ket{ab}=\ket{ab}$ et $\mathrm{SWAP}\ket{ab}=\ket{ba}$) :
\begin{align*}
  \mathrm{C\text{-}SWAP}\ket{000}&=\ket{0}\braket{0}{0}\otimes I\ket{00}+\ket{1}\braket{1}{0}\otimes\mathrm{SWAP}\ket{00}
     =\ket{0}\otimes\ket{00}+0=\ket{000},\\[1.2ex]
  \mathrm{C\text{-}SWAP}\ket{101}&=\ket{0}\braket{0}{1}\otimes I\ket{01}+\ket{1}\braket{1}{1}\otimes\mathrm{SWAP}\ket{01}
     =0+\ket{1}\otimes\ket{10}=\ket{110},\\[1.2ex]
  \mathrm{C\text{-}SWAP}\ket{110}&=\ket{0}\braket{0}{1}\otimes I\ket{10}+\ket{1}\braket{1}{1}\otimes\mathrm{SWAP}\ket{10}
     =0+\ket{1}\otimes\ket{01}=\ket{101},\\[1.2ex]
  \mathrm{C\text{-}SWAP}\ket{010}&=\ket{0}\braket{0}{0}\otimes I\ket{10}+\ket{1}\braket{1}{0}\otimes\mathrm{SWAP}\ket{10}
     =\ket{0}\otimes\ket{10}+0=\ket{010}.
\end{align*}
Les huit états de base donnent :
\begin{center}
\renewcommand{\arraystretch}{1.3}
\begin{tabular}{@{}ccccccccc@{}}
\toprule
Entrée & $\ket{000}$ & $\ket{001}$ & $\ket{010}$ & $\ket{011}$ & $\ket{100}$ & $\ket{101}$ & $\ket{110}$ & $\ket{111}$ \\
\midrule
Sortie & $\ket{000}$ & $\ket{001}$ & $\ket{010}$ & $\ket{011}$ & $\ket{100}$ & $\ket{110}$ & $\ket{101}$ & $\ket{111}$ \\
\bottomrule
\end{tabular}
\end{center}
Seuls $\ket{101}$ et $\ket{110}$ changent (contrôle à $1$ et $y\neq z$) ; pour $\ket{100}$ et
$\ket{111}$ l'échange de deux qubits égaux ne change rien. Dans la base
$\ket{000},\ket{001},\dots,\ket{111}$ (ordre binaire), la matrice est
$\mathrm{diag}(I_4,\mathrm{SWAP})$ (deux blocs $4\times4$) :
\[
  \mathrm{C\text{-}SWAP}=\begingroup\setlength{\arraycolsep}{3pt}\begin{pmatrix}
    1&0&0&0&0&0&0&0\\
    0&1&0&0&0&0&0&0\\
    0&0&1&0&0&0&0&0\\
    0&0&0&1&0&0&0&0\\
    0&0&0&0&1&0&0&0\\
    0&0&0&0&0&0&1&0\\
    0&0&0&0&0&1&0&0\\
    0&0&0&0&0&0&0&1
  \end{pmatrix}\endgroup .
\]
C'est l'identité dont les lignes $6$ et $7$ sont échangées.

\paragraph{Contrôle en superposition.} Avec $\ket{+}_x\otimes\ket{0}_y\otimes\ket{1}_z
=\frac{1}{\sqrt2}\big(\ket{001}+\ket{101}\big)$ :
\[
  \mathrm{C\text{-}SWAP}\,\frac{1}{\sqrt2}\big(\ket{001}+\ket{101}\big)
  =\frac{1}{\sqrt2}\big(\ket{001}+\ket{110}\big).
\]
Le résultat n'est pas un produit : les coefficients de $\ket{x}\otimes\ket{yz}$ forment la
matrice (lignes $x=0,1$, colonnes $yz=00,01,10,11$)
\[
  \frac{1}{\sqrt2}\begin{pmatrix}0&1&0&0\\0&0&1&0\end{pmatrix},
\]
dont les deux lignes ne sont pas proportionnelles. Le qubit $x$ est intriqué avec la paire
$(y,z)$.

\subsubsection{Contrôlé-contrôlé-NOT : la porte de Toffoli}

\begin{definition}{Porte de Toffoli (CCNOT)}
Sur trois qubits $(x,y,z)$, avec $x$ et $y$ pour contrôles et $z$ pour cible,
\[
  \mathrm{CCNOT}=\ket{0}\bra{0}_x\otimes I_y\otimes I_z
   +\ket{1}\bra{1}_x\otimes\Big[\ket{0}\bra{0}_y\otimes I_z+\ket{1}\bra{1}_y\otimes\sigma_x\Big].
\]
Le crochet est un CNOT (contrôle $y$, cible $z$), appliqué seulement si $x=1$ : si $x=0$, rien ;
si $x=1$ et $y=0$, rien ; si $x=1$ et $y=1$, le qubit $z$ est inversé.
\end{definition}

Calcul de $\mathrm{CCNOT}\ket{110}$ :
\begin{align*}
  \mathrm{CCNOT}\ket{110}&=\ket{0}\braket{0}{1}\otimes I\ket{1}\otimes I\ket{0}
     +\ket{1}\braket{1}{1}\otimes\Big[\ket{0}\braket{0}{1}\otimes I\ket{0}
       +\ket{1}\braket{1}{1}\otimes\sigma_x\ket{0}\Big]\\
  &=0+\ket{1}\otimes\Big[0+\ket{1}\otimes\ket{1}\Big]=\ket{111}.
\end{align*}
Pour les huit états de base, $\mathrm{CCNOT}\ket{x,y,z}=\ket{x,\,y,\,z\oplus xy}$ :
\begin{center}
\renewcommand{\arraystretch}{1.3}
\begin{tabular}{@{}ccccccccc@{}}
\toprule
Entrée & $\ket{000}$ & $\ket{001}$ & $\ket{010}$ & $\ket{011}$ & $\ket{100}$ & $\ket{101}$ & $\ket{110}$ & $\ket{111}$ \\
\midrule
Sortie & $\ket{000}$ & $\ket{001}$ & $\ket{010}$ & $\ket{011}$ & $\ket{100}$ & $\ket{101}$ & $\ket{111}$ & $\ket{110}$ \\
\bottomrule
\end{tabular}
\end{center}
Seuls $\ket{110}$ et $\ket{111}$ sont échangés : la matrice est $\mathrm{diag}(I_6,X)$
(un bloc $6\times6$ égal à l'identité, puis $X$),
\[
  \mathrm{CCNOT}=\begingroup\setlength{\arraycolsep}{3pt}\begin{pmatrix}
    1&0&0&0&0&0&0&0\\
    0&1&0&0&0&0&0&0\\
    0&0&1&0&0&0&0&0\\
    0&0&0&1&0&0&0&0\\
    0&0&0&0&1&0&0&0\\
    0&0&0&0&0&1&0&0\\
    0&0&0&0&0&0&0&1\\
    0&0&0&0&0&0&1&0
  \end{pmatrix}\endgroup .
\]

\begin{remarque}[Portes réversibles]
Les portes de Fredkin et de Toffoli sont des permutations des états de base qui échangent
seulement deux d'entre eux : elles sont leurs propres inverses ($\mathrm{C\text{-}SWAP}^2=\mathrm{CCNOT}^2=\mathrm{Id}$)
et unitaires. Avec $z=0$, la Toffoli calcule le ET logique $x\wedge y$ sans rien effacer :
$\ket{x,y,0}\to\ket{x,y,xy}$ ; avec $z=1$, elle calcule le NON-ET. Elle sert à écrire tout calcul
classique de façon réversible.
\end{remarque}

\clearpage
% ---------------------------------------------------------------------
\subsection{Récapitulatif du chapitre}\label{sec:recap4}
% ---------------------------------------------------------------------

\begin{center}
\adjustbox{max width=\linewidth}{%
\begin{tikzpicture}[>=Stealth, font=\footnotesize,
    boite/.style={draw=insarouge, fill=insarouge!6, rounded corners=2pt, thick,
                  minimum width=4.4cm, minimum height=1.5cm, align=center, inner sep=3pt},
    lien/.style={->, very thick, insagris},
    etiq/.style={font=\scriptsize\itshape, text=insagris, inner sep=2pt}]
  \node[boite] (uni) at (0,0)
    {\textbf{Porte unitaire}\\ $U^\dagger U=\mathrm{Id}$, \ $\ket{\Psi}\to U\ket{\Psi}$\\ {\scriptsize\S\,\ref{sec:portes}}};
  \node[boite] (un) at (5.9,0)
    {\textbf{Portes à un qubit}\\ $X$, $Y$, $Z$, $H$, $S$, $T$\\ {\scriptsize\S\,\ref{sec:pauli}--\ref{sec:phase}}};
  \node[boite] (blo) at (11.8,0)
    {\textbf{Rotations de Bloch}\\ $X,Y,Z,H$ : angle $\pi$ ; $P_\theta$ autour de $z$\\
     {\scriptsize\S\,\ref{sec:portes-bloch}}};
  \node[boite] (cir) at (11.8,-2.5)
    {\textbf{Circuits}\\ fils, boîtes, $U_2U_1$\\ {\scriptsize\S\,\ref{sec:circuits}, \ref{sec:op2}}};
  \node[boite] (cn) at (5.9,-2.5)
    {\textbf{CNOT et SWAP}\\ $\ket{a,b}\to\ket{a\oplus b,b}$\\ {\scriptsize\S\,\ref{sec:cnot}, \ref{sec:swap}}};
  \node[boite] (cu) at (0,-2.5)
    {\textbf{Portes contrôlées}\\ $\mathrm{C}U=\mathrm{diag}(I,U)$, Fredkin, Toffoli\\
     {\scriptsize\S\,\ref{sec:controlees}}};
  \draw[lien] (uni) -- (un);
  \draw[lien] (un) -- (blo);
  \draw[lien] (blo) -- (cir);
  \draw[lien] (cir) -- (cn) node[etiq, midway, above=1pt] {2 qubits};
  \draw[lien] (cn) -- (cu) node[etiq, midway, above=1pt] {contrôle};
\end{tikzpicture}}
\end{center}

\begin{aretenir}
\begin{itemize}
  \item Une porte est un opérateur unitaire : elle conserve la norme, elle est réversible
        ($U^{-1}=U^\dagger$) et agit sur la sphère de Bloch comme une rotation. C'est l'évolution de
        Schrödinger, $U=\ee^{-\ii\Hop t/\hbar}$.
  \item $X$ échange $\ket{0}$ et $\ket{1}$, $Z$ change le signe de $\ket{1}$, $Y=\ii XZ$ fait les deux ;
        $H$ échange les bases $Z$ et $X$ et crée des superpositions ; $P_\theta$, $S=P_{\pi/2}$ et
        $T=P_{\pi/4}$ font tourner autour de $z$ ($T^2=S$, $S^2=Z$).
  \item Dans un circuit, le temps va de gauche à droite et la matrice est le produit dans l'ordre
        inverse, $U_2U_1$. Le fil du haut porte le dernier chiffre du ket. Des portes en parallèle
        forment un produit tensoriel $U_A\otimes U_B$.
  \item Le CNOT inverse la cible si le contrôle vaut $1$ : $\ket{a,b}\to\ket{a\oplus b,b}$. Le SWAP
        échange deux qubits ; $\mathrm{CNOT}_{21}=\mathrm{SWAP}\,\mathrm{CNOT}_{12}\,\mathrm{SWAP}$.
  \item Une porte contrôlée a pour matrice $\mathrm{diag}(I,U)$. Fredkin (contrôlé-SWAP) et
        Toffoli (CCNOT) sont des permutations des états de base : réversibles et universelles
        pour le calcul classique.
\end{itemize}
\end{aretenir}

\begin{center}
\renewcommand{\arraystretch}{1.45}
\begin{tabular}{@{}l l r@{}}
\toprule
\textbf{Notion} & \textbf{Formule} & \textbf{§} \\
\midrule
Porte unitaire & $U^\dagger U=\mathrm{Id}$, \quad $\bra{\Psi}U^\dagger U\ket{\Psi}=1$ & \ref{sec:portes} \\
Pauli & $X\ket{0}=\ket{1}$, \quad $Z\ket{1}=-\ket{1}$, \quad $X^2=Y^2=Z^2=\mathrm{Id}$ & \ref{sec:pauli} \\
Hadamard & $H\ket{0}=\ket{+}$, \quad $H\ket{1}=\ket{-}$, \quad $H^2=\mathrm{Id}$ & \ref{sec:hadamard} \\
Phase & $P_\theta=\mathrm{diag}(1,\ee^{\ii\theta})$, \quad $S=P_{\pi/2}$, \quad $T=P_{\pi/4}$ & \ref{sec:phase} \\
Circuit & $U_1$ puis $U_2$ : \quad $U_2U_1$ & \ref{sec:circuits} \\
Portes locales & $(U\otimes I)(\ket{u}\otimes\ket{v})=U\ket{u}\otimes\ket{v}$ & \ref{sec:op2} \\
CNOT & $\ket{a,b}\to\ket{a\oplus b,b}$, \quad $\mathrm{CNOT}=I\otimes\ket{0}\bra{0}+X\otimes\ket{1}\bra{1}$
  & \ref{sec:cnot} \\
SWAP & $\mathrm{SWAP}\ket{a,b}=\ket{b,a}$, \quad $\mathrm{SWAP}=\mathrm{CNOT}_{12}\mathrm{CNOT}_{21}\mathrm{CNOT}_{12}$
  & \ref{sec:swap} \\
Contrôlé-$U$ & $\mathrm{C}U=\ket{0}\bra{0}\otimes I+\ket{1}\bra{1}\otimes U=\mathrm{diag}(I,U)$ & \ref{sec:controlees} \\
Toffoli & $\mathrm{CCNOT}\ket{x,y,z}=\ket{x,y,z\oplus xy}$ & \ref{sec:controlees} \\
\bottomrule
\end{tabular}
\end{center}

\begin{savoirfaire}
\begin{itemize}
  \item Appliquer une porte à un état (par la matrice ou par les ket-bra) et identifier sa rotation
        de Bloch (exercices~\ref{ex:conj} et~\ref{ex:phases}).
  \item Lire un circuit, écrire sa matrice dans le bon ordre et suivre l'état à chaque étape
        (exercice~\ref{ex:kickback}).
  \item Écrire la matrice de $U\otimes I$, du CNOT, du SWAP ou d'une porte contrôlée, et calculer
        leur action sur les états de base.
  \item Reconnaître une porte réversible qui calcule une fonction booléenne
        (exercice~\ref{ex:toffoli}).
\end{itemize}
\end{savoirfaire}
